%&program=xetex
%&encoding=UTF-8 Unicode

\TeXXeTstate=1
\frenchspacing
%\font\title="Nafees Web Naskh pakistani BTK:script=arab" at 72pt
%\font\heading="Nafees Web Naskh pakistani BTK:script=arab" at 18pt
%\font\body="Nafees Web Naskh pakistani BTK:script=arab" at 15pt
%\font\bodybold="Nafees Web Naskh pakistani BTK:script=arab" at 15pt

\font\title="Nafees Tahreer Naskh v1.0:script=arab" at 72pt
\font\heading="Nafees Tahreer Naskh v1.0:script=arab" at 18pt
\font\body="Nafees Tahreer Naskh v1.0:script=arab" at 15pt
\font\bodybold="Nafees Tahreer Naskh v1.0:script=arab" at 15pt
\parindent=0.5in \baselineskip=22pt \lineskiplimit=-1000pt

\def\ttl#1{\beginR\title #1\endR}
\def\hdr#1{\beginR\heading #1\endR}
\def\pb#1{\beginR\body #1\endR}


\vfil
\centerline{\ttl{جڑوں کی تلاش}}
\bigskip\bigskip\bigskip

\centerline{\hdr{ابن صفی}}
\break

\everypar={\setbox0=\lastbox \beginR \box0 }
\centerline{\hdr{(۱)}}
\bigskip\bigskip\bigskip
\beginR
\body
ڈھمپ اینڈ کو کا دفتر بڑے مزے میں چل رہا تھا مگر اس کی منیجری کم از کم خاور کے بس کا روگ نہیں تھی کیونکہ بزنس کے چکروں کے لئے اس کا ذہن موزوں نہیں تھا۔ ذہن موزوں رہا ہو یا نہ رہا ہو لیکن صورت تو ضرور ایسی تھی کہ وہ کسی فرم کا منیجر معلوم ہوسکتا تھا! بھاری بھرکم بارعب چہرے والا۔۔!

چونکہ وہ بزنس کے معاملہ میں اناڑی تھا اس لئے اس کے کمرے میں لکڑی کی ایک دیوار سے پارٹیشنز کردیئے گئے تھے ایک طرف جولیانا بیٹھی ٹائپ رائٹر کھٹکا یا کرتی تھی اور دوسری طرف خاور اپنی مینجری سمیت براجمان رہا کرتا تھا!

اگر کبھی کوئی نیا گاہک آجاتا اور خاور کو اسے ڈیل کرنے میں کچھ دشواری محسوس ہوتی تو جولیا کاغذات کا پلندہ دبائے دستخط کرانے کے بہانے اس کی میز پر آجاتی اور دوران گفتگو میں دخل اندازی کرکے خاور کو سہارا دیئے رہتی۔۔!

آج بھی کوئی بڑا گاہک خاور کی میز پر موجود تھا اور اپنے کام کے سلسلے میں بعض امور کی وضاحت چاہتا تھا! جولیا نے محسوس کیا کہ خاور رک رک کر گفتگو کر رہا ہے اور گاہک کے ٹوکنے پر بعض اوقات گڑبڑا بھی جاتا ہے۔۔!

وہ کچھ کاغذات سنبھالے ہوئے خاور کی میز پر جا پہنچی!

"اوہو۔ اچھا ہوا تم آگئیں۔۔" خاور نے کہا اور پھر گاہک سے بولا۔ " یہ میری اسسٹنٹ ہیں سرسوکھے! میرا داہنا ہاتھ۔۔ اب دیکھیئے آپ جو کچھ چاہتے ہیں اس کام کا تعلق زیادہ تر انہیں کی ذات سے ہوگا!۔۔ حسابات وغیرہ کی پڑتال یہی کرتی ہیں"۔

جولیا نے اس گول مٹول آدمی پر ایک اچٹتی سی نظر ڈالی۔۔ یہ کچھ وجہیہ ضرور رہا ہوگا! مگر اب مٹاپے نے اس کے ساتھ جو سلوک کیا تھا اس کا اظہار الفاظ میں ناممکن ہے! بس دیکھنے اور محسوس کرنے کی چیز تھی! قد تو متوسط ہی تھا مگر پھیلاؤ نے اس توسط کی ریڑھ مار کر رکھ دی تھی! صرف کناروں پر تھوڑے سے سیاہ بال تھے جو اگر سفید ہوتے تو اتنے برے نہ معلوم ہوتے۔

اس کے پیروں کے پاس ہی ایک ننھا منا سا خوبصورت کتا بیٹھا سرخ زبان نکالے ہانپ رہا تھا! جولیا نے اسے تعریفی نظروں سے دیکھا۔ اس کے بال بڑے اور سفید تھے۔ کان البتہ گہرے کتھئی تھے اور یہی اس کا حسن تھا۔

"سرسوکھے رام۔۔ اور مس جولیانافٹنرواٹر۔۔!" خاور نے تعارف کرایا۔

سرسوکھے رام نے مسکرا کر سر کو خفیف سی جنبش دی۔

اور جولیانے ہاتھ ملتے ہوئے کہا۔ " میں آپ کی کیا خدمت کرسکتی ہوں جناب!"

وہ دل ہی دل میں ہنس رہی تھی۔ اتنی اردو تو سمجھتی ہی تھی کہ اس کے نام اور حبثہ کے تضاد سے لطف اندوز ہوسکتی!۔۔ 
کتنی ستم ظریفی تھی! یہ ہاتھی سا آدمی سوکھے رام کہلاتا تھا۔۔ یہی نہیں بلکہ خطاب یافتہ بھی تھا! وہ سوچ رہی تھی نہ ہوا عمران ورنہ مزہ آجاتا۔

"دیکھیئے بات دراصل یہ ہے کہ میں مستقل طور پر آپ لوگوں سے معاملہ کرنا چاہتا ہوں"۔ سرسوکھے نے کہا۔

"ہم ہر خدمت کے لئے حاضر ہیں"۔

" وہ۔۔ تو۔۔ تو۔۔ ٹھیک ہے"۔ سرسوکھے نے کرسی کی پشت سے ٹکنے کی کوشش کرتے ہوئے کہا! "مگر آپ کو اس سلسلہ میں تھوڑی سی دردسری بھی مول لینی پڑے گی! دیکھیئے بات دراصل یہ ہے۔۔!"

وہ سانس لینے کے لئے رک گیا اور جولیا جھک کر اس کے کتے کا سر سہلاتے ہوئی بولی۔ "بڑا پیارا کتا ہے!"

سرسوکھے نے اس طرح چونک کر کتے کی طرف دیکھا جیسے اس کی موجودگی کا خیال ہی نہ رہا ہو۔

" آپ کو پسند ہے!" اس نے مسکرا کر پوچھا۔

"بہت زیادہ۔۔"

"تو میری طرف سے قبول فرمائیے۔۔!"

"اوہ۔۔ ارے نہیں۔۔!" جولیا خواہ مخواہ ہنس پڑی۔

"نہیں! اب میں اسے اپنے ساتھ نہیں لے جاؤں گا۔ سرسوکھے نے کہا اور کتے سے بولا۔ " لکی۔۔ یہ دیکھو اب یہ تمہاری مالکہ ہیں"۔

وہ دم ہلانے لگا اور سرسوکھے نے پھر اپنے بزنس کی بات شروع کردی۔

"قصہ دراصل یہ ہے کہ۔۔ اوہ ٹھہرئیے میں پہلے اپنا پورا تعارف تو کرادوں! میری فرم کا نام "سوکھے انٹرپرائزس" ہے۔۔!"

"اوہ۔۔ اچھا میں سمجھ گئی۔۔!"

"آپ جانتی ہیں!" وہ خوش ہو کر بولا۔ "خیر تو۔۔ میرا فارورڈنگ اور کلیرنگ کا الگ سے اسٹاف تھا! لیکن اب اس پر غیر ضروری مصارف ہونے لگے تھے! میں نے حساب لگایا تو اندازہ ہوا کہ اگر یہ کام کسی دوسری فرم کے سپرد کردیا جائے تو نسبتاًسستے میں ہوگا"۔

"جی ہاں۔۔ عموماً یہیں ہوتا ہے۔۔" جولیا سرہلا کر بولی۔

"بس تو پھر میں نے اپنے یہاں وہ سیکشن توڑ دیا ہے! "سرسوکھے نے کہا۔ "اور اب اس کے لئے آپ کی فرم سے معاملات طے کرنا چاہتا ہوں"۔

"غالباً مینجر صاحب آپ کو یہاں کے قواعد وضوابط سے آگاہ کرچکے ہیں"۔

"جی ہاں۔ اور میں ان سے کلی طور پر متفق ہوں۔ سرسوکھے نے کہا۔ "قواعد وضوابط کی بات نہیں تھی! میں تو دراصل آپ کے لئے تھوڑی سی دردسری بڑھانا چاہتا ہوں۔۔!"

"فرمائیے۔۔!"

"آپ کو ایک ایسا حساب بھی تیار کرنا ہوگا جس سے یہ ظاہر ہو کہ یہ کام میری ہی فرم کے ایک سیکشن نے کیا ہے"۔

خاور نے جولیاکی طرف دیکھا! اور جولیا جلدی سے بولی"یہ کوئی ایسی بات نہیں ہے جس کے لئے آپ کو زیادہ تشویش ہو۔ ایسا بھی ہوجائے گا"۔

"بس تو پھر ٹھیک ہے! کیا آپ کسی وقت میرے دفتر آنے کی زحمت گوارا کرسکتی ہیں؟"

"جب آپ فرمایئے۔۔!"

"نہیں بھئی جب آپ کو فرصت ملے۔ بس آنے سے پہلے فون کردیئے گا"۔

"بہتر ہے! میں آ کر دیکھ لوں گی کہ اب تک آپ کے یہاں حسابات کس طرح رکھے جاتے رہے ہیں"۔

"اوہ۔ شکریہ! یہ تو بڑی اچھی بات ہوگی! اس کے لئے آپ جو بھی حق المحنت تجویز کریں مجھے اس پر اعتراض نہ ہوگا۔۔!"

"حق المحنت کیسا"۔ جولیا نے حیرت سے کہا! "یہ تو میں اپنی فرم کے انٹرسٹ میں کروں گی۔ ہمارے لئے یہی کیا کم ہے کہ ہمیں اتنا بڑا اور مستقل کام مل رہا ہے"۔

"یہی بات۔۔!" سرسوکھے نے میز پر اس طرح گھونسہ مار کر کہا کہ اس کا سارا جسم تھلتھلا گیا!" یہی بات۔۔ یہی اسپرٹ کام کرنے والوں میں ہونی چاہیئے"۔ پھر خاور سے بولا۔ " آپ خوش قسمت ہیں جناب کہ اتنے اچھے ساتھی آپ کے حصے میں آئے ہیں!"

"شکریہ۔۔" خاور نے سگار کا ڈبہ اسے پیش کیا۔

"بس جناب! اب اجازت دیجیئے!"۔۔ وہ اٹھتا ہوا بولا۔ پھر جولیا سے کہا۔ "میں آپ کا منتظر رہوں گا"۔۔ ساتھ ہی دم ہلاتے کتے سے بولا۔ "نہیں لکی تم میرے ساتھ نہیں جاسکتے! تمہاری مالکہ وہ ہیں!"

کتا جولیا کی طرف مڑا اور وہ متحیر رہ گئی کیونکہ اب وہ اس کی کرسی پر دونوں اگلے پنجے ٹیک کر کھڑا ہوگیا تھا اور اس کی ران سے اپنی تھوتھنی رگڑ رہا تھا!

اس نے پھر اس کے سر پر ہاتھ پھیرا اور اس کی ننھی سی دم بڑی تیزی سے ہلنے لگی۔

"کمال ہے!۔۔ جولیا اور خاور نے بیک وقت کہا۔

"کتوں کو ٹرینڈ کرنا میری ہابی ہے"۔ سرسوکھے مسکرایا۔ "میرے سارے کتے بڑے سمجھدار ہیں! اب یہ میرے ساتھ واپس جانے کی کوشش نہیں کرے گا۔ اور صرف آپ ہی کے ساتھ جائے گا! آپ کے دفتر کا کوئی دوسرا آدمی اسے اپنے ساتھ نہیں لے جاسکتا۔۔ اچھا بس اجازت دیجیئے!۔۔"

وہ ان دونوں سے مصافحہ کرکے رخصت ہوگا۔ اس کی چال بھی عجیب تھی بس ایسا معلوم ہورہا تھا جیسے کوئی گیند اچھلتا کودتا ہوا چل پڑا ہو۔

"کیا خیال ہے۔۔!" اس کے چلے جانے کے بعد خاور نے جولیا کی طرف دیکھا۔

"حیرت انگیز۔۔"

"ہر اعتبار سے۔۔ ہماری بڑی بدقسمتی ہے کہ اس شہر میں ایسے ایسے عجوبے موجود ہیں لیکن ہمیں ان کے دیدار نہیں ہوتے۔۔ تم نے اس کی چال پر غور کیا؟"

"ہاں! وہی تو میرے لئے حیرت انگیز تھی۔ میں سوچ بھی نہیں سکتی تھی کہ اتنا موٹا آدمی اتنی تیز رفتاری سے چل سکے گا"۔
"اس کی آنکھیں کتنی چمکیلی ہیں"۔ خاور نے کہا۔

"اور یہ کتا۔۔" جولیا نے کتے کی طرف دیکھ کر کہا۔ جو اب اس کے پیروں کے قریب بیٹھا زبان نکالے ہانپ رہا تھا۔۔!

\endR
\bigskip\bigskip\bigskip

\centerline{\hdr{(۲)}}
\bigskip\bigskip\bigskip
\beginR
\body

جوزف رانا پیلس ہی کا ہو کر رہ گیا تھا! آتشدان کا بت والے کیس کے بعد اس نے فلیٹ کی شکل نہیں دیکھی تھی۔ عمران کی تاکید تھی کہ وہ ادھر کا رخ بھی نہ کرے۔۔! اس طرح سلیمان یہ فیصلہ کرنے کے قابل ہوسکا تھا کہ وہ بدستور عمران ہی کی خدمت کرتا رہے گا۔

رانا پیلس میں سب ہی تھے۔ نوکر چاکر، ڈرائیور، جوزف۔ حتیٰ کہ بلیک زیروبھی (بوڑھے آدمی کے میک اپ میں)۔ لیکن رانا تہور علی صندوقی کا کہیں پتہ نہ تھا۔۔!

بلیک زیرو بوڑھے طاہر صاحب کے روپ میں رانا تہور علی صندوقی کا منیجر تھا، سمجھا جاتا تھا کہ وہ اس کی جائیداد کی دیکھ بھال کرتا ہے۔

جوزف ہر وقت فوجی وردی میں رہتا تھا اور اس کے دونوں پہلوؤں سے ریوالور لٹکے رہتے تھے۔ اس کا خیال تھا کہ فوجی وردی میں اس کی مارشل اسپرٹ ہر وقت بیدار رہتی ہے اور شراب نہ ہونے پر اسپرٹ ہی میں پانی ملا کر پینے سے بھی نہیں مرتی۔۔!

جوزف بلانوش تھا! لیکن اسے معینہ مقدار سے زیادہ شراب نہیں ملتی تھی اس لئے وہ اکثر اسپرٹ میں پانی ملا کر پیا کرتا تھا۔۔!

اس وقت وہ اسپرٹ کے نشے کی جھونک میں پورچ میں "اٹینشن" تھا!۔بالکل کسی بت کی طرح بےحس وحرکت۔ پلکیں ضرور جھپکتی رہتی تھیں۔ مگر بالکل ایسا ہی معلوم ہوتا تھا جیسے کسی الّو کو پکڑ کر دھوپ میں بٹھا دیا گیا ہو۔۔! اور وہ خاموشی سے مجسم احتجاج بن کر تن بہ تقدیر ہوگیا ہو۔۔!

دفعتاً ایک آدمی پشت پر ایک بہت بڑا تھیلا لادے ہوئے پھاٹک میں داخل ہوا لیکن جوزف کی پوزیشن میں کوئی فرق نہ آیا بلکہ وہ تو اس کی طرف دیکھ بھی نہیں رہا تھا!۔۔

مگر جیسے ہی وہ پورچ کے قریب آیا۔ اچانک جوزف دہاڑا۔

"ہالٹ۔۔!"

اور وہ آدمی بھڑک کر دوچار قدم کے فاصلے پر تھیلے سمیت ڈھیر ہوگیا!

"گٹ اپ۔۔!" جوزف اپنی جگہ سے ہلے بغیر پھر دہاڑا۔۔

"ارے مار ڈالا۔۔!" وہ مفلوک الحال آدمی دونوں ہاتھوں سے سر تھام کر کراہا!

"کی در جاتا!"۔۔ جوزف غرایا۔

"بھیر جاتا۔۔ رانا صاحب کے پاس۔۔ ایسی ایسی جڑی بوٹیاں ہیں میرے پاس!۔۔"

"کیا باکتا۔۔!" جوزف پھر غرایا!

"آؤں۔۔ آجاؤں۔۔ پاس آجاؤں!" وہ آدمی خوف زدہ انداز میں ہاتھ ہلا ہلا کر پوچھتا رہا۔

اب جوزف خود ہی اپنی جگہ سے ہلا اور وہ آدمی تھیلا سمیٹتا ہوا پیچھے پھدک گیا! یہ دبلا تپلا اور چیچڑ جسم والا ایک بوڑھا آدمی تھا۔ آنکھیں اندر کو دھنسی ہوئی اور دھندلی تھیں!۔۔ لیکن ہاتھ پاؤں میں خاصی تیزی معلوم ہوتی تھی۔

"کیا بکتا۔۔!" جوزف اس کے سر پر پہنچ کر دہاڑا۔

"شش شش۔۔ شقاقل۔۔ مصری!" وہ تھیلےسے کوئی چیز نکال کر اسے دکھاتا ہوا پیچھے کھسکا!

"یہ کیا ہائے۔۔!" جوزف غرایا۔

"اجی بس۔۔ کیا بتاؤں۔۔" وہ بہت تیزی سے بول رہا تھا! رر ۔۔ رانا صاحب قدر کریں گے"۔

"رانا صاحب نائیں ہائے۔۔ بھاگ جیاؤ۔۔!

"تو آپ ہی لڑائی کیجیئے صاحب۔۔ مزہ آجائے گا۔۔ جڑی بوٹیاں۔۔ ہا ہا۔۔ رانا صاحب کہاں ہیں!"

"ام نائیں۔۔ جیان تا۔۔ جیاؤ۔۔!"

اتنے میں بلیک زیرو شور سن کر باہر آگیا۔

"کیا بات ہے۔۔" اس نے جوزف سے انگریزی میں پوچھا!

"باس کو پوچھتا ہے! میں کہتا ہوں باس نہیں ہیں! وہ مجھے کوئی چیز دکھاتا ہے"۔

بلیک زیرو نے بوڑھے کی طرف دیکھا! وہ جھک جھک کر اسے سلام کر رہا تھا۔

"حضور۔۔ حضور۔۔ حضور عالی۔۔ سرکار۔ جڑی بوٹیاں ہیں میرےپاس۔ بڑی دور سے رانا صاحب کا سن کر آیا ہوں"۔

بلیک زیرو نے جلدی میں کچھ سوچا اور آہستہ سے بولا۔ "ہاں کہو ہم سن رہے ہیں"۔

"جو کچھ کہیئے۔ حاضر کروں سرکار۔۔!"

"ہم کیا کہیں! ہم نے تمہیں کب بلایا تھا؟"

"سرکار حضور۔۔ رانا صاحب بڑے معرکے کی بوٹیاں ہیں۔ بس طبعیت خوش ہوجائے گی!"

"کیا ہمارے کسی دوست نے تمہیں بھیجا ہے؟"

"جی حضور۔۔ ہم نے اس سرکار کی بڑی تعریف سنی ہے!"

"خیر اندر چل کر۔۔ ہمیں کچھ بوٹیاں دکھاؤ! اور ان کے خواص بتاؤ"۔

بوڑھا خوش نظر آنے لگا تھا اس نے تھیلا سمیٹ کر کاندھے پر رکھا اور بلیک زیرو کے پیچھے چلنے لگا۔

جوزف کھڑا احمقانہ انداز میں پلکیں جھپکاتا رہا!۔۔ پھر یک بیک وہ چونک کر اس بوڑھے آدمی کے پیچھے جھپٹا!

بلیک زیرو اور بوڑھا آدمی اندرداخل ہوچکے تھے! بلیک زیرو اسے ایک کمرے میں بٹھانے کا ارادہ کر ہی رہا تھا کہ اس نے جوزف کو اس پر جھپٹے دیکھا!۔۔

"ارے۔۔ ارے حضور"۔ بوڑھا بوکھلا گیا۔

بلیک زیرو بھی بھونچکا رہ گیا!۔۔

لیکن بوڑھا دوسرے ہی لمحے میں زمین پر تھا! اور جوزف نے اس کی میلی اور سال خوردہ پتلون کی جیب سے ایک چھوٹا سال پستول نکال لیا تھا۔۔!

بوڑھا اس اچانک حملے سے بری طرح بوکھلا گیا تھا۔ اس لئے جوزف کی گرفت سے آزاد ہونے کے بعد بھی اسی طرح بےحس وحرکت پڑا رہا البتہ اس کی آنکھیں کھلی ہوئی تھیں اور وہ پلکیں بھی جھپکا رہا تھا۔

"کیوں! تم کون ہو!۔۔" بلیک زیرو نے آنکھیں نکال کر بولا۔

"مم۔۔ میں نہیں جانتا صاحب!۔۔ کہ یہ خطرناک ۔۔ چیز میری جیب میں کس نے ڈالی تھی"۔ وہ ہانپتا ہوا بولا۔

"بکواس مت کرو"۔ بلیک زیرو غرایا! تم کون ہو؟"

"جی میں جڑی بوٹیاں تلاش کرکے بیچتا ہوں۔۔ شوقین رئیس میری قدر کرتے ہیں"۔

"مگر تم پہلے تو کبھی یہاں نہیں آئے۔۔!" بلیک زیرو اسے گھورتا ہوا بولا۔

"جی بےشک میں پہلے کبھی نہیں آیا"۔

"کیوں نہیں آئے تھے؟" بلیک زیرو نے غصیلے لہجے میں کہا! اس کے ذہن میں اس وقت عمران رینگنے لگا تھا اور اس نے یہ سوال بالکل اسی کے سے انداز میں کیا تھا۔

"جج۔۔ جی۔۔ ای۔۔ کیا بتاؤں مجھے اس سرکار کا پتہ نہیں معلوم تھا! وہ تو ابھی ابھی ایک صاحب نے سڑک ہی پر بتایا تھا کہ اس محل میں جاؤ۔ یہاں رانا صاحب رہتے ہیں! بہت بڑی سرکار ہے!۔۔"

"اس پستول کی بات کرو۔۔!"

"صص۔۔ صاحب! میں نہیں جانتا! بھلا میرے پاس پستول کا کیا کام! پتہ نہیں کس نے کیوں یہ حرکت کی ہے۔ میں کچھ نہیں جانتا!۔۔ خدا کے لئے ان کالے صاحب کو یہاں سے ہٹا دیجیئے ورنہ میرا دم نکل جائے گا"۔

جوزف اسے خونخوار نظروں سے گھورتا ہوا بڑبڑا رہا تھا۔ "مسٹرٹائر۔ یہ کیا کہہ رہا ہے! مجھے بھی بتائیے"۔

"اس کو گردن سے پکڑ کر ٹانگ لو"۔ بلیک زیرو نے کہا۔

جوزف پستول کو بائیں ہاتھ میں سنبھال کر اس کی طرف بڑھا! لیکن اچانک ایسا معلوم ہواجیسے آنکھوں کے سامنے بجلی سی چمک گئی ہو! بوڑھا چکنے فرش پر پھسلتا ہوا کمرے سے باہر نکل گیا تھا۔

"خبردار فائر نہ کرنا جوزف۔۔" بلیک زیرو چیخا۔

جوزف نے بوڑھے پر چھلانگ لگائی تھی اور اب فرش سے اٹھ رھا تھا کیونکہ بوڑھا تو چھلاوہ تھا چھلاوہ۔

جب تک جوزف اٹھتا وہ بیرونی برآمدے میں تھا۔

"فائر مت کرنا۔ بلیک زیرو پھر چیخا! ساتھ ہی اب وہ بھی تیزی دکھانے پر آمادہ ہوگیا تھا جوزف کو پھلانگتا ہوا وہ بھی بیرونی برآمدے میں آیا۔

یہاں دو ملازم کھڑے چیخ رہے تھے۔

"۔۔ صاحب وہ چھت پر ہے"۔ دونوں نے بیک وقت کہا۔

بلیک زیرو چکرا گیا! بھلا یہ کیسے ممکن تھا کہ وہ اتنی جلدی چھت پر بھی پہنچ جاتا!۔۔

نوکروں نے قسمیں کھا کر یقین دلایا کہ انہوں نے اسے بندروں کی سی پھرتی سے اوپر جاتے دیکھا ہے۔ انہوں نے گندے پانی کے ایک موٹے پائپ کی طرف اشارہ کیا تھا۔ جس سے ملی ہوئی پورچ کی کارنس تھی اور پورچ کی چھت بہت زیادہ اونچی نہیں تھی کوئی بھی پھرتیلا آدمی کم از کم پورچ کی چھت تک تو اتنے وقت میں پہنچ ہی سکتا تھا۔

پھر ذرا ہی سی دیر میں پوری عمارت چھان ماری گئی لیکن اس کا کہیں پتہ نہیں تھا!۔۔

اندر پہنچ کر بلیک زیرو نے محسوس کیا کہ اس چھلاوے نے اپنا تھیلا بھی نہیں چھوڑا تھا۔

"ٹائر صاب"۔ جوزف نے غصیلی آواز میں کہا۔ "مجھے فائر کرنے سے کیوں منع کیا تھا؟"

"باس کا حکم ہے کہ اس محل میں کبھی گولی نہ چلائی جائے"۔

"چاہے کوئی یہاں آکر جوزف دی فائٹر کے منہ پر تھوک دے"۔

"خاموش رہو! باس کے حکم میں بحث کی گنجائش نہیں ہوا کرتی"۔

جوزف فوجیوں کے سے انداز میں اسے سلیوٹ کرکے اپنے کمرے کی طرف مڑ گیا۔ اس کا موڈ خراب ہوگیا تھا اس لئے وہ شراب کی بوتل پر ٹوٹ پڑا۔۔

\endR

\bigskip\bigskip\bigskip
\centerline{\hdr{(۳)}}
\bigskip\bigskip\bigskip
\beginR
\body


آج صفدر تین دن بعد آفس میں داخل ہوا تھا۔ مگر اس حال میں کہ اس کے بال گردآلود تھے۔ لباس میلا اور شیو بڑھا ہوا تھا۔
دوسروں نے اسے حیرت سے دیکھا! اور اس نے ایک بہت بری خبر سنائی!

"عمران مار ڈالا گیا!"

اور یہ خبر بم کی طرح ان پر گری! جولیا تو اس طرح اچھلی جیسے اس کی کرسی میں اچانک برقی رو دوڑا دی گئی ہو!

"کیا بک رہے ہو!۔۔ اس نے کانپتے ہوئے سسکی سی لی۔

وہ سب صفدر کے گرد اکھٹے ہوگئے! اس وقت یہاں صرف سیکرٹ سروس کے آدمی تھے۔ چونکہ چھٹی کا وقت ہوچکا تھا اس لئے دوڑ دھوپ کے کام کرنے والے جاچکے تھے۔

"ہاں! یہ حادثہ مجھے زندگی بھر یاد رہے گا!" صفدر بھرائی ہوئی آواز میں بولا۔ میں تین دن سےاس کے ساتھ ہی تھا! ہم دونوں کیپٹن واجد والی تنظیم کے بقیہ افراد کی فکر میں تھے۔ تین دن سے ایک آدمی پر نظر تھی! آج اس کا تعاقب کرتے ہوئے ندی کی طرف نکل گئے! مقبرے کے پاس جو سرکنڈوں کی جھاڑیاں ہیں وہاں ہمیں گھیر لیا گیا! حملہ اچانک ہوا تھا! پھر یہ بات میری سمجھ میں آئی کہ ہمیں دھوکے میں رکھا گیا تھا! ہم تو دراصل یہ سمجھتے رہے تھے کہ اس تنظیم کا ایک آدمی ہماری نظروں میں آگیا ہے لیکن حقیقت یہ تھی کہ وہ ہمیں نہایت اطمینان سے ختم کرنا چاہتے تھے۔ کسی ایسی جگہ گھیرنا چاہتے تھے جہاں سے بچ کر ہم نکل ہی نہ سکیں یعنی انہوں نے بھی وہ طریقہ اختیار کیا تھا جسے واجد کو پکڑنے کے لئے عمران کام میں لایا تھا"۔

"پھر کیا ہوا۔۔ باتوں میں نہ الجھاؤ!" جولیا مضطربانہ انداز میں چیخی۔

"ہم پر چاروں طرف سے فائرنگ ہو رہی تھی اور ہم کھلے میں تھے۔ اچانک میں نے عمران کی چیخ سنی۔ وہ ٹیکرے سے ندی میں گر رہا تھا! میں نے اسے گرتے اور غرق ہوتے دیکھا۔ تم جانتے ہی ہو کہ ندی کا وہ کنارہ کتنا گہرا ہے جس کنارے پر مقبرہ ہے۔۔!"

"تم کیسے بچ گئے؟"

"بس موت نہیں آئی تھی!" صفدر نے پھیکی سی مسکراہٹ کے ساتھ کہا۔

"تب تو پھر تم آفس ناحق آئے۔۔! تمہیں ادھر کا رخ ہی نہ کرنا چاہیئے تھا! جاؤ جتنی جلدی ممکن ہو اپنی قیام گاہ پر پہچنے کی کوشش کرو"۔

جولیا میز سے ٹکی کھڑی تھی۔ اس کا سر چکرارہا تھا!

"نہیں میں یقین نہیں کرسکتی!۔۔ کبھی نہیں"۔ وہ کچھ دیر بعد ہذیانی انداز میں بولی۔ " عمران نہیں مر سکتا! بکواس ہے۔ کبھی نہیں! تم جھوٹے ہو!"

وہ خواہ مخواہ ہنس پڑی! اس میں اس کے ارادے کو دخل نہیں تھا!۔۔

وہ سب اسے عجیب نظروں سے دیکھنے لگے۔ ان میں تنویر بھی تھا۔

"مرنے کو تو ہم سب ہی اسی وقت مرسکتے ہیں!" اس نے کہا۔

"ہم سب مرسکتے ہیں! مگر عمران نہیں مرسکتا! اپنی بکواس بند کرو"۔

پھر جولیا نے کانپتے ہوئے ہاتھ سے ایکس ٹو کے نمبر ڈائیل کئے لیکن دوسری طرف سے جواب نہ ملا!

" تمہیں سرسوکھے کے ہاں جانا تھا"۔ خاور نے کہا۔

"جہنم میں گیا سرسوکھے"۔ جولیا حلق پھاڑ کر چیخی۔ " کیا تم سب پاگل ہوگئے ہو۔ کیا عمران کا مرجانا کوئی بات ہی نہیں ہے!"

"اس کی موت پر یقین آجانے کے بعد ہی ہم سوگ مناسکیں گے!" خاور نے پھیکی سی مسکراہٹ کے ساتھ کہا۔

دفعتاً لیفٹننٹ چوہان نے صفدر سے سوال کیا! "تمہیں وہ آدمی ملا کہاں تھا!۔۔ اور تمہیں یقین کیسے آیا تھا کہ وہ اسی تنظیم سے تعلق رکھتا ہے"۔

"عمران نے مجھے یہ نہیں بتایا تھا"۔

"آخر وہ تمہیں ہی کیوں ایسے مہمات کے لئے منتخب کرتا ہے؟"

"وہ کیوں کرنے لگا! مجھے ایکسٹو کی طرف سےہدایت ملی تھی۔۔!"

وہ سب پھر خاموش ہوگئے۔ جولیا میز پر سر ٹیکے بیٹھی تھی! اور تنویر غصیلی نظروں سے اسے دیکھ رہا تھا!

پھر وہ اٹھی اوراپنا بیگ سنبھال کر دروازے کی طرف بڑھی۔

"تم کہاں جارہی ہو؟" تنویر نے اسے ٹوکا۔

"شٹ اپ۔۔" وہ مڑ کر تیز لہجے میں بولی۔ "میں ایکسٹو کے علاوہ اور کسی کو جواب دہ نہیں ہوں"۔

وہ باہر نکل کر اپنی چھوٹی سی ٹوسیٹر میں بیٹھ گئی! لیکن وہ نہیں جانتی تھی کہ اسے کہاں جانا ہے۔۔!

صفدر کو وہ ایک دیانت دار اور سنجیدہ آدمی سمجھتی تھی۔ اس قسم کی جھوٹ کی توقع اس کی ذات سے نہیں کی جاسکتی! اس نے سوچا ممکن ہے عمران نے اسے بھی ڈاج دیا ہو!۔۔ لیکن کیا ضروری ہے کہ وہ ہمیشہ بچتا ہی رہے۔

کچھ دیر بعد ٹوسیٹر ایک پبلک فون بوتھ کے قریب رکی اور بوتھ میں آکر عمران کے نمبر ڈائیل کئے! اور دوسری طرف سے سلیمان نے جواب دیا! لیکن اس نے عمران کے متعلق لاعلمی ظاہر کرتے ہوئے بتایا کہ وہ پچھلے تین دنوں سے گھر نہیں آیا۔

جولیا نے سلسلہ منقطع کرتے ہوئے ٹھنڈی سانس لی۔

کیسے معلوم ہو کہ صفدر کا بیان کہاں تک درست ہے! آخر یہ کمبخت کیوں بچ گیا! پھر ذرا ہی سی دیر میں اسے ایسا محسوس ہونے لگا جیسے صفدر ہی عمران کا قاتل ہو!۔۔

پھر اس نے غیر ارادی طور پر اپنی گاڑی ندی کی طرف جانے والی سڑک پر موڑ دی۔۔

سورج غروب ہونے والا تھا۔ مگر وہ دن رہے وہاں پہنچنا چاہتی تھی اس لئے کار کی رفتار خاصی تیز تھی۔ گھاٹ کی ڈھلان شروع ہوتے ہی اس نے بائیں جانب والے ایک کچے راستے پر گاڑی موڑ دی۔ اسی طرف سے وہ اس ٹیکرے تک پہنچ سکتی تھی جہاں ایک قدیم مقبرہ تھا۔ اور دور تک سرکنڈوں کا جنگل پھیلا ہواتھا۔

کچے راستے کی دونوں جانب جھنڈ بیریوں سے ڈھکے ہوئے اونچے اونچے ٹیلے تھے۔

مقبرے تک گاڑی نہیں جاسکتی تھی کیونکہ وہاں تک پہچنے کا راستہ ناہموار تھا! اس نے گاڑی روکی، انجن بند کیا اور نیچے اتر کر خالی خالی آنکھوں سے افق میں دیکھتی رہی جہاں سورج آسمان کو چھوتی ہوئی درختوں کی قطار کے پیچھے جھک چکا تھا!

پھر وہ چونکی اور مقبرے کی طرف چل پڑی۔

ابھی دھندلکا نہیں پھیلا تھا!۔۔ دریا کی سطح پر ڈھلتی ہوئی روشنی کے رنگین لہرئے مچل رہے تھے۔۔ وہ ٹیکرے کے سرے کی جانب بڑھتی چلی گئی!

مگر کیا یہ حماقت ہی نہیں تھی!۔۔ اس نے سوچا! آخر وہ یہاں کیوں آئی ہے؟

ٹیکرے کے نیچے پانی پر ایک موٹر بوٹ نظر آئی جس میں کوئی نظر نہیں آرہا تھا! ہوسکتا ہے کچھ لوگ اس کے چھوٹے سے کیبن میں رہے ہوں۔

اچانک موٹر بوٹ سے ایک فائر ہوا۔ پانی پر ایک جگہ بلبلے اٹھے تھے اور گولی بھی ٹھیک اسی جگہ پڑی تھی۔

کیبن کی کھڑکی سے رائفل کی نال پھر اندر چلی گئی اور اس کے بعد ایک آدمی سر نکال کر پانی کی سطح پر دیکھنے لگاجہاں ایک بڑی سی مردہ مچھلی ابھر آئی تھی!

پھر کیبن کی دوسری کھڑکی سے ایک سیاہ رنگ کا بڑا سا کتا پانی میں کودا اور تیرتا ہوا مچھلی تک جاپہنچا! اس کی دم منہ میں دبا کر وہ پھر موٹر بوٹ کی طرف مڑا تھا۔

دوسری بار جب موٹر بوٹ میں بیٹھے ہوئے آدمی نے اپنے دونوں ہاتھ کھڑکی سے نکال کر مچھلی کو سنبھالا۔ اس وقت جولیا نے اسے صاف پہچان لیا! وہ سرسوکھے تھا!

اس نے مچھلی اندر کھینچ لی اور کتا بھی کھڑکی سے کیبن میں چلاگیا۔

تو وہ مچھلیوں کا شکار کھیل رہا تھا۔۔ جولیا ٹیکرے سے پرے کھسک آئی۔ اس نے سوچا اچھا ہی ہوا سرسوکھے کی نظر اس پر نہیں پڑی! ورنہ خواہ مخواہ تھوڑی دیر تک رسمی قسم کی گفتگو کرنی پڑتی! مگر اب وہ یہاں کیوں ٹھہرے! آئی ہی کیوں تھی؟ یہاں کیا ملتا!۔۔ اگر عمران مارا بھی گیا تو۔۔! وہ۔۔ وہ یک بیک چونک پڑی!اگر وہ یہاں مارا گیا ہوگا تو ایک آدھ بار لاش سطح پر ضرور ابھری ہوگی! مگر اسے کیا؟ ضروری نہیں ہے کہ کسی نے اسے دیکھا بھی ہو!۔۔

پھر وہ کیا کرے۔۔ کیا کرے۔۔!

غیر ارادی طور پر وہ سرکنڈوں کی جھاڑیوں میں گھس پڑی! یہ ایک پتلی سی پگڈنڈی تھی جو سرکنڈوں کی جھاڑیوں سے گذر کر کسی نامعلوم مقام تک جاتی تھی۔

کچھ دور پر اسے ریوالور کے چند خالی کارتوس پڑے ملے اور صفدر کے بیان کی تصدیق ہوگئی! ویسے وہ تو اس پر یوں بھی اعتماد کرتی تھی۔

مگر سوال یہ تھا کہ اب جولیا کیا کرے۔۔ یہ بات تو خود صفدر کو بھی نہیں معلوم تھی کہ عمران نے اس آدمی کو کہاں سے کھود نکالا تھا جس کے تعاقب میں وہ دونوں یہاں تک آئے تھے اور یہ حادثہ پیش آیا تھا!

اچانک کوئی چیز اس کی پشت سے ٹکرائی اور وہ اچھل پڑی۔ بس غنیمت یہی تھااس کے حلق سے کسی قسم کی آواز نہیں نکلی تھی ورنہ وہ چیخ ہی ہوتی۔

اس نے جھک کر اس کاغذ کو اٹھایا جو شاید کسی وزنی چیز پر لپیٹ کر پھیکا گیا تھا۔

کاغذ کی تہوں کے درمیان ایک چھوٹی سی کنکری تھی۔

کاغذ پر تحریر تھا:

"جولیا۔ دفع ہوجاؤ یہاں سے۔۔ کھیل مت بگاڑو!"

ایک بےساختہ قسم کی مسکراہٹ اس کے ہونٹوں پر پھیل گئی! دل پر سے بوجھ سا ہٹ گیا!۔۔ اور وہ تیزی سے واپسی کے لئے مڑ گئی! طرز تحریر عمران ہی کا سا تھا!

\endR

\bigskip\bigskip\bigskip
\centerline{\hdr{(۴)}}
\bigskip\bigskip\bigskip
\beginR
\body

واپسی بڑے سکون کے ساتھ ہوئی۔ جولیا کا دل چاہ رہا تھا کہ قہقہے لگائے، ہنستی ہی رہے!۔۔ لیکن وہ صرف ذہنی مسرت پر ہی قناعت کئے ہوئے کار ڈرائیو کرتی رہی!

گھر پہنچ کر اس نے ٹھنڈی پھواروں سے غسل کیا اور ڈریسنگ گاؤن پہنے ہوئے خواب گاہ میں چلی گئی۔ آج کی تھکن اسے بڑی لذت انگیز محسوس ہو رہی تھی!

اس نے ہیٹر پر چائے کے لئے پانی رکھتے ہوئے سوچا! اگر اس وقت آجائے عمران۔۔؟ اچھی طرح خبر لوں اس کی۔

دفعتاً فون کی گھنٹی بجی!

جولیا نے ہاتھ بڑھا کر رسیور اٹھا لیا۔

"ہیلو۔"

"ایکس ٹو۔" دوسری طرف سے آواز آئی۔

"یس سر۔"

"تم ندی کی طرف کیوں گئی تھیں؟"

"اوہ۔۔ جناب۔۔ وہ۔۔ عمران۔۔"

"ہاں مجھے علم ہے۔۔ مگر تم کیوں گئی تھیں؟"

"صص۔۔ صفدر۔۔!"

تمہارے علاوہ۔۔ اور کوئی کیوں نہیں گیا؟"

"پتہ نہیں جناب!" جولیا جھنجلا گئی۔

"وہ جانتے ہیں کہ انہیں اتنا ہی کرنا ہے جتنا کہا جائے۔۔!"

"یعنی میں۔۔ اس کی موت کی خبر سنتی۔۔ اور۔۔!"

"تجہیز و تکفین کی فکر نہ کرتی!" ایکس ٹو نے طنزیہ لہجے میں جملہ پورا کردیا!

"تم کون ہوتی ہو اس کی فکر کرنے والی! اپنی حدود سے باہر قدم نہ نکالا کرو"۔

"بہت بہتر جناب!" جولیا کسی سلگتی ہوئی لکڑی کی چٹخی!

"تمہارا لہجہ!۔۔ تم ہوش میں ہو یا نہیں!" ایکس ٹو اپنے مخصوص خونخوار لہجے میں غرایا!

"میرا محکمہ عشقیہ ڈراموں کے ریہرسل کے لئے نہیں ہے! سمجھیں۔۔!"

"جج۔۔ جی۔۔ ہاں"۔ جولیا بوکھلا گئی۔

دوسری طرف سے سلسلہ منقطع ہوگیا۔

وہ رسیور رکھ کر آرام کرسی کی پشت سے ٹک گئی۔ اس کی آنکھیں پھیلی ہوئی تھیں۔ اور دل بہت شدت سے دھڑک رہا تھا۔

پھر آہستہ آہستہ سکون ہوتا گیا اور اسے ایکس ٹو پر اس زور سے غصہ آیا کہ ذہنی طور پر ناچ کر رہ گئی۔۔ اسے کیا حق حاصل ہے۔ وہ کون ہوتا ہے۔ میرے نجی معاملات میں دخل دینے والا ظالم۔۔ کمینہ۔۔ ذلیل۔۔!

فون کی گھنٹی پھر بجی!

اس نے برا سا منہ بنا کر رسیور اٹھا لیا! اور "ہیلو" کہتے وقت بھی اس کا لہجہ زہریلا ہی رہا!

"مس فٹز واٹر پلیز۔۔!" دوسری طرف سے آواز آئی۔

"ہاں۔۔! جولیا نے بھرائی ہوئی آواز میں جواب دیا۔۔ وہ بولنے والے کی آواز نہیں پہچان سکی تھی۔

"میں سوکھے رام بول رہا ہوں!"

"اوہ۔۔! فرمائیے۔۔ جناب۔۔!"

"میں اس وقت اپنے آفس میں تنہا ہوں! کیا آپ تکلیف کریں گی"۔

"اس وقت!" جولیا نے حیرت سے کہا اور پھر کسی سوچ میں پڑ گئی!

"آپ نہیں سمجھ سکتیں مس فٹز واٹر۔۔ میں دراصل آپ کو اپنے اعتماد میں لینا چاہتا ہوں! میری بدنصیبی کی داستان طویل ہے"۔

"میں بالکل نہیں سمجھی! سرسوکھے۔۔ پلیز!"

"فون پر کچھ نہیں کہہ سکتا!"

"اچھا سر سوکھے! میں آرہی ہوں! مگر آپ کو میرے گھر کانمبر کیسے ملا؟"

"بس اتفاق ہی سے میں مچھلیوں کا شکار کھیل کر واپس آرہا تھا کہ آپ کے دفتر کے ایک صاحب نظر آگئے۔ انہوں نے اپنا نام بتایا تھا لیکن صرف صورت آشنائی کی حد تک میری یادداشت قابل رشک ہے! نام وغیرہ البتہ یاد نہیں رہتے بہرحال میں نے ان سے آپ کےمتعلق پوچھا تھا انہوں نے بتایا کہ آپ اس وقت گھر ہی پر ملیں گی۔ انہوں نے فون نمبر بھی بتایا !"

"خیر۔ میں آرہی ہوں"۔ جولیا نے کہا اور سلسلہ منقطع کرکے خاور کے نمبر ڈائیل کئے۔ وہ گھر ہی پر موجود تھا۔

"سر سوکھے مجھے اس وقت اپنے آفس میں طلب کر رہا!" جولیا نے کہا۔

"ضرور جاؤ۔۔ ذرہ برابر بھی ہچکچاہٹ نہ ہونی چاہیئے۔ تمہاری حفاظت کا انتظام بھی کردیا جائے گا"۔

"مگر میں نہیں سمجھ سکتی؟"

"ٹھہرو!" خاور نے جملہ پورا نہیں ہونے دیا۔ " ایکس ٹو کی ہدایت ہے کہ اگر آج کل کوئی نیا گاہک بنے تو اسے ہر ممکن رعایت دی جائے! میں سر سوکھے کا معاملہ اس کے علم میں لاچکا ہوں"۔

"اور اگر میں جانے سے انکار کردوں تو۔۔"؟

"میں اسے محض مذاق سمجھوں گا! کیونکہ تم ناسمجھ نہیں ہو"۔

جولیا نے اپنی اور سرسوکھے رام کی گفتگو دہراتے ہوئے کہا۔ "وہ آدمی اب تک میری سمجھ میں نہیں آیا۔۔!"

"پرواہ مت کرو! ایکس ٹو اس کے معاملے میں بہت زیادہ دلچسپی لے رہا ہے"۔

جولیا نے پھر برا سا منہ بنایا اور سلسلہ منقطع کردیا۔

تھوڑی دیر بعد پھر اس کی ٹوسیٹر شہر کے بارونق بازاروں میں دوڑ رہی تھی۔

تقریباً پندرہ منٹ بعد اس نے عمارت کے سامنے کار روکی جس کی دوسری منزل پر سرسوکھے انٹرپرائزس کا دفتر تھا۔ کھڑکیوں میں اسے روشنی نظر آئی۔ چوتھی یا پانچویں منزل کی بات ہوتی تو وہ لفٹ ہی استعمال کرتی! لیکن دوسری منزل کے لئے تو زینے ہی مناسب تھے!

سر سوکھے نے بڑی گرم جوشی سے اس کا استقبال کیا! لیکن جولیا محسوس کر رہی تھی کہ وہ کچھ خائف سا نظر آرہا ہے!
"بیٹھیئے بیٹھیئے! مس فٹز واٹر میں بےحد مسرور ہوں کہ آپ میری درخواست پر تشریف لائیں۔۔!" وہ ہانپتا ہوا بولا۔ جولیا ایک کرسی کھسکا کر بیٹھ گئی!

میں آپ کا زیادہ وقت نہیں برباد کروں گا مس فٹز واٹر!" سوکھے رام پھر بولا۔ "اوہ۔۔ ٹھہرئیے! آپ کیا پئیں گی۔ اس وقت تو میں ہی آپ کو سرو کروں گا کیوں کہ اس وقت یہاں ہم دونوں کے علاوہ اور کوئی بھی نہیں ہے"۔

"اوہ شکریہ! میں کسی چیز کی بھی ضرورت نہیں محسوس کر رہی اور پھر میں تو ویسے بھی شراب نہیں پیتی!"

"گڈ!۔۔" سرسوکھے کی آنکھیں بچگانے انداز میں چمک اٹھیں! وہ اسے تحسین آمیز نظروں سے دیکھتا ہوا بولا۔ "اگر آپ شراب نہیں پیتیں تو میں یہی کہوں گا کہ آپ پر اعتماد کیا جاسکتا ہے! بڑی پختہ قوت ارادی رکھتی ہیں وہ لڑکیاں جو شراب نہیں پیتیں"۔

"شکریہ ! جی ہاں میں بھی سمجھتی ہوں! خیر آپ کیا کہنے والے تھے؟"

جواب میں سرسوکھے نے پہلے تو ایک ٹھنڈی سانس لی اور پھر بولا۔ " میں نے اپنا فارورڈنگ اور کلیرنگ کا شعبہ بلا وجہ نہیں ختم کیا! میں مجبور تھا! نہ کرتا تو بہت بڑی مصیبت میں پڑ جاتا! لیکن ٹھہریئے ۔۔ میں آپ پر یہ بھی واضح کرتا چلو ں مس فٹز واٹر کہ آپ کو یہ سب باتیں کیوں بتا رہا ہوں! میں جانتا ہوں کہ عورتیں طبعاً رحم دل اور ہمدرد ہوتی ہیں"۔

وہ خاموش ہو کر کچھ سوچنے لگا! اور جولیا سوچنے لگی کہ اس گفتگو کا ماحصل کیا ہوگا جس کے سر پیر کا ابھی تک تو پتہ نہیں چل سکا!

"اوہ۔۔ میں خاموش کیوں ہوگیا!" سرسوکھے چونک کر بولا! پھر خفیف سی مسکراہٹ اس کے ہونٹوں پر نظر آئی اور اس نے کہا۔ 
"میری باتیں اکثر بےربط ہوجاتی ہیں مس فٹز واٹر! مگر ٹھہریئے میں ایک نقطے کی وضاحت کرنے کی کوشش کروں گا! میرے فاورڈنگ اینڈ کلیرنگ سیکشن میں کوئی بہت ہی بدمعاش آدمی آگھسا تھا اور ایسے انداز میں اسمگلنگ کررہا تھا کہ آئی گئی میرے ہی سرجاتی۔ لکڑی کی پیٹیوں میں باہر سے مال پیک ہو کر آتا تھا لیکن اس کے بعد پتہ نہیں چلتا تھا کہ خالی پیٹیاں کہاں غائب ہوجاتی تھیں!"

"میں نہیں سمجھی! "

"خالی پیٹیاں۔۔ غائب ہوجاتی تھیں!"

"تو اس کا یہ مطلب ہے کہ آپ فرم رٹیل بھی کرتی ہے۔۔!" جولیا نے حیرت سے کہا۔ پیٹیوں کا کھول ڈالا جانا تو یہی ظاہر کرتا ہے!"

"گڈ! آپ واقعئی ذہین ہیں! مجھ سے اندازے کی غلطی نہیں ہوئی"۔ سرسوکھے خوش ہو کر بولا! "میں ساری پیٹیوں کی بات نہیں کر رہا تھا۔ بلکہ میری مراد صرف ان بڑی پیٹیوں سے تھی جن میں مشینوں کے پرزے پیک ہو کر آتے ہیں! وہ پیٹیاں تو لامحالہ کھولی جاتی تھیں کیوں کہ ان مشینوں کی تیاری فرم ہی کراتی ہے! یعنی وہ یہیں اسمبل ہوتی ہیں"۔

"خیر۔۔ اچھا! "جولیا سرہلا کر بولی۔ " لیکن آپ خالی پیٹیوں کے متعلق کچھ کہہ رہے تھے!"

"وہ پیٹیاں غائب ہوجاتی تھیں"!

"اچھا چلیئے!" جولیا مسکرا کر بولی۔ "اگر وہ پیٹیاں غائب ہوجاتی ہیں تو اس میں پریشانی کی کیا بات ہے۔ کوئی غریب آدمی انہیں بیچ کر اپنا بھلا کرلیتا ہوگا"۔

"اوہ یہی تو آپ نہیں سمجھتیں مس فٹز واٹر۔۔ بات دراصل یہ ہے کہ وہ پیٹیاں فائیو پلائی وڈ کی ہوتی ہیں۔۔مطلب سمجھتی ہیں نا آپ۔۔ خیر میں شروع سے بتاتا ہوں! ۔۔ مجھے کبھی ان پیٹیوں کا خیال بھی نہ آتا۔ مجھے بھلا اتنی فرصت کہاں کہ کاروبار کی ذرا ذرا سی تفصیل ذہن میں رکھتا پھروں۔۔ بات دراصل یہ ہوئی کہ ایک دوران میں کوٹھی پر لکڑی کا کام ہو رہا تھا۔ ایک جگہ لکڑی کا پارٹیشن ہونا تھا! خیال یہ تھا کہ دیوار کے فریم میں ہارڈ بورڈ لگا دیا جائے۔ لیکن کسی نے فائیو پلائی وڈ کی ان پیٹیوں کا خیال دلا دیا! میں نے سوچا کہ ہارڈ بورڈ سے بہتر وہی رہے گی پلائی وڈ۔۔ لہذا میں اتفاق سے خود ہی گوڈاؤن کی طرف جا نکلا وہاں اسی دن کچھ پیٹیاں کھولی گئی تھیں۔ چوکیدار تنہا تھا اور وہ خود ہی پیٹیاں کھول کر ان میں سے پرزے نکال رہا تھا۔ مجھے بڑی حیرت ہوئی! کیونکہ یہ کام تو کسی ذمہ دار آدمی کے سامنے ہونا چاہیئے تھا اور پھر یہ چوکیدار کی ڈیوٹی نہیں تھی۔ میں نے اس سے اس کے متعلق استفسار کیا اور اس نے بوکھلا کر جواب دیا کہ گوڈاؤن انچارج نے اسے یہی ہدایت دی تھی!۔۔ میں نے سوچا کہ انچارج سے جواب طلب کروں گا۔ اور چوکیدار سے کہا کہ وہ ایک ٹھیلا لائے اور جتنی جتنی بھی پیٹیاں خالی ہوگئی ہیں انہیں کوٹھی میں بھجوادے۔۔ وہ ٹھیلا لینے کے لئے دوڑا گیا۔ لیکن پھر اس کی واپسی نہ ہوئی! اوہ۔۔ خوب یاد آیا مس فٹز واٹر۔۔ لکی تو ٹھیک ہے نا!۔۔ وہ ایک فرمانبردار کتا ہے۔۔ آپ کو یقیناً اس سے کوئی شکایت نہ ہوگی۔۔!"

"بہترین ہے۔۔!" جولیا نے کہا۔

"میرے پاس کئی قسم کے بہترین کتے ہیں! بہتری کمیاب نسلیں بھی ہیں! کسی دن کوٹھی آئیے آپ انہیں دیکھ کر بہت خوش ہوں گی"۔

"آپ یہ فرما رہے تھے کہ چوکیدار غائب ہوگیا۔۔"

"اوہ۔۔ دیکھیئے! بس اسی طرح ذہن بہک جاتا ہے! ہاں تو وہ مردودبھاگ گیا۔ میں نے ایک دوسرے گوڈاؤن کے چوکیدار سے ٹھیلا منگوایا! اس دوران میں، میں نے ایک پیٹی کا ڈھکن اٹھایا اور اندازہ کرنے لگا کہ وہ ہارڈ بورڈ سے بہتر ثابت ہوگا یا نہیں! اچانک اس کے ایک گوشے پر نظر رک گئی اور میری آنکھیں حیرت سے پھیل گئیں۔ جانتی ہیں! میں نے کیا دیکھا!۔۔ لکڑیوں کی پرت میں ایک پرت سونے کی بھی تھی! سونے کا پتر۔۔ اسے بڑی خوبصورتی سے لکڑی کے پرتوں کے درمیان جمایا گیا تھا۔۔ شائد پیٹی کی کیلیں نکالتے وقت ایک گوشے کی لکڑی ادھڑ گئی تھی اور پرت ظاہر ہوگئی تھی! میں نے فوراً ہی گودام میں تالا ڈال دیا اور کوٹھی پر فون کرکے چار معتبر اور مسلح چوکیدار وہاں طلب کتے اور انہیں ہدایت کردی کہ کسی کو گودام کے قریب بھی نہ آنے دیں!۔۔ میں آپ سے کیا بتاؤں مس فٹز واٹر! ان تختوں سے تقریباً اٹھائیس سیر سونا برآمد ہوا تھا!۔۔ لیکن میں نے کسی کو بھی اس کی خبر نہ ہونے دی۔ آپ خود ہی سوچیئے اگر یہ بات کھل جاتی تو کون یقین کرتا کہ سرسوکھے کے ہاتھ صاف ہیں! کون یقین کرتا!۔۔ گوڈاؤن انچارج سے پوچھ گچھ کی تو معلوم ہوا کہ ہمیشہ یہی ہوتا ہے چوکیدار کسی بڑے آفیسر کا حوالہ دے کہ اسے مطمئن کردیتا تھا! چونکہ اس سلسلے میں کبھی کوئی پوچھ گچھ نہیں ہوئی تھی اس لئے اس نے بھی اس پر دھیان نہیں دیا۔ اس طرح وہ ایک دردسری سے بچا رہتا تھا ورنہ اسے بھی کھولی جانے والی پیٹیوں کا باقاعدہ طور پر ریکارڈ رکھنا پڑتا! میں نے اس سے پہلے کی خالی پیٹیوں کے بارے میں پوچھا تو اس نے جنرل منیجر کی درجنوں چھٹیاں دکھائیں جن میں وقتاً فوقتاً خالی پیٹیاں طلب کی گئی تھیں! اس نےبتایا کہ کچھ کباڑی قسم کے لوگ آتے تھے اور پیٹیاں وصول کرکے رسیدیں دے جاتے تھے! اس نے رسیدیں بھی دکھائیں!۔۔ میں نے جنرل منیجر سے انکوائری کی! مگر اس نے چھٹیوں کے دستخط اپنے نہیں تسلیم کئے! اس پر میں نے ایک ایکسپرٹ کی خدمات حاصل کیں جس نے جنرل منیجر کے بیان کی تصدیق کردی! یعنی وہ دستخط سچ مچ جعلی تھے! بس یہیں سے انکوائری کا خاتمہ ہوگیا! میں اب کس کے گریبان میں ہاتھ ڈالتا!"۔۔

"آپ نے پولیس کو اطلاع دی ہوتی!" جولیا نے کہا۔

"شائد آپ میری دشواریوں کو ابھی تک نہیں سمجھیں! یقین کیجیئے کہ میں قانونی معاملات میں بےحد ڈرپوک قسم کا آدمی ہوں۔ اگر کہیں پولیس نے الٹا مجھ پر ہی نمدہ کس دیا تو کیا ہوگا؟ میں تو کسی کو منہ دکھانے کے قابل بھی نہ رہوں گا! اوہ مس فٹز واٹر۔ بہرحال مجھے اپنے فاورڈنگ اینڈ کلیرنگ کے عملہ پر شبہ تھا اس لئے میں نے وہ سیکشن ہی توڑ دیا! اور اس کے پورے عملے کو برطرف کردیا!"۔

"چوکیدار کا کیا ہوا تھا؟" جولیا نے پوچھا۔

"اوہ۔ اس کا آج تک پتہ نہیں لگا سکا! وہ مل جاتا تو اتنی درسری ہی کیوں مول لی جاتی۔ اس سے تو سب کچھ معلوم ہوسکتا تھا! اب آپ میری مدد کیجیئے!"۔

"مگر میں اس سلسلے میں کیا کرسکتی ہوں؟"

سرسوکھے کی ٹھنڈی سانس کمرے میں گونجی اور وہ تھوڑی دیر بعد مسکرا کر بولا! "اب مجھے پوری بات شروع سے بتانی پڑے گی۔۔ بات دراصل یہ ہے مس واٹر۔ میرے یہاں ایک اینگلو برمیز ٹائپسٹ تھی مس روشی۔ وہ آج کل رنگون گئی ہوئی ہے۔ اس نے ایک بار کسی مسٹر عمران کا تذکرہ کیا تھا جو پرائیویٹ سراغرساں ہیں!۔۔ اتفاق سے ایک دن مجھے اس نےدور سے مسٹر عمران کی زیارت بھی کرائی تھی اور مجھے اچھی طرح یاد ہے کہ آپ ان کے ساتھ تھیں"۔

"میں۔۔؟"

"جی ہاں۔ آپ۔۔ دیکھیئے مجھے شکلیں ہمیشہ یادر رہتی ہیں یہ اور بات ہے کبھی کبھی نام بھول جاتا ہوں مگر یہ بھی کم ہی ہوتا ہے! اس دوران میں جب یہ واقع پیش آیا مجھے مسٹر عمران کا خیال آیا تھا! مگر افسوس کہ مجھے ان کا پتہ نہیں معلوم تھا! اچانک ایک دن آپ نظر آگئیں! آپ اس وقت آفس میں داخل ہو رہی تھیں! میں نہیں جانتا تھا کہ آپ وہیں کام کرتی ہیں! میں نے پوچھ گچھ کی تو معلوم ہوا کہ آپ وہیں کام کرتی ہیں! میں سوچا واہ سرسوکھے تو بہت خوش نصیب ہو۔ تمہارا فارورڈنگ اور کلیرنگ کا کام بھی ہوتا رہے گا اور عمران ساحب تک پہنچ بھی ہوجائے گی۔۔ واہ۔۔ اور آج کل میرے ستارے بھی اچھے ہیں مس فٹز واٹر۔۔ اگر میں آپ کو صرف واٹر کہوں تو آپ کو کوئی اعتراض تو نہ ہوگا۔ فٹز واٹر کہنے میں زبان لڑکھڑاتی ہے"۔

"آپ مجھے صرف جولیانا کہہ سکتے ہیں!" جولیا بڑے دلآویز انداز میں مسکرائی۔

"اوہ۔ بہت بہت شکریہ!"۔ وہ خوش ہو کر بولا۔ "میں آپ کا بےحد ممنون ہوں اس وقت میرے دل پر سے ایک بہت بڑا بوجھ ہٹ گیا ہے! صرف آپ ہی سے میں یہ بات کہہ سکا ہوں!۔۔ اوہ مس فٹز واٹر میں کتنا خوش نصیب ہوں دراصل اسی گفتگو کے لئے میں نے آپ کو تکلیف دی تھی اور نہ حسابات تو سب جگہ کے یکساں ہوتے ہیں"۔

"پھر آپ کیا چاہتے ہیں۔۔؟"

"مجھے عمران صاحب سے ملائیے! ان سے سفارش کیجیئے۔ انہیں مجبورکیجیئے کہ اس معاملہ کا پتہ لگائیں۔ حالانکہ میں نے فارورڈنگ اینڈ کلیرنگ کے عملے کو الگ کردیا ہے مگر کون جانے اصل چور اب بھی یہیں موجود ہو اور کبھی اس کی ذات سے مجھے کوئی بڑا نقصان پہنچ جائے۔ میں نجی طورپر اس کی تحقیقات چاہتا ہوں۔ پولیس کو کانوں کان خبر نہ ہونی چاہیئے"۔

"دیکھیئے میں کوشش کروں گی! ویسے بہت دنوں سے عمران سے ملاقات نہیں ہوئی"۔

"کوشش نہیں! بلکہ یہ کام ضرور کیجیئے گا مس جولیانا۔۔ اخراجات کی پروا مجھے نہ ہوگی"۔

"آج آپ مقبرے کے نیچے مچھلیوں کا شکار کھیل رہے تھے؟" جولیا مسکرا کر بولی اور آپ کا اسپینیئل شکار کی ہوئی مچھلیاں گھسیٹ رہا تھا"۔

"شکار تو میں یقیناً کھیل رہا تھا!"۔ اس نے حیرت سے کہا۔ "مگر آپ کو یہ کیسے معلوم ہوا کہ مقبرے کے نیچے کھیل رہا تھا"۔

"میں نے آپ کو دیکھا تھا۔۔"!

"کمال ہے! آپ وہاں کہاں۔۔؟"

"میں بھی اوپر جھاڑیوں میں تیتر تلاش کر رہی تھی! کچھ فائر بھی کئے تھے! کیا آپ نے میرے فائروں کی آوازیں نہیں سنی تھیں؟"

"قطعی نہیں یا پھر ہوسکتا ہے میں نے دھیان نہ دیا ہو۔ اور تو کیا آپ بندوق چلاتی ہیں۔۔؟"

"مجھے بندوق سے عشق ہے"۔

"شاندار!۔۔" سرسوکھے بچکانہ انداز میں چیخا۔ اس کی آنکھوں کی چمک میں بھی بچپن ہی جھلک رہا تھا!۔۔ "آپ بندوق چلاتی ہیں! شاندار۔۔ آپ واقعئی خوب ہیں۔ مگر آپ نے مجھے آواز کیوں نہیں دی تھی!۔۔ آہا کبھی میرے ساتھ شکار پر چلیئے"۔

"فرصت کہاں ملتی ہے مجھے۔۔!" جولیا مسکرائی۔

"اوہ۔۔ تو آپ کو بہت کام کرنا پڑتا ہے"۔۔!

"بہت زیادہ۔۔!"

"بدتمیزی ضرور ہے مگر کیا پوچھ سکتا ہوں کہ آپ کو تنخواہ کتنی ملتی ہے؟"

"مجھے فی الحال وہاں ساڑھے چار سو مل رہے ہیں"۔

"بس۔۔ یہ تو کچھ بھی نہیں ہے! آپ پر اتنی ذمہ داریاں ہیں! اور تنخواہ! آپ جانتی ہیں روشنی کو یہاں کتنا ملتا تھا؟"

جولیا نے نفی میں سر ہلادیا!

"چھ سو!"

"اوہ۔۔!" جولیا نے خواہ مخواہ حیرت ظاہر کی۔ وہ سرسوکھے کو بد دل نہیں کرنا چاہتی تھی کیونکہ "چھ سو" کہتے وقت اس کا لہجہ فخریہ تھا!

"اور آپ کی خدمات کا معاوضہ تو ایک ہزار سے کسی طرح بھی کم نہ ہونا چاہیئے؟"

جولیا صرف مسکرا کر رہ گئی۔ انداز خاکسارانہ تھا!

"میں اسے بیہودگی تصور کرتا ہوں کہ آپ کو آفر دوں!۔۔ بہرحال جب بھی آپ وہاں سے بددل ہوں۔ سوکھے انٹرپرائزس کے دروازے اپنے لئے کھلے پائیں گی"۔

"بہت بہت شکریہ جناب!"

دفعتاً سر سوکھے نے انگلی اٹھا کر اسے خاموش رہنے کاا شارہ کیا اور اس کے چہرے پر ایسے آثار نظر آئے جیسے کسی کی آہٹ سن رہا ہو! جولیا بھی ساکت ہوگئی اس نے بھی کسی قسم کی آواز سنی تھی!

اچانک سرسوکھے خوف زدہ انداز میں دہاڑا۔ "کون ہے؟"

کسی کمرے میں کوئی وزنی چیز گری اور بھاگتے ہوئے قدموں کی آواز آئی ایسا لگا جیسے کوئی دوڑتا ہوا زینے طے کر رہا ہو۔۔!

سرسوکھنے جیب سے پستول نکال لیا! لیکن جولیا اس کے چہرے پر خوف کے آثار دیکھ رہی تھی!

"ٹھہریئے"۔ جولیا اٹھتے ہوئی بولی۔ "میں دیکھتی ہوں"۔

"اوہ۔۔ نہیں! پتہ نہیں کون تھا؟ بہرحال آپ نے دیکھ لیا تھا!" اس نے کہا اور دروازے کی طرف بڑھا۔ جولیا بھی اس کے پیچھے بڑھی! 
انہوں نے سارے کمرے دیکھ ڈالے۔ برابر والے کمرے میں دیوار کے قریب ایک چھوٹی سی میز گری ہوئی نظر آئی!

"یہ دیکھیئے۔۔" سرسوکھے نے کہا۔ "کوئی اس میز پر کھڑا ہو کر روشندان سے ہماری گفتگو سن رہا تھا!"

جولیا نے میز کی سطح پر ربر سول جوتے کے نشانات دیکھے۔

"آپ اس میز کو کسی کمرے میں مقفل کرادیجیئے! یہ نشانات عمران کے لئے کارآمد ہوسکتے ہیں"، جولیا نے کہا۔

"گڈ۔۔!" وہ خوش ہو کر بولا! "اب دیکھیئے یہ آپ کی ذہانت ہی تو ہے! مجھے اس کا خیال نہیں آیا تھا۔ اوہ مس جولیانا مجھے یقین ہے کہ اب میری پریشانیوں کا دور ختم ہوجائے گا۔ مجھے یقین ہے"۔

"آپ بالکل فکر نہ کریں۔۔" جولیا نے کچھ سوچتے ہوئے کہا۔ "آپ کے پاس بلڈ ہاؤنڈز بھی ہیں"!

"نہیں۔۔ کیوں۔۔!"

"اگر کوئی ہوتا تو اسے اس آدمی کی راہ پر بہ آسانی لگایا جاسکتا تھا جو اس وقت ہماری گفتگو سن رہا تھا!"

سرسوکھے کی آنکھیں حیرت سے پھٹی رہ گئیں!

"اوہ۔۔ مس جولیانا! آپ کی ذہانت کی کہاں تک تعریف کی جائے۔۔۔ آپ تو بہت گریٹ ہیں! عمران صاحب کی صحبت نے آپ کو بھی اچھا خاصہ جاسوس بنادیا ہے۔ کاش آپ ہمارے ساتھ ہوتیں! میں چین کی نیند لے سکتا! ساری تشویش ختم ہوجاتی۔۔!"

سرسوکھے نے خاموش ہو کر ٹھنڈی سانس لی۔

\endR

\bigskip\bigskip\bigskip
\centerline{\hdr{(۵)}}
\bigskip\bigskip\bigskip
\beginR
\body

اندھیری رات تھی۔ سڑک پر ویرانیاں رقص کر رہی تھیں! اور ان کا رقص دراصل جوزف کے وزنی جوتوں کی تال پر ہو رہا تھا! وہ اونٹ کی طرح سر اٹھائے چلا جا رہا تھا۔ گو اس وقت وہ فوجی لباس میں نہیں تھا! اور اس کے دونوں ریوالور بھی ہولسٹروں کی بجائے جیب میں تھے۔
اس سڑک پر الیکٹرک پول اتنے فاصلے پر تھے کہ دو روشنیوں کے درمیان میں ایک جگہ ایسی ضرورملتی تھی جہاں اندھیرا ہی رہتا تھا۔ درمیان میں دو پول چھوڑ کر بلب لگائے گئے تھے۔ یہ شہر سے باہر کا حصہ تھا۔ اگر ان اطراف میں دو چار فیکٹریاں نہ ہوتیں تو یہ سڑک بالکل ہی تاریک ہوتی۔

جوزف اس وقت کتھئی سوٹ اور سفید قمیض میں تھا! ٹائی تو وہ کبھی استعمال ہی نہیں کرتا تھا! آج کل وہ بالکل ہی دیو معلوم ہوتا تھا! عمران کی ڈنڈبیٹھکوں نے اس کا جسم اور زیادہ نمایاں کردیا تھا!

وہ یکساں رفتار سے چلتا رہا اور اس کے وزنی جوتوں کی آوازیں دور دور تک گونجتی رہیں۔ فیکٹریوں کے قریب پہنچ کر وہ بائیں جانب مڑ گیا!۔۔ یہ فیکٹریوں کی مخالف سمت تھی! ادھر دور تک ویرانہ ہی تھا۔ ناہموار اور جھاڑیوں سے ذھکی ہوئی زمین میلوں تک پھیلی ہوئی تھی۔

اچانک جوزف رک گیا! وہ اندھیرے میں آنکھیں پھاڑ رہا تھا۔۔ تقریباً سو گز کے فاصلے پر مشرق کی طرف اسے کوئی ننھی سی چمکدار چیز دکھائی دی اور وہ دوسرے ہی لمحے زمین پر تھا! اب وہ گھنٹوں اور ہتھیلیوں کے بل بالکل اسی طرح آہستہ آہستہ چل رہا تھا جیسے کوئی تیندوا شکار کی گھات میں ہو!

رخ اسی جانب تھا جہاں وہ ننھی سی چمکدار چیز نظر آئی تھی۔

"جوزف۔۔!" اس نے ہلکی سی سرگوشی سنی!۔۔ اور وہ کسی وفادار کتے کی طرح اچھل کر ادھر ہی پہنچ گیا!

"شش۔۔"!

جوزف جھاڑیوں میں دبک گیا پھر کوئی اس کے کاندھے پر ہاتھ رکھ کر بولا۔ " چند منٹ یہیں رکو"۔

جوزف جس پوزیشن میں تھا اسی میں رہ گیا! یہ اس کی عجیب وغریب عادت تھی۔ جب بھی اسے مخاطب کیا جاتا تو ہو اسی طرح ساکت ہوجاتاکہ اٹھا ہوا ہاتھ اٹھا ہی رہ جاتا! جماہی آ رہی ہوتی تو منہ پھیلا کر ہی رہ جاتا اور تاوقتیکہ کوئی نہ کہہ دی جاتی پھیلا ہی رہتا۔۔!

تھوڑی دیر بعد کہا گیا۔

"جوزف کیا تم اس وقت بہت خوش ہو"؟

"ہاں۔ باس بہت زیادہ۔۔ کیونکہ میں آج ایک نئی چیز دریافت کی ہے"۔

"اچھا۔۔!"

"ہاں باس! اگراسپرٹ اور پانی میں تھوڑا سا جنجر ایسنس بھی ملا لیا جائے تو بس۔۔ مزہ ہی آجاتا ہے"۔

"تم نے پھر اسپرٹ شروع کردی ہے؟"

"ہاں۔۔ باس۔۔"۔

"ایک ہزار ڈنڈ۔۔!"

"نن۔۔ نہیں۔۔ باس!" جوزف بوکھلا کر بولا! "نشہ اتر جائے گا! کھوپڑی بالکل خالی ہوجائے گی! اور میں کیچوا بن کر رہ جاؤں گا۔۔!"

"چلو اٹھو!۔۔" عمران نے اسےٹھوکا دیا۔

"ہم کہاں چلیں گے باس۔۔؟"

"کالا گھاٹ۔۔ تم نے دیکھا نا؟"

"ہاں۔۔ باس۔۔"۔

"وہاں ایک شراب خانہ ہے!"

"میں جانتا ہوں باس!"۔۔ جوزف خوش ہو کر بولا! "وہاں تاڑی بھی ملتی ہے"۔

"ہوم۔۔! اس شراب خانہ کے پاس ندی کی سمت جو ڈھلان شروع ہوتی ہے!تمہیں وہاں رکنا ہوگا!"

"ڈھلان پر رک کر کیا کروں گا باس! کہ آپ شراب خانہ میں جائیں اور میں ڈھلان پر کھڑا رہوں"۔

"چلتے رہو۔۔!"

وہ اندھیرے ہی میں ناہموار راستے طےکرتے رہے! کبھی کبھی محدود روشنی والی چھوٹی سی ٹارچ روشن کرلی جاتی!

جوزف کچھ بڑبڑا رہا تھا!

"خاموشی سے چلتے رہو"۔ کہا گیا۔

آدھے گھنٹے بعد وہ ایک ڈھلوان راستے پر چل رہے تھے جہاں سے ندی کے کنارے والے چراغوں کے سلسلے صاف نظر آنے لگے تھے!

"ایک بار پھر سنو جوزف!" اس سے کہا گیا۔ "تم شراب خانے کی پشت پر ندی والی ڈھلان پر ٹھہرو گے"۔

"اچھا باس!" جوزف نے بےحد اداس لہجے میں کہا۔

"مگر تم وہاں کیوں ٹھہرو گے؟"

"جماہیاں لینے اور آنسو بہانے کے لئے!" جوزف کی آواز دردناک تھی!

عمران ہنس پڑا۔

"مگر باس ! تم اپنے محل میں کیوں نہیں آتے۔۔؟" جوزف نے کہا!

"یہ ایک درد بھری کہانی ہے۔۔ جوزف!" عمران غمناک لہجے میں بولا۔ " میری آخری بیوی کے رشتے دار مجھے قتل کردینا چاہتے ہیں۔۔!"

"اف۔فوہ! "جوزف چلتے چلتے رک گیا۔ اسے وہ پھرتیلا بوڑھا یاد آگیا تھا جس نے دو تین دن پہلے رانا پیلس میں اپنی چلت پھرت کا مظاہرہ کیا تھا!

بلیک زیرو کو علم ہی نہیں تھا کہ عمران کہاں ہوگا اس لئے یہ کہانی عمران تک نہیں پہنچ سکی تھی! اتفاق سے آج صبح جوزف ہوا خوری کو نکلا تھا۔ راستے میں ایک لڑکے نے اسے ایک خط دیا جو عمران کی طرف سے ٹائپ کیا گیا تھا اور جس میں جوزف کے لئے ہدایت تھی کہ وہ رات کو فلاں وقت فلاں مقام پر پہنچ جائے۔

جوزف اس معاملہ میں اتنا محتاط ثابت ہوگا کہ اس نے اس کا تذکرہ بلیک زیرو (طاہر صاحب) سے بھی نہیں کیاتھا۔ حالانکہ وہ خود بھی دھوکا کھا سکتا تھا۔ کیونکہ وہ خط ٹائپ کیا ہوا تھا اور اس کے نیچے بھی عمران کے دستخط نہیں تھے بلکہ نام ہی ٹائپ کردیا گیا تھا! لیکن اس نے کسی وفادار کتے کی طرح اس میں عمران کی بو محسوس کی تھی اور نتیجے کے طور پر وہ اس وقت یہاں موجود تھا!

"کیوں رک گئے؟" عمران نے ٹوکا۔

اس پر اس نے جڑی بوٹیاں فروخت کرنے والے بوڑھے کی داستان دہرائی اور بتایا کہ کس طرح اس نے اس کی جیب سے پستول نکال لیا تھا۔

عمران سوچ میں پڑ گیا اس کی سمجھ میں نہیں آ رہا تھا کہ اس پر عمران ہونے کی بنا پر حملے ہو رہے تھے یا اس لئے کوئی اس کے پیچھے پڑ گیا تھا کہ رانا تہور علی صندیقی کا راز معلوم کرسکے۔ یا پھر حملہ آواروں کی نظروں میں بھی تہور علی اور عمران ایک ہی شخصیت کے دو مختلف روپ تھے۔۔!

"بس اسی سے اندازہ کرلو۔ جوزف۔۔ کہ آج کل میں کتنی الجھنوں میں گھرا ہوا ہوں۔۔!"

"مجھے ان کا پتہ بتاؤ باس! ایک کو بھی زندہ نہ چھوڑوں گا"۔ جوزف بھرائی ہوئی آواز میں بولا!

"چلتے رہو۔۔!" عمران بولا۔ وہ سوچ رہا تھا کہ اب وہ اپنے ماتحتوں کو اپنے قریب بھی نہیں آنے دے گا ورنہ اس کا امکان بھی ہے کہ اسی سلسلے میں ڈھمپ اینڈ کو کا راز ہی فاش ہوجائے۔

"ہاں تو باس! مجھے ڈھلان پر کیا کرنا ہوگا؟"

"اگر میری عدم موجودگی وہاں کوئی سبز رنگ کا موٹر بوٹ آئے تو تم فوراً ہی ایک ہوائی فائر کردینا"۔

"بس صرف ہوائی فائر کردوں گا!" جوزف نے پھر مایوسانہ انداز میں پوچھا۔

"تم پر خون کیوں سوار رہتا ہے جوزف؟"

"نہیں تو باس!۔۔ وہ دراصل میں سوچتا ہوں کہ مجھے پھانسی کیوں نہ ہوجائے میں نے سنا ہے کہ اب اسپرٹ میں لائسنس کے بغیر نہیں ملا کرے گی۔ مجھے کون لائسنس دے گا! اس لئے بہتر یہی ہے کہ میں کسی کو قتل کرکے جیل چلا جاؤں!"

"اور اگر میں ہی تمہیں قتل کردوں تو۔۔!"

"نہیں! اس کی بجائے میری بوتلوں میں اضافہ کردو۔ باس!" جوزف گھگھیا۔

"اب روزانہ پانچ ہزار ڈنڈ۔۔!"

"مم۔۔ مرا۔۔ نہیں۔۔ نہیں باس میرےپھیپھڑے پھٹ جائیں گے"۔

"خاموش رہو۔ ہم شراب خانے کے قریب ہیں! تم یہیں سے اسی پگڈنڈی پر مڑ جاؤ! آگے چل کر یہ دو مختلف سمتوں میں تقسیم ہوگئی ہے مگر تم بائیں جانب مڑ جانا۔ پگڈنڈی نہ چھوٹنے پائے۔ اس طرح تم ٹھیک اسی جگہ پہنچو گے جہاں ٹھہر کر تمہیں میرا انتظار کرنا ہے"۔

"اچھا باس!"جوزف کسی بہت ہی ستم رسیدہ آدمی کی طرح ٹھنڈی سانس لے کر پگڈنڈی پر مڑ گیا۔۔!

عمران جو اب روشنی میں آچکا تھا یعنی طور پر جوزف کے لئے ایک مسئلہ بن کر رہ جاتا!۔۔ اسی لئے اور بھی اس نے اسے اندھیرے ہیں میں رخصت کردیا تھا! وہ دراصل ایک بوڑھے بھکاری کے روپ میں تھا اور اس کے جسم پر چیتھڑے جھول رہے تھے!

جوزف چلتا رہا! اس مقام کو پہچاننے میں بھی اسے کوئی دشواری نہیں پیش آئی۔ جہاں پگڈنڈی دو شاخوں میں بٹ کر مخالف سمتوں میں مڑ گئی تھی! وہ عمران کی بتائی ہوئی سمت پر چلنے لگا!۔۔

ہوٹل کی پشت پر پہنچ کر اس نے چاروں طرف نظریں دوڑائیں! گہرا اندھیرا فضا پر مسلط تھا! کہیں کہیں روشنی کے نقطے سے نظر آرہے تھے!

جوزف لاکھ ڈفر سہی لیکن خطرات کے معاملہ میں وہ جانوروں کی سی حس رکھتا تھا! اس نے سوچا کہ فائر کرنے کے بعد وہ کیا کرے گا! اگر کچھ لوگ آگئے اور وہ پکڑ لیا گیا تو۔۔! کیا باس اسے پسند کرے گا!۔۔

اب وہ کوئی ایسا درخت تلاش کرنے لگا جسے فائر کرنے کے بعد اپنے بچاؤ کے لئے استعمال کرسکے!

اچانک ایک موٹر بوٹ گھاٹ سے آلگی۔۔

جوزف نے تیزی سے جیب میں ہاتھ ڈالا لیکن پھر آنکھیں پھاڑ کر رہ گیا! بھلا اندھیرے میں موٹر بوٹ کا رنگ کیسے نظر آتا! ہیڈ لیمپ کی روشنی بھی اسے نہ ظاہر کرسکتی تھی!۔۔

"او۔۔ باس!" جوزف دانت پیس کر بڑبڑایا۔ "تم نشے میں تھے یا مجھے ہی ہوش نہیں تھا! سبز رنگ۔۔ ہائے سبز رنگ۔۔ زرد نکلے تو کیا ہوگا۔۔ نیلا۔۔ اودا۔۔ کتھئی۔۔ زعفرانی۔۔ اب میں کیا کروں۔۔؟ او باس!۔۔

وہ کھڑا دانت پیستا رہا پھر اپنے سر پر مکے مارنے لگا!

بہرحال اب اس کے لئے ضروری ہوگیا تھا کہ وہ عمران کو تلاش کرکے پوچھتا کہ اندھیرے میں موٹر بوٹ کا رنگ کیسے دیکھا جائے؟

وہ شراب خانے کے صدر دروازے کی طرف چل پڑا۔ اسے یقین تھا کہ عمران شراب خانے ہی میں ملے گا!۔۔ شاید اس نے کہا بھی تھا!۔۔

شراب خانہ پوری طرح آباد ملا! اس کی چھت زیادہ اونچی نہیں تھی! دیواریں اور چھت سفید آئل پینٹ سے رنگی گئی تھیں! بس ایسا ہی معلوم ہوتا تھا جیسے وہ کسی بہت بڑے بحری جہاز کا شراب خانہ ہو! لیکن یہاں اتنی صفائی اور خوش سلیقگی کو دخل نہیں تھا۔

لوگ میلی کچیلی میزوں پر بیٹھے تاڑی یا دیسی شراب پی رہے تھے! ویسے بھی یہاں قیمتی شرابیں شاذونادر ہی ملتی تھیں!

یہاں پہنچ کر جوزف کی پیاس بری طرح جاگ اٹھی۔ وہ ہونٹوں پر زبان پھیرتا اور چندھیائی ہوئی آنکھوں سے چاروں طرف دیکھتا رہا! لیکن یہاں کہیں اسے عمران نہ دکھائی دیا!

وہ جو ابھی زیادہ نشے میں نہیں تھے اسے گھورنے لگے تھے!

دفعتاً ایک بوڑھا آدمی جھومتا ہوا اپنی میز سے اٹھا اور جوزف کی طرف بڑھنے لگا! اس کے ہاتھ میں گلاس تھا!

اس کی ہیت کذائی پر جوزف کو ہنسی آگئی۔ یہ ایک پست دبلا پتلا آدمی تھا! چہرے پر اگر ڈاڑھی نہ ہوتی تو بالکل گلہری معلوم ہوتا! آنکھیں دھندلی تھیں۔

جوزف کے قریب پہنچ کر وہ رک گیا اور اس طرح سر اٹھا کر اس کی شکل دیکھنے لگا جیسے کسی منارہ کی چوٹی کا جائزہ لے رہا ہو!۔۔

"کیا ہے۔۔؟" جوزف کھسیانے انداز میں ہنس کر پوچھا۔

"مجھے ڈر ہے کہ کہیں تمہارے کانوں تک اپنی آواز پہنچانے کے لئے مجھے۔۔لاؤڈ اسپیکر نہ استعمال کرنا پڑے!"

"ہام!" جوزف اسے پکڑنے کے لئے جھکا اور وہ اچھل کر پیچھے ہٹ گیا!

"خفا ہونے کی ضرورت نہیں ہے۔ میں بہت غم زدہ آدمی ہوں"۔ بوڑھے نے رونی آواز میں کہا۔ وہ انگریزی ہی میں گفتگو کر رہا تھا!
"کیا ہوا ہے تمہیں"۔ جوزف غرایا۔

"ادھر چلو۔ میں تمہیں پلاؤں گا! تمہیں اپنی دکھ بھری داستان سناؤں گا! مجھے یقین ہے کہ تم میری مدد کرو گے! بہت زیادہ لمبے آدمی عموماً مجھ پر رحم کرتے ہیں"۔

"میں نہیں پیوں گا۔۔!" جوزف نے احمقانہ انداز میں کہا اور پھر چاروں طرف دیکھنے لگا۔

"کیا تمہیں کسی کی تلاش ہے"۔ بوڑھے نے پوچھا۔

"نہیں۔!"

"تو پھر آؤ۔ نا۔۔ غم غلط کریں۔ تم مجھے کوئی بہت شریف آدمی معلوم ہوتے ہو"۔

"ہاں۔!" جوزف نے سر ہلا کر پلکیں جھپکائیں۔

"آؤ۔۔ دوست آؤ۔ تمہار دل بہت نورانی ہے!"

جوزف سچ مچ خوش ہوگیا! اپنی صفائے دل کے متعلق کسی سےکچھ سن کر وہ نہال ہوجاتا تھا۔ ایسے مواقع پر اسے فادر جوشوا یاد آجتے جنہوں نے اسے عیسائی بنایا تھا اور جو اکثر کہا کرتے تھے کہ "تم سفید فاموں سے افضل ہو کیونکہ تم کالوں کے دل بڑے نورانی ہوتے ہیں"۔

بوڑھا اسے اپنی میز پر لے آیا۔

"اوہ۔۔ شکریہ! میں گھر سے باہر کبھی کچھ نہیں پیتا!" جوزف نے کہا۔

"یہ بہت بری عادت ہے دوست! گھر پر پینے سے کیا فائدہ۔ کیا دیواروں سے دل بہلاتے ہو!"

"عادت ہے۔۔! جوزف نے خواہ مخواہ دانت نکال دیئے۔

"نہیں میری خاطر! پیو! میں بہت غم زدہ آدمی ہوں۔۔ میری بات نہ ٹالو! ورنہ میرے غموں میں ایک کا اور اضافہ ہوجائے گا!"
"تمہیں کیا غم ہے؟"

"ایک دو۔ نہیں۔۔ ہزاروں میں!۔۔ بس تم پیو پیارے۔۔ یہی میرے غم کا علاج ہے۔ تم بہت نیک آدمی ہو ضرور پیو گے مجھے یقین ہے۔۔!"

"کیا میرے پینے سے تمہارے غم دور ہوجائیں گے!" جوزف نے بڑی معصومیت سے پوچھا!

"قطعی دور ہوجائیں گے۔۔!"

"اچھا تو پھر میں پیوں گا! خدا تمہاری مشکل آسان کرے!" جوزف نے انگلیوں سے کراس بنایا۔

"کیا پیو گے؟"

"تاڑی۔۔ سالہاسال ٍذرے کہ میں تاڑی نہیں پی۔۔!"

"مذاق مت کرو پیارے۔!" بوڑھے نے کہا۔

"میں مذاق نہیں کر رہا"۔ جوزف کو غصہ آگیا!

"اچھا۔۔ اچھا۔۔ تاڑی ہی سہی"۔ بوڑھے نے کہا اور اٹھ کر کاؤنٹر کی طرف چلا گیا۔ واپسی پر اس کے ہاتھوں میں تاڑی کی بوتل اور گلاس تھے۔

جوزف نے حلق تر کرنا شروع کیا! جب کھوپڑی کچھ گرم ہوئی تو میز پر گھونسہ مار کر بولا۔ "بتاؤ کس کی وجہ سے تمہیں اتنے دکھ پہنچے ہیں؟"

"ابھی بتاؤں گا۔۔ سن سے پہلے آج کا غم دہراؤں گا!"

تھوڑی دیر تک خاموشی رہی پھر بوڑھے نے کہا۔ "هزاروں روپے کی شراب برباد ہوجائے گی۔ اگر میں نے دو گھنٹے کے اندر ہی اندر کوئی قدم نہ اٹھایا!"

"شراب برباد ہوجائے گی!" جوزف نے متحیرانہ انداز میں پلکیں جھپکائیں۔

"ہاں! پانچ بیرل۔ یہاں سے تقریبا ایک میل کے فاصلے پر جنگل میں پڑے ہوئے ہیں۔ میں نے ہی انہیں وہاں چھپایا تھا۔ اب اطلاع ملی ہے کہ پولیس کو شبہ ہوگیا ہے! اس لئے وہ عنقریب وہاں گھیرا ڈالنے والی ہے۔ کاش میرے بازؤوں میں اتنی قوت ہوتی کہ میں ان بیرلوں کو قریب ہی کے ایک کھڈ میں لڑھکا سکتا!"

"یہ کون سی بڑی بات ہے"۔ جوزف اکڑ کر بولا۔ "میں چل کر لڑھکا دوں گا!"

"اوہ۔۔ اگر تم ایسا کرسکو تو ایک بیرل تمہارا انعام۔۔!"

"لاؤ۔۔ ہاتھ"۔ جوزف میز پر ہاتھ مار کر بولا! "بات پکی ہوگئی! میں لڑھکاؤں گا اور تم اس کے عوض مجھے ایک بیرل دو گے!"

پھر تاڑی کی مزید دو بوتلیں ختم ہونے تک بات بالکل ہی پکی ہوگئی اور جوزف لڑکھڑاتا ہوا اٹھا۔۔ بوڑھا آدمی کسی ننھے سے بچے کی طرح اس کی انگلی پکڑے چل رہا تھا!۔۔

یہ جوڑا دیکھ کر لوگ بےتحاشہ ہنسے تھے۔۔ اور جوزف تو اب اسے قطعی فراموش کرچکا تھا کہ یہاں کیوں آیا تھا!۔۔

\endR

\bigskip\bigskip\bigskip
\centerline{\hdr{(۶)}}
\bigskip\bigskip\bigskip
\beginR
\body


ایکس ٹو نے اپنے ماتحتوں کو باقاعدہ طور پر ہدایت کردی تھی کہ وہ عمران کے متعلق کسی چکر میں نہ پڑیں۔ نہ تو اس کے فلیٹ کے فون نمبر رنگ کئے جائیں اور نہ کوئی ادھر جائے! جولیا کو اس قسم کی ہدایت دیتے وقت اس کا لہجہ بےحد سخت تھا!

جولیا اس پر بری طرح جھلا گئی تھی! لیکن کرتی بھی کیا! ایکس ٹو بہرحال اپنے ماتحتوں کے اعصاب پر سوار تھا! وہ اس سے اسی طرح خائف رہتے تھے جیسے ضعیف الاعتقاد لوگ ارواح کے نام پر لرزہ براندام ہوجاتے ہیں!

مگر جولیا الجھن میں مبتلا تھی۔ آج کل ایک ناقابل فہم سی خلش ہر وقت ذہن میں موجود رہتی اور اس کا دل چاہتا کہ وہ شہر کی گلیوں میں بھٹکتی پھرے! چھتوں اور دیواروں کے درمیان گھٹن سی محسوس ہوتی تھی!

آج صبح اس نے فون پر بڑے جھلائے ہوئے انداز میں ایکس ٹو سے گفتگو کی تھی۔ اسے بتایا تھا کہ سرسوکھے کی بھاگ دوڑ کا اصل مقصد کیا ہے! پھروہ اس کے لئے عمران کو تلاش کرے یا نہ کرے!۔۔

"بس اسی حد تک جولیا ناکہ وہ مطمئن ہوجائے!" ایکس ٹو نے جواب دیا تھا! " اسے یہ شبہ نہ ہونا چاہیئے کہ تم اسے ٹال رہی ہو! بلکہ عمران کی گمشدگی پر پریشانی بھی ظاہر کرو!"

جولیا برا سا منہ بنا کر رہ گئی تھی!

سرسوکھے کی فرمائش کے مطابق آج اسے عمران کی تلاش میں اس کا ساتھ دینا تھا! سب سے پہلے وہ عمران کے فلیٹ میں پہنچے لیکن سلیمان سے یہی معلوم ہوا کہ عمران پچھلے پندرہ دنوں سے غائب ہے! پھر جولیا نے ٹپ ٹاپ نائٹ کلب کا تذکرہ کرتے ہوئے کہا کہ عمران وہاں کا مستقل ممبر ہے۔ ہوسکتا ہے کہ وہاں اس کےمتعلق کچھ معلومات حاصل ہوسکیں۔

وہ ٹپ ٹاپ کلب پہنچے۔ یہاں بھی کوئی امید افزا صورت نہ نکل سکی! آخر سرسوکھے نے تھکے ہوئے لہجے میں کہا۔ "اب کہاں جائیں۔ میں واقعی بڑا بدنصیب ہوں مس جولیانا۔ آئیے کچھ دیر یہیں بیٹھیں!"

جولیا کو اس پہاڑ نما آدمی سے بڑی الجھن ہوتی تھی! اس کے ساتھ کہیں نکلتے ہوئے اس کے ذہن میں صرف یہی ایک خیال ہوتا تھا کہ وہ بڑی مضحکہ خیز لگ رہی ہوگی۔ آس پاس کے سارے لوگ انہیں گھور رہے ہوں گے!

مگر اس کمبخت ایکس ٹو کو کیا کہیئے جس کا حکم موت کی طرح اٹل تھا!

وہ سرسوکھے کے ساتھ بیٹھی اور بور ہوتی رہی! لیکن پھر اس نے ریکرئیشن ہال میں چلنے کی تجویز پیش کی!

مقصد یہ تھا کہ وہاں کوئی نہ کوئی اس سے رقص کی درخواست ضرور کرے گا اور سرسوکھے سے پیچھا چھوٹ جائے گا! سرسوکھے اس تجویز پر خوش ہوا تھا!

وہ ریکریشن ہال میں آئے۔ یہاں ابھی آرکسٹرا جاز بجا رہا تھا! اور چند باوردی منتظمین چوبی فرش پر پاؤڈر چھڑکتے پھر رہے تھے۔
وہ گیلری میں جا بیٹھے! تھوڑی دیر بعد رقص کے لئے موسیقی شروع ہوئی!

"کیا میں آپ سے رقص کی درخواست کرسکتا ہوں!" سرسوکھے نے ہچکچاتے ہوئے کہا!

"آپ!" جولیا نے متحیرانہ لہجے میں سوال کیا! اس کا سر چکرا گیا تھا!

"اوہ"۔ دفعتاً سرسوکھے بےحد مغوم نظر آنے لگا! کرسی کی پشت سے ٹکتے ہوئے اس نے چھت پر نظریں جما دیں! جولیا کو اپنے رویے پر افسوس ہونے لگا کیونکہ سرسوکھے کی آنکھیوں میں آنسو تیر رہے تھے! جولیا نے محسوس کیا کہ اس کا وہ "آپ" گویا ایک تھپڑ تھا جو سرسوکھے کے دل پر پڑا تھا! کیونکہ "آپ" کہتے وقت جولیا کےلہجے میں تحیر سے زیادہ تضحیک تھی!

"اوہو۔۔ تو پھر۔۔ آپ اٹھیئے نا!" جولیا نے بوکھلائے ہوئے لہجے میں کہا۔

وہ ہنسنے لگا۔ بےتکی سی ہنسی! ایسا معلوم ہو رہا تھا جیسے خود اسے بھی احساس ہو کہ وہ یوں ہی احمقانہ انداز میں ہنس پڑا ہے۔ پھر وہ آنکھیں ملنے لگا!

"نہیں ۔!" وہ کچھ دیر بعد بھرائی ہوئی آواز میں بولا۔ "میں اپنی اس بےتکی درخواست پر شرمندہ ہوں! میں آپ کو بھی مضحکہ خیز نہیں بنانا چاہتا !"

وہ پھر ہنسا مگر جولیا کو اس کی ہنسی دردناک معلوم ہوئی تھی! ایسا لگا تھا جیسے متعدد کراہوں نے ہنسی کی شکل اختیار کرلی ہو!

"مس فٹز واٹر!" اس نے اپنے سینے پر ہاتھ مار کر کہا۔ "ہڈیوں اور گوشت کا یہ بنجر پہاڑ ہمیشہ تنہا کھڑا رہے گا۔ میں نے نہ جانے کس رو میں آپ سے درخواست کردی تھی! اداس اور تنہا آدمی بچوں کی سی ذہنیت رکھتے ہیں"۔ گوشت اور ہڈیوں کے اس بےہنگم سے ڈھیر میں چھپا ہوا سرسوکھے رام بچہ ہی تو ہے جو بڑی لاپروائی سے اس بدنما ڈھیر کو اٹھائے پھرتا ہے۔ اگر باشعور ہوتا تو۔۔"

"اور دیکھیئے! آپ بالکل غلط سمجھے سرسوکھے! میرا یہ مطلب ہرگز نہیں تھا! دراصل مجھے اس پر حیرت تھی کہ۔۔!"

"نہیں۔ مس جولیانا! میں خود بھی تماشہ بننا پسند نہیں کروں گا!" وہ ہاتھ اٹھا کر دردناک آواز میں بولا۔

جولیا خاموش ہوگئی! رقص شروع ہوکا تھا! سرسوکھے رقاصوں کو کسی بچے ہی کے سے انداز میں دیکھتا رہا۔۔! نہ جانے کیوں جولیا سچ مچ اس کے لئے مغوم ہوگئی تھی!

\endR

\bigskip\bigskip\bigskip
\centerline{\hdr{(۷)}}
\bigskip\bigskip\bigskip
\beginR
\body


جوزف بس چلتا ہی رہا! اسے احساس نہیں تھا کہ وہ کتنا چل چکا ہے۔ اور کب تک چلتا رہے گا۔ اس کے ساتھ ہی اس کی زبان ہی چل رہی تھی ۔ نوجوانی کے قصے چھیڑ رکھے تھے!

نوجوانی کے قصے بھی جوزف کی ایک کمزوری تھی۔ وہ مزے لے لے کر اپنے کارنامے بیان کرتا تھا اور ان کہانیوں کے درمیان قبیلے کی ان لڑکیوں کا تذکرہ ضرور آتا تھا جو اس پر مرتی تھیں۔ اس مرحلہ پر جوزف کے ہونٹ سکڑ جانے اور آواز میں سختی پیدا ہوجاتی۔ ایسا لگتا جیسے حقیقتاً اسے کبھی ان کی پرواہ نہ ہوئی ہو! اس وقت وہ بوڑھے سے کہہ رہا تھا۔ "بھلا بتاؤ۔ مجھے ان باتوں کی فرصت کہاں ملتی تھی۔ میں تو زیادہ تر رائفلوں اور نیزوں کے کھیل میں الجھا رہتا تھا۔ جب بھی سفید فام شکاری میرے علاقہ میں داخل ہوتے تو انہیں تندوے کی تلاش ضروری ہوتی تھی! میں ہی ان کی رہنمائی کرتا تھا۔ ان کی زندگیاں میری مٹھی میں ہوتی تھیں۔۔ اب بتاؤ۔۔ تم ہی بتاؤ۔۔ میں کیا کرتا! نگانہ جو قبیلے کی سب سے حسین لڑکی تھی! اس نے مجھے بددعائیں دی تھیں۔۔ آہ۔۔ آج میں اسی لئے بھٹکتا پھر رہا ہوں۔ مگر بتاؤ! اس کے لئے کہاں سے وقت نکالتا۔۔!"

جوزف نے پھر بکواس شروع کردی۔ تاڑی کی تین بوتلیں ہٹلر بھی بن سکتی ہیں اور علم الکلام کی ماہر بھی۔۔!

اچانک بوڑھا چلتے چلتے رک گیا۔ اور خوش ہو کر بولا! "واہ۔۔ اب تو وہ بیرل یہاں سے لے جائے بھی جاسکتے ہیں! میرے آدمی ٹرک لے آئے ہیں لیکن پولیس کا کہیں پتہ نہیں ہے۔۔!"

"ہائیں!" جوزف منہ پھاڑ کر رہ گیا۔ پھر بولا! "اب میرے انعام کا کیا ہوگا!"

"ایک بیرل تمہارا ہے دوست!" بوڑھے نے اس کی کمر تھپتھپا کر کہا! "تم اب انہیں ٹرک میں چڑھانے میں مدد دو گے"۔

ٹرک قریب ہی موجود تھا۔ اس کا پچھلا ڈھکنا زمین پر ٹکا ہوا تھا۔ جوزف نے چندھائی ہوئی آنکھوں سے چاروں طرف دیکھا! یہ ایک ویرانہ تھا۔ گھنیرے درخت اور جھاڑ جھنکار قریب وجوار کے اندھیرے میں کچھ اور اضافہ کرتے ہوئے سے معلوم ہو رہے تھے۔
"چلو۔ اندازہ کرلو کہ تم بیرل اوپر چڑھا سکو گے یا نہیں!" بوڑھے نے کہا اور ٹرک پر چڑھ گیا۔

جوزف کی رفتار سست تھی۔ لیکن وہ بھی اوپر پہنچ ہی گیا! ٹرک تین طرف سے بند تھا اور اس کی چھت کافی اونچی تھی! لیکن جوزف جیسے لمبے تڑنگے آدمی کو تو جھکنا ہی پڑا تھا۔

"چڑھا سکو گے نا؟" بوڑھے نے پوچھا۔

"بل۔۔ بل۔۔ بلکول۔۔" جوزف لڑکھڑایا اور آندھی سے اکھڑتے ہوئے کسی تناور درخت کی طرح ڈھیر ہوگیا! اسے اس پر بھی غور کرنے کا موقعہ نہیں مل سکا تھا کہ کھوپڑی پر ہونے والے تین بھرپور وار زیادہ نشہ آور ہوتے ہیں۔۔ یا تاڑی کی تین بوتلیں۔۔!

اس کا ذہن تاریکی کی دلدل میں ڈوبتا چلا گیا! پھر دونوں ٹرک کے اگلے حصے میں چلے گئے!

تھوڑی دیر بعد ٹرک چل پڑا!

\endR

\bigskip\bigskip\bigskip

\centerline{\hdr{(۸)}}
\bigskip\bigskip\bigskip
\beginR
\body


صفدر نے اس دن کے بعد سے اب تک ڈھمپ اینڈ کو کے دفتر کی شکل نہیں دیکھی تھی۔ جب وہاں عمران کی موت کی اطلاع لے کر گیا تھا! ایکس ٹو کی طرف سے اسے یہی ہدایت ملی تھی!

لیکن وہ عمران کے متعلق الجھن میں تھا! کبھی یقین کرنے پر مجبور ہوتا کہ اب عمران اس دنیا میں نہیں! اور کبھی پھر کئی طرح کے شبہات سر اٹھاتے! مگر یہ تو اس کی آنکھوں کے سامنے کی بات تھی کہ عمران چیخ مار کر ندی میں جا پڑا تھا! کچھ بھی ہو دل نہیں چاہتا تھا کہ عمران کی موت پر یقین کرے!

جولیا نے کسی کو بھی نہیں بتایا تھا کہ عمران زندہ ہے اور اسے اس واقعہ کے بعد اس کی کوئی تحریر ملی تھی! ایکس ٹو تو اسے یقینی طور پر صحیح حالات کا علم تھا۔ ورنہ وہ جولیا کو فون پر سرزنش کیوں کرتا۔ یہی سوچ کر جولیا نے اس سے بھی اس مسئلہ پر کسی قسم کی گفتگو نہیں کی تھی!

بہرحال صفدر آج کل زیادہ تر گھر ہی میں پڑا رہتا تھا۔۔ اس وقت بھی وہ آرام کرسی میں پڑا اونگھ رہا تھا! اچانک فون کی گھنٹی بجی جو ان دونوں شاذونادر ہی بجتی تھی!

وہ اچھل پڑا۔۔!

"ہیلو۔۔!" اس نے ماؤتھ پیس میں کہا۔

"ہائیں۔۔!" دوسری طرف سے آواز آئی۔ "کیا تم زندہ ہو؟"

"ارے!" صفدر پرمسرت لہجے میں چیخا! "آپ۔۔!"

اس نے عمران کی آواز صاف پہچان لی تھی۔

"اتنی زور سے نہ چیخو کہ تمہاری لائن کو شادیٴ مرگ ہوجائے۔ ویسے میں عالم بالا سے بو ل رہا ہوں!"

"عمران صاحب خدا کے لئے بتایئے کہ وہ سب کیا تھا؟"

"یار بس کیا بتاؤں"۔ دوسری طرف سے مغوم لہجے میں کہا گیا! " میں تو یہی سمجھ کر مرا تھا کہ گولی لگ چکی ہے۔ مگر فرشتوں نے پھر دھکا دے دیا! کہنے لگے کھسکو یہاں سے۔ یہاں چارسو بیسی نہیں چلے گی۔ گولی وولی نہیں لگی۔ آئندہ اچھی طرح مرے بغیر ادھر کا رخ بھی نہ کرنا۔ نہیں تو اب کی دم لگا کر واپس کئے جاؤ گے!"

صفدر ہنسنے لگا! وہ بےحد خوش تھا۔ اس کی بیک بہت بڑی الجھن رفع ہو گئی تھی!

"جولیا بے حد پریشان تھی۔۔!" صفدر نے کہا۔

"پچھلے سال میں اس سے ساڑھے پانچ روپے ادھار لیئے تھے نا۔۔ آج تک واپس نہیں کرسکا۔۔!"

"عمران صاحب خدا آپ کو جمالیاتی حس بھی عطا کردے تو کتنا اچھا ہو!"

"تب پھر لوگ مجھے جمال احمد کہیں!" عمران خوش ہو کر بولا۔ "اور میں جمالی تخلص کرنے لگوں! خیر اس پر کبھی سوچیں گے۔ اس وقت تمہیں ایک آدمی کا تعاقب کرنا ہے جو ٹپ ٹاپ نائٹ کلب کے بلیرڈ روم نمبر۳ میں بلیرڈ کھیل رہا ہے۔ اس کے جسم پر سرمئی آئیرن کا سوٹ ہے اور گلے میں نیلی دھاریوں والی زرد ٹائی۔ اگر وہ تمہارے پہنچنے تک وہاں سے جاچکا ہو تو پھر وہیں ٹھہرنا"۔

دوسری طرف سے سلسلہ منقطع ہوگیا!

صفدر کو ٹپ ٹاپ نائٹ کلب پہچنے میں بیس منٹ سے زیادہ نہیں لگے تھے! وہ آدمی اب بھی بلیرڈ روم میں موجود تھا جس کے متعلق عمران نے بیس منٹ پہلے اس سے فون پر گفتگو کی تھی۔ یہ ایک لمبا تڑنگا اور صحت مند جوان تھا۔ جبڑوں کی بناوٹ اس کی سخت دلی کا اعلان کر رہی تھی۔ البتہ آنکھیں کاہلوں اور شرابیوں کی سی تھیں۔ آنکھوں کی بناوٹ اور جسم کے پھرتیلے پن میں بڑا تضاد تھا۔

صفدر اس طرح ایک خالی کرسی پر جا بیٹھا جیسے وہ بھی کھیلنے کا ارادہ رکھتا ہو! یہاں چار بلیرڈ روم تھے اور ہر کمرے میں دو دو میزیں تھیں! اس کمرے کی دونوں میزوں پر کھیل ہو رہا تھا!

بھاری جبڑے والے کا ساتھی تھوڑی دیر بعد ہٹ گیا! اور بھاری جبڑے والے صفدر سے پوچھا!

"کیا آپ کھیلیں گے؟"

"جی ہاں۔۔!" صفدر اٹھ گیا۔

دونوں کھیلنے گے! کچھ دیر بعد صفدر نے محسوس کیا کہ اس کی باتیں بڑی دلچسپ ہوتی ہیں۔ پتہ نہیں کیسے وہ عورتوں اور آرائشی مصنوعات کا تذکرہ نکال بیٹھا تھا!

"کیا خیال ہے آپ کا یہ عورتیں سال میں کتنی لپ اسٹک کھا جاتی ہوں گی؟"

اس نے پوچھا!

"ابھی تک میں عورتوں کے معاملات سمجھنے کے قابل نہیں ہوا"۔ صفدر نے جواب دیا۔

"اوہو۔۔ تو کیا ابھی تک سنگل ہی ہو یار۔۔!"

"بالکل سنگل۔۔!"

"یہ تو بہت بری بات ہے کہ تمہاری آمدنی کا بہت بڑا حصہ لغویات پر نہیں صرف ہوتا"۔

"تم شائد بہت زیادہ زیربار ہوجاتے ہو"۔ صفدر مسکرایا۔

"دو بیویاں ہیں! لیکن ایک کو دوسری کی خبر نہیں۔۔!"

"یہ کیسے ممکن ہے؟"

"دن ایک کے ہاں گزرتا ہے، رات دوسری کے ہاں"۔ ایک سمجھتی ہے کہ میں فلموں کے لئے کہانیاں لکھتا ہوں! وہی جس کے ہاں رات بسر ہوتی ہے۔ اور دوسری سمجھتی ہے کہ میں ایک مل میں اسسٹنٹ ویونگ ماسٹر ہوں اور ہمیشہ رات کی ڈیوٹی پر رہتا ہوں"۔

"تو تم حقیقتاً کیا کرتے ہو؟"

"فلموں کے لئے کہانیاں لکھتا ہوں۔۔!" اس نے جواب دیا۔ "اور یہ کہانیاں کہیں بھی بیٹھ کر لکھی جاسکتی ہیں! اور کبھی ناوقت سیٹ پر جانا پڑا تو اس وقت والی بیوی سمجھتی ہے کہ اوورٹائم کر رہا ہوں۔ یا شوٹنگ طویل ہوگئی ہے۔۔!"

"کمال کے آدمی ہو۔۔!"

"بیویوں کو دھوکا دینا میری تفریح ہے!۔ اب تیسری کے امکانات پر غور کر رہا ہوں لیکن وقت کیسے نکالوں گا"۔

"واہ۔۔ تیسری بھی کرو گے۔۔!"

"کرنی ہی پڑے گی۔ دیکھو یار قصہ دراصل یہ ہے کہ زیادہ سے زیادہ شادیاں کرنے سے سالیوں کی تعداد میں اضافہ ہوتا ہے۔۔ اور سالیاں۔۔ ہا۔۔ اگر سالیاں نہ ہوں تو دنیا ویران ہوجائے!"

"مجھے تو اس نام ہی سے گھن آتی ہے"۔ صفدر نے کہا۔

"آہا۔ تو تم انہیں سالیوں کی بجائے بتاشیاں یا جلیبیاں کہہ لیا کرو! کیا فرق پڑتا ہے"۔

صفدر ہنسنے لگا اور تھوڑی دیر بعد یہ بھول ہی گیا کہ وہ یہاں کس لئے آیا تھا۔

کھیل ختم ہوجانے کے بعد وہ ڈائننگ روم میں آ بیٹھے۔ بھاری جبڑے والا ایک لاپرواہ اور فضول خرچ آدمی معلوم ہوتا تھا۔
کافی پیتے وقت اس نے صفدر سے کہا۔"یار مجھ پر ایک احسان کرو"۔

"کیا؟" صفدر چونک پڑا۔

اس نے کلائی کی گھڑی دیکھتے ہوئے کہا۔ "چھ بج رہے ہیں لیکن میں رات والی بیوی سے آج پیچھا چھڑانا چاہتا ہوں۔ میں اس سے کہوں گا کہ تم اسسٹنٹ ڈائریکٹر ہو۔ آج رات بھی شوٹنگ ہوگی۔ اس لئے ڈائریکٹر نے تمہیں ساتھ کردیا ہے تاکہ تم مجھے اپنے ساتھ ہی لے جاؤ! ساڑھے سات بجے ہم گھر ہی پر رات کا کھانا کھائیں گے۔ تم برابر کہتے رہنا، بھئی جلدی چلو اور بس ہم آٹھ بجے تک گھر سے نکل آئیں گے۔کیوں؟ پھر ہم دونوں دوست ہوجائیں گے۔ اور تم آئندہ بھی ایسے مواقع پر میرے کام آیا کرنا!"

صفدر ہنسنے لگا۔ مگر بھاری جبڑے والے کی سنجیدگی میں ذرہ برابر بھی فرق نہ آیا!

"میں سنجیدہ ہوں دوست!" اس نے کہا۔ "اگر تم یہ کام نہ کرسکو تو صاف جواب دو۔ تاکہ میں کسی دوسرے کو پھانسوں! بس کسی اور کے ساتھ کچھ دیر کھیلنا پڑے گا! سارے ہی آدمی تمہاری طرح ٹھس تھوڑا ہی ہوں گے۔ ایڈوینچر کا شوق کسے نہیں ہوتا! بہتیرے پھنسیں گے!"

صفدر نے سوچا چلو دیکھا ہی جائے گا کہ یہ آدمی کس حد تک بکواس کر رہا ہے اور اسے بہرحال اس کے متعلق معلومات فراہم کرنی تھیں! پہلے چوری چھپے یہ کام سرانجام دینا پڑتا۔ مگر اب۔۔ اب تو وہ اسے کھلی ہوئی کتاب کی طرح پڑھ سکے گا۔

اس نے حامی بھر لی۔

باہر نکل کر بھاری جبڑے والے نے کہا۔ "یہ تو اور اچھی بات ہے کہ تمہاری کار بھی موجود ہے! اب وہ شبہ بھی نہ کرسکے گی کہ میں اسے الو بنا رہا ہوں۔ وہ تمہارے اسسٹنٹ ڈائریکٹر پر ایمان لے آئے گی"۔

"قطعی۔!" صفدر یوں ہی بولنے کے لئے بولا۔

وہ صفدر کی رہنمائی کرتا رہا اور پھر ماڈل کالونی کی ایک دور افتادہ عمارت کے سامنے کار روکنے کو کہا۔ عمارت نہ خوبصورت تھی اور نہ بڑی تھی۔ پائیں باغ ابتر حالت میں تھا۔ جس سے مالک مکان کی لاپرواہی یا مفلوک الحالی ظاہر ہو رہی تھی!

اس نے اسے نشست کے کمرے میں بٹھایا اور خود اندر چلا گیا!

صفدر سوچ رہا تھا کہ اسے فلموں یا فلموں کی شوٹنگ کے متعلق بالکل کچھ نہیں معلوم! اگر اس کی بیوی اس سلسلے میں اس سے کچھ پوچھ بیٹھی تو کیا ہوگا۔۔!

لیکن اس کے کچھ پوچھنے سے پہلے تین چار آدمی اس پر ٹوٹ پڑے۔ حملہ پشت سے ہوا تھا۔ اس لئے اسے سنبھلنے کا موقع نہ مل سکا۔

ایک نے اس کا منہ دبالیا تھا اور دو بری طرح جکڑے ہوئے دروازے کی جانب کھینچ رہے تھے۔ لیکن جب وہ اس طرح اسے کمرے سےباہر نہ لے جاسکے تو تین مزید آدمی ان کی امداد کے لئے وہاں آپہنچے۔ اور صفدر کشاں کشاں ایک تہہ خانے میں پہنچا دیا گیا۔ تہہ خانے کا علم تو اسے اس وقت ہوجا جب اس کی آنکھوں پر سے پٹی کھولی گئی۔ بعد میں آنے والے تین آدمیوں میں سے ایک نے اس کی آنکھوں پر رومال باندھ دیا تھا اور کسی نے دونوں ہاتھ پشت پر جکڑ دیئے تھے۔

لیکن جب آنکھوں پر سے رومال کھولا گیا تو اس کے سامنے صرف ایک ہی آدمی تھا اور یہ تھا وہی بھاری جبڑے والا جو اسے ٹپ ٹاپ نائٹ کلب سے یہاں تک لایا تھا!

"مجھے افسوس ہے دوست!" اس نے سر ہلا کر مغوم لہجے میں کہا۔ "اس وقت دونوں بیویاں یہاں موجود ہیں! اس لئے یہ ابتری پھیلی ہے۔ سالیوں کی بجائے دونوں طرف کے سالے اکھٹے ہوگئے ہیں اور انہیں شبہ ہے کہ تم ہی مجھے بہکایا کرتے ہو!"

صفدر نچلا ہونٹ دانتوں میں دبائے ہوئے اسے گھورتا رہا!

وہ کوشش کر رہا تھا کہ پشت پر بندھے ہوئے ہاتھ آزاد ہوجائیں! لیکن کامیابی کی امید کم تھی۔ اگر کسی طرح وہ اپنے ہاتھ استعمال کرنے کے قابل ہوسکتا تو اس بھاری جبڑے کے زاویوں میں کچھ نہ کچھ تبدیلیاں ضرور نظر آتیں کیونکہ وہ ایک بےجگر فائٹر تھا!
دفعتاً بائیں جانب دیوار میں ایک دروازہ نما خلاء نمودار ہوئی اور جوزف جھکا ہوا اندر داخل ہوا۔ اس کے سر پر پٹی چڑھی ہوئی تھی اور اس کے دونوں ہاتھ بھی پشت پر بندھے ہوئے تھے! سر شاید زخمی تھا! شاید یہ صفدر کی چھٹی حس ہی تھی جس نے اس کے چہرے پر حیرت کے آثار نہ پیدا ہونے دیئے اور جوزف تو پہلے ہی سے سر جھکائے کھڑا ہوا تھا! اس نے کسی طرف دیکھنا بھی نہیں تھا! اس کے چہرے پر نظر آنے والے آثار اکھڑے ہوئے نشے سے پیدا ہونے والی بوریت کی غمازی کر رہے تھے۔ زیادہ دیر تک شراب نہ ملنے پر اس کی پلکیں ایسی ہی بوجھل ہوجاتی تھیں کہ وہ کسی کی طرف دیکھنے میں بھی کاہلی محسوس کرتا تھا!

اچانک بھاری جبڑے والے نے صفدر سے پوچھا۔ " یہ کون ہے؟"

"میں کیا جانوں!" صفدر غرایا۔ "کہیں تمہارا دماغ تو نہیں خراب ہوگیا!"

بھاری جبڑے والے کا قہقہہ کافی طویل تھا لیکن جوزف اب بھی سر جھکائے کسی بت کی طرح کھڑا رہا!ایسا معلوم ہو رہا تھا جیسے یہ آوازیں اس کے کانوں تک پہنچی ہی نہ ہوں۔ جو آدمی اسے یہاں لایا تھا اس کی رائفل کی نال اب بھی اس کی کمر سے لگی ہوئی تھی!

"تم بکواس کرکے کامیاب نہیں ہوسکتے دوست"۔ بھاری جبڑے والے نے کہا۔ "تم عمران کے آدمی ہو! اور اس وقت بھی اس کے ساتھ تھے۔ جب وہ ندی پر مقبرہ کے قریب گھیرا گیا تھا"۔

"مجھے اس سے کب انکار ہے مگر میں اس آدمی کو نہیں جانتا"۔ صفدر نے لاپروائی سے کہا۔

"یہ عمران کا ملازم نہیں ہے؟" بھاری جبڑے والے نے غرا کر کہا۔

"میں نے تو کبھی عمران کے ساتھ نہیں دیکھا"۔ صفدر نے جواب دیا! وہ جانتا تھا کہ جوزف اب عمران کے ساتھ اس کے فلیٹ میں نہیں رہتا بلکہ مستقل طور پر رانا پیلس ہی میں اس کا قیام ہے۔ اس لئے وہ اس کے معاملے میں محتاط ہو کر زبان کھول رہا تھا!

"رانا تہور علی کو جانتے ہو؟"

"یہ نام میرے بالکل نیا ہے"۔ صفدر نے متحیرانہ لہجے میں کہا۔

"او۔۔ حبشی۔۔!" دفعتاً وہ جوزف کی طرف مڑکر گرجا! "اب تم اپنی زبان کھولو۔ ورنہ تمہارے جسم کا ایک ایک ریشہ الگ کر دیا جائے گا"۔

"جاؤ۔۔" جوزف سر اٹھائے بغیر بھرائی سی آواز میں بولا! "پہلے میری پیاس بجھاؤ! پھر میں بات کروں گا۔ تم لوگ بہت کمینے ہو۔ تمہیں شاید نہیں معلوم کہ شراب ہی میری زبان کھلوا سکے گی"۔

"شراب نہیں مل سکے گی"۔

"تب پھر مجھے کسی کی بھی پروا نہیں! جو تمہارا دل چاہے کرو"۔

"ادھر دیکہو۔ کیا تم اس آدمی کو پہچانتے ہو؟" اشارہ صفدر کی طرف تھا۔

"کیوں دیکھوں؟ کیسے دیکھوں؟ میری آنکھوں کے سامنے غبار اڑ رہا ہے۔ مجھے اپنے پیر بھی صاف نہیں دکھائی دیتے۔ شراب لاؤ۔ یا مجھے گولی مار دو"۔

"پلاؤ۔ اسے۔ پلاؤ"۔ دفعتاً بھاری جبڑے والا دونوں ہاتھ ملا کر غرایا ۔ "اتنی پلاؤ کہ اس کا پیٹ پھٹ جائے"۔

رائفل والا جوزف کے پاس سے ہٹ کر پچھلے دروازے سے نکل گیا۔

"عمران کہاں ہے؟" وہ پھر صفدر کی طرف متوجہ ہوا۔

"اگر تم یہ جانتے ہو کہ میں اس دن عمران کے ساتھ تھا جب ہم پر چاروں طرف سے گولیاں برس رہی تھیں تو یہ بھی جانتے ہو گے کہ عمران کام آگیا تھا اور میں بچ کر نکل گیا تھا"۔

"ہمیں تو اس پر یقین تھا کہ تم بھی نہ بچے ہوگے! لیکن آج تم ہاں میرے سامنے موجود ہو! تم اتنی چالاکی سے نکل گئے تھے کہ ہمیں پتہ ہی نہ چل سکا تھا"۔

"عمران گولی کھا کر دریا میں گر گیا تھا"۔ صفدر نے بھرائی ہوئی آواز میں کہا! لیکن وہ ڈر رہا تھا کہ کہیں جوزف یہ جملے سن کر چونک نہ پڑے۔ اس وقت کی گفتگو سے اچھی طرح اندازہ کرچکا تھا کہ وہ رانا تہور علی اور عمران کی الجھن میں پڑ گئے ہیں۔
لیکن صفدر کے اندیشے بےبنیاد ثابت ہوئے کیونکہ جوزف کے کانوں پر جوں تک نہیں رینگی تھی اس نے نہ تو سر اٹھایا اور نہ کسی طرف دیکھا۔

تھوڑی دیر بعد قدموں کی آہٹ سنائی دی اور رائفل والا دیسی شراب کی دو بوتلیں لئے دروازے سے اندر داخل ہوا۔

"ایک بوتل کھول کر اس کے منہ سے لگا دو"۔ بھاری جبڑے والے نے کہا۔ تعمیل کی گئی! جوزف کے موٹے موٹے ہونٹ بوتل کے منہ سے چپک کر رہ گئے! بڑا مضحکہ خیز منظر تھا۔ ایسا ہی لگ رہا تھا کہ جیسے کسی بھوکے شیرخوار بچے نے دودھ کی بوتل سے منہ لگا کر چسر چسر شروع کردی ہو۔

آدھی بوتل غٹا غٹ پی جانے کے بعد اس نے بوتل کا منہ چھوڑ کر دو تین لمبی لمبی سانسیں لیں اور مسکرا کر بولا۔

"تم بڑے اچھے ہو! بڑے پیارے آدمی ہو! تم پر آسمان سے برکتیں نازل ہوتی رہیں! اور آسمانی باپ تمہیں اچھے کاموں کی توفیق دے"۔

بھاری جبڑے والا کینہ توز نظروں سے اسے دیکھ رہا تھا۔ ایسا معلوم ہو رہا تھا جیسے سالہا سال سے اسے مار ڈالنے کی خواہش پال رہا ہو! جوزف نے بقیہ آدھی بوتل بھی ختم کردی!

اب وہ کسی جاگتے ہوئے آدمی کی سی حالت میں آگیا تھا۔ آنکھیں سرخ ہوگئیں تھیں اور چہرے کی سیاہی چمکنے لگی تھی!

"ارے۔۔ یہ آدمی۔۔" دفعتاً اس نے بھرائی ہوئی آواز میں کہا۔ "ہاں! مجھے یاد پڑتا ہے کہ میں نے اسے ایک آدھ بار مسٹر عمران کے ساتھ دیکھا تھا"۔

"لیکن میں نے تو تمہیں کبھی نہیں دیکھا"۔ صفدر نے غصیلی آواز میں کہا۔

"یہ بھی ممکن ہے مسٹر کہ تمہاری نظر مجھ پر کبھی نہ پڑی ہو"۔

"عمران کہاں ملے گا؟" بھاری جبڑے والا غرایا۔

"میں کیا بتا سکتا ہوں مسٹر"۔ جوزف نے متحیرانہ انداز میں پلکیں جھپکائیں۔ "بہت دنوں کی بات ہے جب میں مسٹر عمران کے ساتھ تھا۔ لیکن وہ میرے پینے پلانے کا بار سنبھالنے کی حیثیت نہیں رکھتے تھے۔ اس لئے انہوں نے خود ہی میرا پیچھا چھوڑ دیا۔۔ اس طرح میں نے اطمینان کا سانس لیا! ورنہ مجھے تو اس کا غلام رہنا ہی پڑتا ہے جو مجھے زیر کرلے۔ اور پھر میرا تو ڈاکٹر طارق والا مقدمہ بھی چل رہا ہے"۔

"کیسا مقدمہ۔۔؟"

اس پر جوزف نے ڈاکٹر طارق کی کہانی دہراتے ہوئے کہا۔ ماسٹر عمران نے مجھے بہت پیٹا تھا۔ وہ شاید پولیس کے لئے کام کرتے ہیں۔۔!"

بھاری جبڑے والا تھوڑی دیر تک کچھ سوچتا رہا پھر بولا! "رانا کون ہے؟"

"باس ہے میرا۔ جوزف نے فخر سے سینہ تان کر کہا۔

"وہ کہاں ملے گا۔۔؟"

"میں نہیں جانتا۔ ان سے تو بس کبھی کبھی ملاقات ہوتی ہے"۔

"عمران سے اس کا کیا تعلق ہے۔۔؟"

"میں کیا بتاسکتا ہوں مسٹر۔ میں کیا جانوں! میں نے کبھی ان کے ساتھ مسٹر عمران کو نہیں دیکھا"۔

"تم رانا کے پاس کیسے پہنچے تھے؟"

"بس یوں ہی میں ایک دن سڑک پر جا رہا تھا کہ ایک کار میرے پاس رکی! اس پر سے رانا صاحب اترے اور کہنے لگے میں نے پچھلے سال شاید تمہیں نیٹال میں دیکھا۔ میں نے کہا کہ میں تو دس سال سے اس ملک میں ہوں! انہوں نے کہا کہ ہوسکتا ہے ان کے ذہن میں کوئی اور ہو۔ پھر وہ مجھ سے میرے متعلق پوچھ گچھ کرنے لگے!۔۔ یہ۔۔ دوسری بوتل بھی مسٹر۔۔ خدا تمہیں ہمیشہ خوش رکھے اور عورت کے سائے سے بچائے۔ تم بہت نیک ہو"۔

بھاری جبڑے والے کے اشارے پر دوسری بوتل بھی کھولی گئی! اور جوزف چوتھائی پینے کے بعد بولا۔ "ہاں تو تم کیا پوچھ رہے تھے۔ برادر۔۔!"

"تم رانا کے پاس کیسے پہنچے تھے؟"

"ہاں۔۔ ہاں۔۔ شاید میں یہی بتا رہا تھا کہ وہ مجھ سے میرے بارے میں پوچھ گچھ کرنے لگے!"

"چلو کہتے رہو! رکو مت!" بھاری جبڑے والا بولا۔

"میں نے انہیں بتایا کہ مجھے نوکری کی تلاش ہے۔ انہوں نے پوچھا باڈی گارڈز کے فرائض انجام دے سکو گے! اوہ۔۔ بڑی آسانی سے میں نے انہیں بتایا اور یہ بھی کہا کہ میرا نشانہ بڑا عمدہ ہے اور میں کبھی ہیوی ویٹ چیمپین بھی رہ چکا ہوں۔ وہ بہت خوش ہوئے اور مجھے نوکر رکھ لیا! میں ان کے پیسنے کی جگہ خون بھی بہا سکتا ہوں۔ لارڈ آدمی ہیں۔ کبھی نہیں پوچھتے کہ میں دن بھر میں کتنی بوتلیں صاف کردیتا ہو"۔

بھاری جبڑے والا پھر کسی سوچ میں پڑ گیا! ایسا معلوم ہو رہا تھا جیسے وہ اس کے بیان پر تذبذب میں پڑ گیا ہو۔

دوسری طرف صفدر پر جوزف کے جوہر پہلی بار کھلے تھے! وہ اب تک اسے پرلے درجے کا ایڈیٹ ہی تصور کرتا رہا تھا! لیکن اس وقت تو عمران ہی کا یہ قول کرسی نشین ہوا تھا کہ جوزف ایک نادر الوجود شکاری کتا ہے۔ سادہ لوحی اور چیز ہے! لیکن بےضرر نظر آنے والے کتے بھی شکار کے وقت اپنی تمام تر صلاحیتوں سے کام لیتے ہیں! بشرطیکہ وہ شکاری ہوں! جوزف پر صحیح معنوں میں یہ مثال صادق آتی تھی۔

"دیکھو میں تمہاری ہڈیاں چور کردوں گا۔ ورنہ مجھ سے اڑنے کی کوشش نہ کرو"۔

"بس یہ بوتل ختم کر لینے دو! اس کے بعد جو دل چاہے کرنا!" جوزف نے ہونٹ چاٹتے ہوئے کہا۔

"صرف ایک دن کی مہلت اور دی جاتی ہے۔ تم عمران کا پتہ بتاؤ اور تم رانا تہور علی کا۔۔!" بھاری جبڑے والا ہاتھ اٹھا کر بولا۔

وہ رائفل والے کو اپنے پیچھے آنے کا اشارہ کرتا ہوا دروازے سے نکل گیا اور پھر وہ دروازہ بھی غائب ہوگیا۔ دیوار برابر ہوگئی تھی۔

جوزف دوسری بوتل کی طرف ندیدوں کی طرح دیکھنے لگا جس میں ابھی تین چوتھائی شراب باقی تھی۔ اس پر کاگ بھی نہیں تھا۔

وہ تھوڑی دیر تک حسرت بھری نظروں سے اسے دیکھتا رہا پھر پشت پر بندھے ہوئے ہاتھوں کے بل فرش پر نیم دراز ہوگیا! دیکھتے ہی دیکھتے بوتل دونوں پیروں میں دبائی اور پیر سر کی طرف اٹھنے لگے۔۔ اور بوتل کا منہ اس کے ہونٹوں سے جالگا!

صفدر کھڑا پلکیں چھپکاتا رہا! "غٹ غٹ" کی صدائیں تہہ خانے کے سکوت میں گونج رہی تھیں۔ بوتل خالی ہوئے بغیر ہونٹوں سے نہ ہٹ سکی۔

دفعتاً کھٹاکے کی آواز آئی اور بھاری جبڑے والا پھر اندر داخل ہوا اس بار اس کے اس کے ہاتھ میں چمڑے کا چابک تھا! نہ جانے کیوں جوزف مسکرا پڑا مگر وہ جوزف کی طرف متوجہ نہیں تھا!

"سرسوکھے رام کو عمران کی تلاش کیوں ہے؟" اس نے صفدر سے پوچھا!

"میں نہیں جانتا"۔

"تم جانتے ہو۔۔!" وہ چابک زمین پر مارتا ہوا دہاڑا۔

"میرے ہاتھ کھول دو۔ پھر اس طرح اکڑوں تو یقیناً مرد کہلاؤ گے"۔

اس بار چابک صفدر کے جسم پر پڑا اور وہ تلملا گیا۔

"بتاؤ!"

صفدر اس کی طرف جھپٹا لیکن اس نے اچھل کر پیچھے ہٹتے ہوئے پھر چابک گہما دیا! اس طرح صفدر نے کئی چابک کھائے! اور یک بیک سست پڑ گیا! یہ حماقت ہی تو تھی کہ وہ اس طرح پٹ رہا تھا! ادھر جوزف کا یہ حال تھا کہ وہ کوشش کے باوجود بھی فرش سے نہیں اٹھ سکتا تھا! پورے چھتیس گھنٹوں کے بعد اسے شراب ملی تھی اور اس نے یہ دو بوتلیں جس طرح ختم کی تھیں اس طرح کوئی دوسرا پانی بھی نہ پی سکتا!

"میں نہیں جانتا۔۔!"

"ڈھمپ اینڈ کو کا اصل بزنس کیا ہے؟"

"فارورڈنگ اینڈ کلیرنگ۔۔!"

"تم وہاں کام کرتے ہو؟"

"ہاں۔۔!"

"پھر عمران کا اور تمہارا کیا ساتھ۔۔؟"

"مجھے شوق ہے سراغرسانی کا"۔ صفدر بولا۔ "عمران کی وجہ سے میں بھی اپنا یہ شوق پورا کرسکتا ہوں کیونکہ وہ پولیس کے لئے کام کرتا ہے"۔

"تمہارے دفتر کی اسٹینو ٹائپسٹ جولیا کا عمران سے کیا تعلق ہے؟"

"یہ وہی دونوں بتا سکیں گے!" صفدر نے ناخوشگوار لہجے میں کہا۔

بھاری جبڑے والا کھڑا دانت پیستا رہا۔ پھر آنکھیں نکال کر آہستہ آہستہ بولا۔ "تم مجھے نہیں جانتے! میں تمہارے فرشتوں سے بھی اگلوالوں گا! خواہ اس کے لئے تمہارا بند بند بھی کیوں نہ الگ کرنا پڑے۔۔!"

وہ پیر پٹختا ہوا چلا گیا! دیوار کی خلاء اس کے گذرتے ہی پر ہوگئی تھی! ایک تختہ سا بائیں جانب کھسک کر دوسری جانب کی دیوارسے جا ملتا تھا!




\endR

\bigskip\bigskip\bigskip

\centerline{\hdr{(۹)}}
\bigskip\bigskip\bigskip
\beginR
\body



جیسے ہی جولیا کی نظر سرسوکھے پر پڑی وہ ستون کی اوٹ میں ہوگئی۔ یہاں پام کا بڑا گملا رکھا ہوا تھا اور پام کے پتے اسے چھپانے کے لئے کافی تھے۔

وہ سرسوکھے سے بھاگنے لگی تھی! کیونکہ وہ اسے بےحد بور کرتا تھا! وہ پرانی کہانی جس کا سلسلہ میں وہ عمران کا تعاون حاصل کرنا چاہتا تھا بار بار دہرائی جاتی! اور پھر اس کے ساتھ سرسوکھے کی اداسی بھی تو تھی! اسے غم تھا کہ اس کے آگے پیچھے کوئی نہیں ہے۔ کوئی ایسا نہیں ہے جسے وہ اپنا کہہ سکے! جوانی ہی میں موٹاپا شروع ہوگیا تھا اور اسی بنا پر خود اس کی پسند کی لڑکیاں اسے منہ لگانا پسند نہیں کرتی تھیں۔۔ وہ جولیا سے یہ ساری باتیں کہتا رہتا! ٹھنڈی سانسیں بھرتا اور کبھی کبھی اس کی آنکھوں میں آنسو تیرنے لگتے! جنہیں وہ چھپانے کے لئے وہ طرح طرح کے منہ بناتا! اور ہزاروں قہقہے جولیا کے سینے میں طوفان کی سی کیفیت اختیار کرلیتے پھر اسے کسی بہانے سے اس کے پاس سے اٹھ جانا پڑتا۔۔ وہ کسی باتھ روم میں گھس کر پیٹ دبا دبا کر ہنستی۔۔! اکثر سوچتی کہ اسے تو اس سے ہمدردی ہونی چاہیئے! پھر آخر اسے اس پر تاؤ کیوں آتا ہے۔۔! وہ غور کرتی تو سرسوکھے کی زندگی اسے بڑی دردناک لگتی! لیکن زیادہ سوچنے پر اسے یا تو ہنسی آتی یا غصہ آتا! کبھی وہ سوچتی کہ کہیں سرسوکھے اس کام کے بہانے اس سے قریب ہونے کی کوشش تو نہیں کر رہا! اس خیال پر غصے کی لہر کچھ اور تیز ہوجاتی! مگر پھر کچھ دیر بعد ہی اس شام کا خیال آجاتا جب وہ اس کے دفتر میں بیٹھی سونے کی اسمگلنگ کی کہانی سن رہی تھی اور دوسرے کمرے کی میز الٹنے کی آواز نے انہیں چونکا دیا تھا! اور پھر اس نے میز کی سطح پر پیروں کے نشانات محفوظ کئے تھے۔۔! وہ سوچتی رہی اور اس نتیجے پر پہنچی کہ وہ حقیقتاً پریشانیوں میں مبتلا ہے۔ یہ اور بات ہے کہ ہر قسم کی پریشانیوں کا تذکرہ بیک بوقت کردینے کا عادی ہو!

وہ روزانہ شام کو عمران کی تلاش میں نکلتے تھے! لیکن آج کے لئے جولیا نے ایک ضروری کام کا بہانہ کرکے اس سے معافی مانگ لی تھی۔۔! لیکن وہ گھر میں نہ بیٹھ سکی! شام ہوتے ہی اس نے سوچا آج تنہا نکلنا چاہیئے! مقصد عمران کی تلاش کے علاوہ اور کچھ نہیں تھا! وہ ٹپ ٹاپ نائٹ کلب کے پورچ میں پہنچی ہی تھی کہ اچانک غیر متوقع طور پر سرسوکھے نظر آگیا تھا! وہ سوچ بھی نہیں سکتی تھی کہ آج وہ بھی وہیں آ مرے گا۔

جیسے ہی وہ پورچ میں پہنچا! جولیا گملے کی آڑ سے نکلی اور جھپٹ کر کلرک روم میں د اخل ہوگئی! یہاں سے ایک راہداری براہ راست ریکریئشن ہال میں جاتی تھی! جہاں آج اسکیٹنگ کا پروگرام تھا۔۔!

وہ بڑی بدحواسی کے عالم میں یہاں پہنچی!

"اف خدا۔۔" وہ بڑبڑائی اور اس کا سر چکرا گیا! کیونکہ سرسوکھے دوسرے دروازے سے ریکریئشن ہال میں داخل ہوا تھا! ویسے اس کی توجہ جولیانا کی طرف نہیں تھی! جولیانا کلرک روم والی راہداری ایک گیلری میں لائی تھی۔ اس نے ذہنی انتشار کے دوران فیصلہ کیا کہ سرسوکھے سے تو کھوپڑی نہیں چٹوائے گی خواہ کچھ ہوجائے۔ پھر؟

وہ جھپٹ کر ایک میز پر جا بیٹھی جہاں ایک اداس آنکھیوں والا نوجوان پہلے ہی سے موجود تھا۔

"معاف کیجیئے گا!" جولیا نے کہا۔ ذرا سر چکرا گیا ہے۔۔۔ ابھی اٹھ جاؤں گی"۔

"کوئی بات نہیں محترمہ!" وہ بھرائی ہوئی آواز میں بولا۔

جولیا نے آنکھوں پر رومال رکھ کر سرجھکا لیا اور چڑھتی ہوئی سانسوں پر قابو پانے کی کوشش کرنے لگی۔۔!

"کیسی طبعیت ہے۔۔ آپ کی؟" تھوڑی دیر بعد نوجوان نے پوچھا!

"اوہ۔۔ جی ہاں۔۔ بس ٹھیک ہی ہے۔۔ اب۔۔!"

"برانڈی منگواؤں۔۔!"

"جی نہیں شکریہ! میں اب بالکل ٹھیک ہوں!" وہ سر اٹھا کر بولی۔

"آج کل موسم بڑا خراب جا رہا ہے!" نوجوان بولا۔

"جی ہاں۔۔ جی ہاں۔۔ یہی بات ہے"۔

یہ دبلے چہرے والا مگر وجیہہ نوجوان تھا! اس کی آنکھوں کی غم آلود نرماہٹ نے اسے کافی دلکش بنادیا تھا۔ پیشانی کی بناوٹ بھی نرم دلی اور اور ایمانداری کا اعلان کر رہی تھی۔۔!

"میں اس شہر میں نوارد ہوں"۔ جولیا نے کہا۔ "مجھے نہیں معلوم تھا کہ یہاں اسکیٹنگ بھی ہوتی ہے! مجھے بےحد شوق ہے۔ اس کا!"

"جی ہاں"۔ اس نے تھکی ہوئی سی مسکراہٹ کے ساتھ کہا۔ "دلچسپ کھیل ہے"۔

"آپ کو پسند ہے؟"

"بہت زیادہ۔۔!" نوجوان کا لہجہ بےحد غم انگیز تھا۔۔!

ٹھیک اسی وقت سرسوکھے ان کے قریب پہنچا! جولیا کی نظر غیر ارادی طور پراس کی طرف اٹھ گئی تھ اور وہ بطور اعتراف شناسائی سر کو خفیف سی جنبش دے کر آگے بڑھ گیا تھا! جولیا بھی بادل ناخواستہ مسکرائی تھی۔

بہرحال اس کے اس طرح آگے بڑھ جانے پر اس کی جان میں جان آئی تھی وہ اس پر یہ بھی نہیں ظاہر کرنا چاہتی تھی کہ اس سے بچنے کی کوشش کر رہی ہے! سرسوکھے آگے بڑھ کر ایک میز پر جا بیٹھا تھا! جولیا سوچ رہی تھی کہ اگر وہ اس میز سے اٹھی اور سرسوکھے کو شبہ بھی ہوگیا کہ وہ تنہا ہے تو وہ تیر کی طرح اس کی طرف آئے گا۔

اتنے میں اسکیٹنگ کے لئے موسیقی شروع ہوگئی! اور جولیا نے اس انداز میں نوجوان کی طرف دیکھا جیسے مطالبہ کر رہی ہو کہ مجھ سے رقص کی درخواست کرو! مگر نوجوان خالی آنکھوں سے اس کی طرف دیکھتا رہا۔۔!

جولیا نے سوچا بدھو ہے لہذا اس نے خود ہی کہا! "اگر آپ کو اسکیٹنگ سے دلچسپی ہے ۔۔ تو۔۔ آئیے۔۔!"

"میں۔۔!" نوجوان کے لہجے میں تحیر تھا! پھر اس کی آنکھوں کی اداسی اور گہری ہوگئی۔۔! اس نے چھبتے ہوئے لہجے میں پوچھا۔ "آپ میرا مذاق کیوں اڑا رہی ہیں محترمہ؟"

"میں نہیں سمجھی!" جولیا بوکھلا گئی!

"ٌیا آپ یہ بیساکھی نہیں دیکھ رہی ہیں!" اس نے ایک کرسی سے ٹکی ہوئی بیساکھی کی طرف اشارہ کیا۔

جولیا کی نظریں اگر پہلے اس پر پڑی بھی ہوگی تو اس نے دھیان نہ دیا ہوگا! بہرحال اب وہ کٹ کر رہ گئی!

"اوہ۔۔ معاف کیجیئے گا!" اس نے لجاجت سے کہا۔ " میں نے خیال نہیں کیا تھا میں بےحد شرمندہ ہوں جناب! کیا آپ معاف نہیں کریں گے؟"

"کو ئی بات نہیں!" وہ ہنس پڑا۔

اس کا بایاں پیر شاید کسی حادثے کی نظر ہو کر گھٹنے کے پاس سے کاٹ دیا گیا تھا اور اب لکڑی کا ایک ڈھانچہ پنڈلی کا کام دے رہا تھا۔

"یہ کیسے ہوا تھا؟" جولیا نے پوچھا۔ وہ سچ مچ اس کے لئے غمگین ہوگئی تھی!

"فوجیوں کی زندگی میں ایسے حادثات کوئی اہمیت نہیں رکھتے"۔ اس نے کہا اور بتایا کہ وہ پچھلی جنگ عظیم میں اطالویوں کے خلاف لڑا تھا اور مورچے پر ہی اس کی بائیں ٹانگ ایک حادثہ کا شکار ہوگئی تھی! وہ سیکنڈ لیفٹنٹ تھا!

بات لمبی ہوتی گئی اور وہ جنگ کے تجربات بیان کرتا رہا۔ تھوڑی ہی دیر بعد جولیا نے محسوس کیا کہ اب اس میز سے اٹھنے کا سوال ہی نہیں پیدا ہوسکتا! اس کے بعد بھی وہ تھوڑی دیر تک ادھر اُدھر کی گفتگو کرتے رہے۔ پھر پہلا دور ختم ہوگیا۔۔! نوجوان نے کافی منگوائی اور جولیا کو انکار کے باوجود بھی پینی ہی پڑی! ویسے بھی وہ اس مغوم نوجوان کی درخواست رد نہیں کرنا چاہتی تھی۔

کچھ دیر بعد کسی جانب سے ایک خوبصورت اور صحت مند نوجوان ان کی طرف آیا اور جولیا سے ساتھی بننے کی درخواست کی۔ جولیا اس کی آواز سن کر چونک پڑی۔

"اگر کوئی حرج نہ ہو تو۔۔!" وہ کہہ رہا تھا!

"ضرور۔۔ ضرور۔۔!" جولیا مسکراتی ہوئی اٹھ گئی تھی! ساتھ ہی اس نے لنگڑے نوجوان کی طرف دیکھ کر سر ہلایا اور یہ بھی محسوس کیا تھا کہ وہ کھسیا سا گیا ہے لیکن یہ کیسے ممکن تھا کہ وہ اس آدمی کی درخواست رد کردیتی جس کے لئے خود اتنے دنوں سے بھٹکتی پھر رہی تھی! صورت سے تو وہ اسے ہرگز نہ پہچان سکتی کیونکہ وہ میک اپ میں تھا لیکن جب اپنی اصلی آواز میں بولا تھا تو جولیا اسے کیوں نہ پپہچان لیتی وہ عمران کے علاوہ اور کوئی نہیں ہوسکتا تھا!

وہ اس جگہ آئے جہاں اسکیٹس ملتے تھے! جلدی جلدی انہیں جوتوں سے باندھا اور چوبی فرش پر پھسل آئے! عمران اس کے دونوں ہاتھ پکڑے ہوئے تھا!

"تم کہاں تھے درندے؟" جولیا نے پوچھا!

"شکار پر۔۔!" عمران نے جواب دیا! پھر بولا۔ "تم اس شام ندی پر کیوں دوڑی آئی تھیں؟"

"یہ اطلاع دینے کےلئے کہ تمہاری موت پر کرائے کے رونے والے بھی نہ مل سکیں گے!"

"لیکن میں تمہیں اس وقت یہ اطلاع دینا چاہتا ہوں کہ تمہارا پورا دفتر ان لوگوں کی نظروں میں آگیا ہے"۔

"پھر کیا کرنا چاہیے؟"

"پرواہ مت کرو!" لیکن فی الحال یہ بھول جاؤ کہ تمہارے ساتھ کبھی کوئی عمران بھی تھا! میں نے انہیں شبہے میں مبتلا کردیا ہے۔ کبھی انہیں میری موت پر یقین سا آنے لگتا ہے اور کبھی وہ میری تلاش شروع کردیتے ہیں"۔

"ایک آدمی اور بھی تمہاری تلاش میں ہے"۔ جولیا نے کہا اور سرسوکھے کا واقعہ بتایا۔

"فی الحال میں اس کے لئے کچھ نہیں کرسکتا!"

"ایکس ٹو تو اس کے کیس میں دلچسپی لے رہا ہے اور میں بڑی شدت سے بور ہو رہی ہوں"۔

"ہوسکتا ہے وہ اس لئے دلچسپی لے رہا ہو کہ تم میری تلاش جاری رکھو! خوب بہت اچھے یہ ایکس ٹو تو یقیناً بھوت ہے وہ شاید مجرموں پر یہی ظاہر کرنا چاہتا ہے کہ عمران کے ساتھیوں کو بھی اس کی موت پر یقین نہیں آیا۔۔ اچھا جولیا تم دن میں تین چار بار میرے فون نمبر پر رنگ کرکے سلیمان سے میرے متعلق پوچھتی رہو! میرا خیال ہے کہ وہ لوگ میرا فون ٹیپ کر رہے ہیں! سرسوکھے کے ساتھ مل کر میری تلاش بھی جاری رکھو!"

"اس کی رام کہانیاں مجھے بور کرکے مار ڈالیں گی!"

"اگر تم اتنی آسانی سے مر سکو تو کیا کہنے ہیں!" عمران نے کہا اور جولیا نے اسے لاکھوں سلواتیں سنا ڈالیں۔

وہ کچھ دیر خاموشی سے اسکیٹنگ کرتے رہے پھر جولیا نے کہا۔

"سرسوکھے یہیں موجود ہے۔۔!"

"کہاں۔۔؟"

جولیا نے بتایا! عمران کنکھیوں سے موٹے آدمی کی طرف دیکھتا ہوا بولا۔ "یہ تو صحیح معنوں میں پہاڑی معلوم ہوتا ہے کیا تم اس کے ساتھ اسکیٹنگ نہیں کرو گی؟"

جولیا نے اسے بتایا کہ کس طرح اس سے پیچھا چھڑانے کے لئے وہ ایک لنگڑے آدمی کے پاس جا بیٹھی تھی!

"بہت بری بات ہے۔۔! موٹاپا اپنے بس کی بات نہیں"۔ عمران نے مغوم لہجہ میں کہا! "تمہیں اس سے شادی کر لینی چاہیئے!"

"میں تمہارا گلا گھونٹ دوں گی۔۔!" جولیا جھلا گئی۔

"آج کل تو سب ہی مجھے مار ڈالنے کی تاک میں ہیں۔۔ ایک تم بھی سہی"۔

جولیا نچلا ہونٹ دانتوں میں دبائے اسکیٹنگ کرتی رہی۔۔! اس غیر متوقع ملاقات سے پہلے اس کے ذہن میں عمران کے متعلق ہزاروں باتیں تھیں جنہیں اس وقت قدری طور پر اس کی زبان میں آنا چاہیئے تھا! لیکن وہ محسوس کر رہی تھی کہ اب اس کے پاس جھنجھلاہٹ کے علاوہ اور کچھ نہیں رہ گیا! ویسے یہ اور بات ہے کہ اس جھنجھلاہٹ کو بھی اظہار کے لئے الفاظ نہ ملتے۔۔!

تو گویا یہ عمران اس کے لئے سوہان روح بن کر رہ گیا تھا! اس کی عدم موجودگی اس کے لئے بےچینی اور اضطراب کا باعث بنتی تھی! لیکن جہاں مشکل نظر آئی تاؤ آگیا۔۔ وہ تاؤ لانے والی باتیں ہی کرتا ھا۔۔!

جولیا کا ذہن بہک گیا تھا اور وہ کسی ننھی سی بچی کی طرف سوچ رہی تھی! یہ بھول گئی تھی کہ وہ کون ہے اور کن ذہنی بلندیوں پر رہتی ہے!

"غالباً۔۔ تم میرے فیصلے پر نظرثانی کر رہی ہو"۔ عمران نے کچھ دیر بعد مسکرا کر کہا!

"کیا مطلب۔۔؟"

"یہی کہ تمہیں سرسوکھے سے شادی کر ہی لینی چاہیئے!" عمران نے سنجیدگی سے کہا۔ "ہوسکتا ہے اس کے بعد ہی وہ صحیح معنوں میں سرسوکھے کہلانے کا مستحق ہوسکے!"

جولیا نے جھٹکا دے کر اپنے ہاتھ اس سے چھڑا لیے اور تھوڑا سا کترا کر تنہا پھسلتی چلی گئی!




\endR

\bigskip\bigskip\bigskip

\centerline{\hdr{(۱۰)}}
\bigskip\bigskip\bigskip
\beginR
\body



گیارہ بجے وہ گھر پہنچی! سرسوکھے سے اس کی گفتگو نہیں ہوئی تھی۔ کیونکہ وہ ٹپ ٹاپ کلب میں زیادہ دیر نہیں بیٹھا تھا!۔۔ جولیا تنہا اسکیٹنگ کرتی رہی تھی! لیکن جب اس نے تقریباً دس منٹ بعد دوبارہ عمران کی تلاش شروع کی تو معلوم ہوا کہ وہ بھی ہال میں موجود نہیں ہے پھر اب وہ وہاں ٹھہر کر کیا کرتی!

گھر پہنچی تو قفل کھولتے وقت کاغذ کی کھڑکھڑاہٹ محسوس ہوئی اور قفل سے ایک رول کیا ہوا کاغذ کا ٹکڑا پھنسا ہوا ملا۔
جولیا نے اسے کھینچ کر ٹارچ کی روشنی میں دیکھا!

اس پر پنسل کی تحریر نظر آئی!

"جولیا ! جب بھی واپس آؤ! فوراً مجھے رنگ کرو"۔

صفدر۔"

"کیا مصیبت ہے؟" وہ تھکے تھکے سے انداز میں بڑبڑائی تھی۔

دروازہ کھول کر وہ خواب گاہ میں آئی یہیں فون تھا! اس پر صفدر کے نمبر رنگ کئے۔

"ہیلو۔۔ کون۔۔ جولیا! دوسری طرف سے آواز آئی! "اوہ۔۔ بس میں تو صرف یہ معلوم کرنا چاہتا تھا کہ تم کب گھر پہنچتی ہو؟"

"کیوں؟"

"چند بہت ہی اہم باتیں ہیں۔ میں وہیں آرہا ہوں! پہچنے میں زیادہ سے زیادہ پندرہ منٹ لگیں گے!"

جولیا نے برا سا منہ بنا کر سلسلہ منقطع کردیا! وہ اب صرف سونا چاہتی تھی لیکن صفدر اتنی رات گئے اس سے کیوں ملنا چاہتا ہے؟

وہ اس کا انتظار کرنے لگی۔۔ پھر صفدر وعدہ کے مطابق پندرہ منٹ کے اندر ہی اندر وہاں پہنچ گیا تھا۔

"کیوں۔۔ اتنی رات گئے؟" جولیا نے متحیرانہ انداز میں پوچھا۔

"صرف ایک بات معلوم کرنا چاہتا ہوں کہ سرسوکھے رام کون ہے اور عمران کو کیوں تلاش کر رہا ہے"۔

"کیوں معلوم کرنا چاہتے ہو؟" یہ سوال غیر ارادی طور پر ہوا تھا۔

"کیونکر کچھ لوگ مجھ سے معلوم کرنا چاہتے ہیں"۔

صفدر نے اپنی کہانی چھیڑ دی۔

"مگر پھر تم یہاں کیسے نظر آرہے ہو"۔ جولیا نے اس کے خاموش ہوجانے پر پوچھا!

"یہ جوزف جیسے گدھے کا کارنامہ ہے! واقعی عمران کا انتخاب بھی لاجواب ہوتا ہے"۔

"مگر میں نے سنا ہے وہ اب عمران کے ساتھ نہیں رہتا!"

"اسی پر تو حیرت ہے!" صفدر نے کہا! حالانکہ اسے ذرہ برابر بھی حیرت نہیں تھی کیونکہ وہ جوزف کی جائے قیام سے اچھی طرح واقف تھا! لیکن ایکس ٹو کی ہدایت کے مطابق اسے پراسرار رانا پیلس کو راز ہی رکھنا تھا!

"خیر تو پھر تم لوگ رہا کیسے ہوئے؟" جولیا نے پوچھا۔

"جوزف نے ایک خالی بوتل پیروں میں دبا کر دیوار پر کھینچ ماری تھی اور پھر اس کا نیک ٹکڑا دانتوں میں دبائے ہوئے میرے پاس آیا تھا۔ ہم دونوں ہی کے ہاتھ پشت پر بندھے ہوئے تھے۔ اس نے اسی شیشے کے ٹکڑے سے میرے ہاتھوں کی ڈور کاٹنی شروع کردی! وہ شیشے کا ٹکڑا منہ میں دبائے کسی نہ تھکنے والے جانور کی طرح اپنے کام میں مشغول رہا۔ آخرکار اسے کامیابی ہی ہوئی۔ رسی کٹتے ہی میرے ہاتھ آزاد ہوگئے! پھر میں نے جوزف کے ہاتھ بھی کھول دیئے لیکن اس خدشے کی بنا پر کچھ دیر پریشان بھی ہونا پڑا کہ کہیں کوئی آ نہ جائے۔ اب ہاتھ پر ہاتھ رکھے بیٹھے رہنا بھی ہمیں کھل رہا تھا اس لئے تہہ خانے سے باہر نکلنے کے سلسلے میں ہم نے اپنی جدوجہد تیز کردی۔ ہمیں وہاں کسی ایسی چیز کی تلاش تھی جس سے دیوار میں دروازہ نما خلاء پیدا کی جاسکتی!"

جولیا کچھ نہ بولی! صفدر نے ایک سگریٹ سلگایا اور دو تین ہلکے ہلکے کش لئے!

لیکن نہ جانے کیوں وہ سوالیہ انداز میں جولیا کی طرف دیکھ رہا تھا۔۔!

کچھ دیر بعد اس نے کہا۔ "یہ ناممکن ہے کہ عمران تم سے نہ ملا ہو"۔

"ابھی تمہاری پچھلی بات پوری نہیں ہوئی"۔ جولیا ناخوشگوار لہجے میں بولی۔

"پھر کوئی بات ہی نہیں رہ گئی تھی! ہم جلد ہی اس دروازے کے میکنزم کا پتہ لگانے میں کامیاب ہوگئے! تہہ خانے کے اوپر۔۔ عمارت سنسان پڑی تھی! کسی جگہ بھی روشنی نہ دکھائی دی۔ وہ لوگ موجود نہیں تھے! ایک کھڑکی سے میں نے کمپاؤنڈ میں جھانکا۔ باہر ایک آدمی موجود تھا اور برآمدے کا بلب روشن تھا! اس آدمی نے چوکیداروں کی سی وردی پہن رکھی تھی! جوزف کسی بلی کی طر برآمدے میں رینگ گیا۔ کمال کا پھرتیلا آدمی ہے۔۔ بالکل کسی تیندوے کی طرح اور تیزی سے جھپٹنے والا! چوکیدار کے حلق سے ہلکی سی آواز بھی نہیں نکل سکی تھی! پھر جلد ہی وہ اپنے ہوش وحواس کھو بیٹھا تھا۔۔ اس طرح ہم وہاں سے فرار ہونے میں کامیاب ہوئے تھے"۔

"پھر کیا کیا تم نے۔۔؟"

"کچھ بھی نہیں! میں اپنی ذمہ داری پر کوئی قدم نہیں اٹھا سکتا"۔

"جولیا نے کچھ کہے بغیر ایکس ٹو کے نمبر ڈائیل کئے۔۔!

اور دوسری طرف سے آواز آئی۔ "دانش منزل پلیز"۔

عمران نے حال ہی میں ایکس ٹو کے پرائیویٹ فون سے ایک ٹیپ ریکارڈ اٹیچ کردیا تھا اور اس کا سسٹم کچھ اس قسم کا تھا کہ رنگ کرنے والے کو ادھر سے ریسور اٹھاے بغیر ہی جواب مل جاتا تھا! اس میں مختلف قسم کے احکامات تھے۔ آج کل کے ٹیپ پر "دانش منزل پلیز" ہی چل رہا تھا کیوں کہ عمران فلیٹ میں ہوتا ہی نہیں تھا! ظاہر ہے کہ ایسے کسی زمانے میں اس کی پناہ گاہ دانش منزل ہی ہوسکتی تھی جب کچھ نامعلوم لوگ اسے مار ڈالنے کے درپے ہوں۔

جولیا نے سلسلہ منقطع کرکے دانش منزل کے لئے ٹرانسمیٹر نکالا! اور بولی۔ " ہیلو۔۔ ہیلو۔۔ ایکس ٹو پلیز۔۔! ایکس ٹو۔۔ ہلو۔۔ ہلو۔۔ ایکس ٹو۔ ایکس ٹو"۔

"ہلو۔۔!" آواز آئی اور یہ ایکس ٹو ہی کی آواز تھی۔

"یہاں صفدر موجود ہے۔۔!"

"تو پھر۔۔!"

"وہ کچھ کہنا چاہتا ہے۔۔ کیا فون استعمال کیا جائے"۔

"میں جانتا ہوں وہ جو کچھ کہنا چاہتا ہے۔ اس سے کہو کہ دو دن کی تھکن بڑی اچھی نیند لاتی ہے"۔

"بہتر ہے!"

"غالباً تم سوچ رہی ہوگی کہ اس عمارت پر چھاپہ کیوں نہ مارا جائے"۔

"جی ہاں قدرتی بات ہے"۔

"لیکن تمہیں معلوم ہونا چاہیئے کہ مجھے سرغنہ کی تلاش ہے۔ وہ اس عمارت میں نہیں تھا! اور اب تو وہاں تمہیں ایک پرندہ بھی نہیں ملے گا!"

"میرے لئے کیا حکم ہے؟"

"وقت آنے پر مطلع کیا جائے گا۔ اور کچھ؟"

"جی نہیں!"

"اوور اینڈ آل۔۔!"

جولیا نے سوئچ آف کردیا اور صفدر کی طرف مڑی جو بہت زیادہ متحیر نظر آرہا تھا!

"یہ سب کچھ جانتا تھا!" صفدر نے آہستہ سے کہہ کر جلدی جلدی پلکیں جھپکائیں اور ختم ہوئے سگریٹ سے دوسرا سگریٹ سلگانے لگا۔ پھر دو تین گہرے کش لے کر بولا۔ " وہ جانتا تھا مگر اس نے مطلق پرواہ نہ کی کہ مجھ پر کیا گذرے گی!"

"مگر تمہیں تو عمران نے اس آدمی کے تعاقب کے لئے کہا تھا"۔

"عمران۔ نتائج کا ذمہ دار تو نہیں ہے!" صفدر نے کہا! "ایکس ٹو کو علم تھا آخر اس نے ہماری مدد کیوں نہیں کی؟"

"صفدر صاحب آپ کو تعاقب کے لئے کہا گیا تھا! اس سے دور رہ کر اس کی نظروں سے بچ کر! عمران نے یہ تو نہ کہا ہوگا کہ آپ اس کے ساتھ بلیرڈ کھیلنا شروع کردیں"۔

"ہاں مجھ سے ہی غلطی ہوئی تھی"۔

"ہوسکتا ہے اسی غلطی کی پاداش میں یہ تمہاری سزا رہی ہو کہ ایکس ٹو نے حالات سے واقف ہونے کے باوجود بھی تمہاری کوئی مدد نہ کی!"

صفدر کچھ نہ بولا! اس کی بھنویں سمٹ گئی تھیں اور پیشانی پر کئی سلوٹیں ابھر آئی تھیں!

کچھ دیر بعد جولیا نے جوزف کا تذکرہ چھیڑدیا!

"وہ عمران ہی کی طرح عجیب ہے! بظاہر ڈیوٹ۔ لیکن۔ بہرحال اس نے مجھے کسی طرح بھی یہ نہیں بتایا کہ وہاں کیسے پہنچا تھا!"

"مگر اب وہ رہتا کہاں ہے؟"

"خدا جانے۔۔!"

"عمران کے فلیٹ میں تو بہت دنوں سے نہیں دیکھا گیا"۔

"ہوں۔ یہ بتاؤ۔ سرسوکھے کا کیا قصہ ہے۔ یہ کون ہے؟" وہ عمران کو کیوں تلاش کر رہا ہے! وہ لوگ یہ بھی جاننا چاہتے ہیں کہ سرسوکھے عمران کی تلاش میں کیوں ہے اور اس نے ہمارے دفتر سے کیوں رابطہ قائم کیا ہے۔۔!"

"سرسوکھے یہاں کا ایک دولت مند آدمی ہے! وہ اس لئے ہمارے فرم سے رجوع ہوا ہے کہ ہم اس کی فرم کے لئے فارورڈنگ اور کلیرنگ کریں! لیکن میں یہ نہیں جانتی کہ اسے عمران کی تلاش کیوں ہے! یہ تو بہت برا ہوا کہ آفس بھی ان کی نظروں میں آگیا ہے"۔

"میرا تو خیال ہے کہ وہ ہمارے چیف ایکس ٹو کے متعلق بھی کچھ نہ کچھ ضرور جانتے ہیں"۔

"اور عمران کے قول کے مطابق یہ لوگ وہی ہیں جن سے آتشدان کے بت والے کیس میں مڈبھیڑ ہوئی تھی۔۔! وہ قصہ وہیں ختم نہیں ہوگیا تھا!" جولیا نے کہا اور کسی سوچ میں پڑ گئی!

دفعتاً فون کی گھنٹی بجی اور اور جولیا نے ریسیور اٹھالیا!

"ہیلو۔۔!"

"میں ہوں"۔ ایکس ٹو کی آواز آئی۔ سرسوکھے کا کیس ایک بار پھر دہراؤ۔ تفصیل سے۔۔!"

جولیا نے شروع سے اب تک کے واقعات دہرانے شروع کردیئے لیکن پھر یک بیک اسے خیال آیا کہ اس نے اصلیت صفدر کو نہیں بتائی! اور وہ اب بھی یہیں موجود ہے۔ لہذا اس نے سونے کی اسمگلنگ کی طرف سے آنے سے پہلے کہا۔ "صفدر یہیں موجود ہے"۔
"پرواہ نہیں"۔ ایکس ٹو کی آواز آئی۔ "صفدر سے اس سلسلے میں کچھ بھی نہ چھپاؤ! وہ ان لوگوں میں سے ہے جن پر میں بہت زیادہ اعتماد کرتا ہوں"۔

پھر جیسے ہی جولیا نے سونے کی اسمگلنگ کی کہانی چھیڑی صفدر اسے گھورنے لگا!

آخر میں جولیا نے پوچھا۔"کیا آپ کو علم ہے کہ جن لوگوں نے صفدر کو پکڑا تھا وہ سرسوکھے میں بھی دلچسپی لے رہے ہیں"۔

"نہیں میں نہیں جانتا"۔

"انہوں نے صفدر سے یہ معلوم کرنے کے لئے سختی برتی تھی"۔

"کیا معلوم کرنے کے لئے۔ جملے ادھورے نہ چھوڑا کرو" ایکس ٹو غرایا۔

"معافی چاہتی ہوں جناب! وہ یہ معلوم کرنا چاہتے تھے کہ سرسوکھے عمران کی تلاش میں کیوں ہے! یہ معلوم کرنے کے لئے انہوں نے صفدر پر چابک برسائے تھے۔ ڈھمپ اینڈ کو اور عمران کا تعلق بھی ان کے لئے الجھن کا باعث بنا ہوا ہے"۔

"اوہ۔۔ اچھا تو۔۔ اب سرسوکھے کو عمران سے ملا دو"۔ ایکس ٹو نے کہا۔

"مگر میں اسے کہاں ڈھونڈوں؟"

"کل صبح سرسوکھے کو گرینڈ ہوٹل میں مدعو کرو! عمران پہنچ جائے گا"۔

"بہت بہتر جناب۔۔!"

دوسری طرف سے سلسلہ منقطع ہوگیا۔





\endR

\bigskip\bigskip\bigskip

\centerline{\hdr{(۱۱)}}
\bigskip\bigskip\bigskip
\beginR
\body



دوسری صبح تقریباً نو بجے جولیا گرینڈ ہوٹل میں سرسوکھے کا انتظار کر رہی تھی اور اسے یقین تھا کہ اب سرسوکھے سے نجات مل جائے گی۔ ظاہر ہے کہ اب تک وہ عمران ہی کے سلسلے میں اس کےساتھ رہی تھی! لیکن اب عمران خود ہی اس سے ملنے والا تھا!

پھر کیا؟ اب بھی اس کی گلوخلاصی نہ ہوگی؟ جولیا کے پاس اس وقت بھی اس سوال کا کوئی واضح جواب نہیں تھا!

ٹھیک نو بج کر دس منٹ پر سرسوکھے ڈائننگ ہال میں داخل ہوا۔ اس کا چہرہ اترا ہوا تھا اور آنکھیں غمگین تھیں! ایسا معلوم ہو رہا تھا جیسے وہ اپنے کسی عزیز کے کریا کرم سے واپس آیا ہو۔۔!

جولیا نے خوش اخلاقی سے اس کا استقبال کیا!

"بس آجائیں گے تھوڑی دیر میں"۔

اس نے غور سے جولیا کی طرف دیکھا۔ ایک ٹھنڈی سانس لی اور دوسری طرف دیکھنے لگا! ایسا کرتے وقت وہ بےحد مضحکہ خیز لگا تھا! جولیا نے نہ جانے کیسے اپنی ہنسی ضبط کی تھی۔

"پچھلی شام آپ مجھ سے ایک منٹ کے لئے بھی نہیں ملی تھیں؟" دفعتاً اس نے سرجھکا کر آہستہ سے کہا!

"میرے چند دوست۔۔"۔

"ٹھیک ہے"! وہ جلدی سے بولا۔ دیکھئیے مجھے غلط نہ سمجھئیے گا! آخر مجھے کیا حق حاصل ہے کہ آپ سے ایسی گفتگو کروں۔ میرے خدا۔۔!"

اس نے دونوں ہاتھوں سے اپنا چہرہ چھپا لیا! اور جولیا کا دل چاہا کہ ایک کرسی اٹھا کر اسی پر توڑ دے۔ گدھا کہیں کا۔ آخر خود کو سمجھتا کیا ہے!

"وہ دیکھئیے"۔ سرسوکھے نے تھوڑی دیر بعد کہا۔ "میں کیا بتاؤں بعض اوقات مجھ سے بچگانہ حرکتیں سرزد ہوجاتی ہیں! بھلا بتائیے یہ بھی کوئی کہنے کی بات تھی مگر زبان سے نکل ہی گئی۔ اسے یوں سمجھئیے۔ دیکھئیے! بالکل بچوں کی طرح۔۔! وہ ٹھہرئیے۔۔ مجھے ایک واقعہ یاد آرہا ہے۔ دیکھئیے شاید آپ اسی سے میرے احساسات کا اندازہ کرسکیں۔ میری ایک بھابی تھیں! میں انہیں بہت پسند کرتا تھا! اور وہ بھی مجھے بہت چاہتی تھیں! ایک دن ان کا ایک کزن آگیا جو میرا ہی ہم سن تھا۔ کچھ دنوں بعد میں نے محسوس کیا کہ اب وہ مجھ پر اتنی مہربان نہیں رہیں جتنی پہلے تھیں۔ بس رو پڑا۔ الگ جاکر۔ کوٹھری میں کھڑا رو رہا تھا کہ بھابی آگئیں۔ میں خاموش ہوگیا۔ وہ رونے کی وجہ پوچھتی رہیں لیکن میں کیا بتاتا! بہرحال مجھے جھوٹ بولنا پڑا۔ میں نے انہیں بتایا کہ میرے پیر میں چوٹ آگئی ہے مجھ سے اٹھا نہیں جاتا۔ انہوں نے مجھے اٹھایا۔ باہر لائیں۔ میرے پیر میں مالش کی۔۔ لیکن میں روتا ہی رہا۔ اب دیکھئیے۔ میں ان سے کیسے کہتا۔ کیسے کہتا کہ وہ اپنے کزن کو مجھ سے زیادہ کیوں چاہتی ہیں۔۔ اسی طرح کل میں کتنا دکھی تھا! بالکل اسی طرح۔ میرا دل چاہ رہا تھا کہ دھاڑیں مار مار کر رونا شروع کردوں! یعنی آپ نے میری طرف آنا بھی گوارہ نہیں کیا۔ اوہ۔۔!"

وہ یک بیک چونک کر خاموش ہوگیا! اس کی آنکھوں سے ندامت کے آثار ظاہر ہو رہے تھے۔ پھر وہ دوبارہ چونک کر بھرائی ہوئی آواز میں بولا۔ "مس جولیانا۔۔ میں آپ سے معافی چاہتا ہوں۔ ایک باکل گدھا اور بےعقل آدمی سمجھ کر معاف کردیجئیے۔ میں آخر یہ ساری بکواس کیوں کر رہا ہوں۔۔ بوائے۔۔"

اس نے بڑے غیر مہذب انداز میں بیرے کو پکارا تھا! ایسا معلوم ہو رہا تھا جیسے وہ اپنی کہی ہوئی باتیں جولیا کے ذہن سے نکال پھینکنے کی کوشش کر رہا ہو۔۔!

"کافی۔۔ اور ایک بڑا پگ وہسکی!" اس نے بیرے سے کہا اور جولیا کی طرف متوجہ ہوا ہی تھا کہ جولیا بولی۔

"پچھلی رات میں نے صرف عمران کے ساتھ اسکیٹنگ کی تھی!"

"نہیں تو۔ میں وہاں موجود تھا میں نے دیکھا پہلے آپ کے ساتھ کوئی اور تھا"۔

"پہلا اور آخری آدمی۔۔!" جولیا مسکرائی۔۔!

"میں نہیں سمجھا!"

"وہ عمران ہی تھا۔۔!"

"نہیں۔۔! مگر۔۔ مجھے اچھی طرح یاد ہے۔ نہیں وہ نہیں ہوسکتے! تم مذاق کر رہی ہو!"

"یقین کیجیئے! وہ میک اپ میں تھا! آج کل وہ کسی چکر میں ہے اور کچھ لوگ اس کے دشمن ہوگئے ہیں اس لئے وہ زیادہ تر خود کو چھپائے رکھتا ہے۔"

"اوہ! بھئی کمال کا آدمی ہے!" سرسوکھے نے بچوں کے سے متحیرانہ لہجے میں کہا۔ "کیا شاندار میک اپ تھا گھنٹوں دیکھتے رہنے کے بعد بھی نہ پہچانا جاسکا"۔

"میں نے بھی اسے صرف آواز سے پہچانا تھا!

"اوہ۔۔!" وہ مضطربانہ انداز میں بولا۔ جس میں دبی ہوئی سی خوشی بھی شامل تھی۔ "تب تو مجھے یقین ہے۔ بالکل یقین ہے کہ میری مشکلات رفع ہوجائیں گی"۔

تھوڑی دیر بعد ایک آدمی تیر کی طرح ان کی طرف آیا اور کرسی کھینچ کر بیٹھ گیا۔

جولیا سٹپٹا گئی! کیونکہ یہ عمران نہیں ہو سکتا تھا اور اگر تھا بھی تو پچھلی رات والے میک اپ میں نہیں تھا!

"فرمائیے جناب!" سرسوکھے غصیلے لہجے میں بولا!

"میرے پیٹ میں درد ہو رہا ہے"۔ آنے والے مسمی صورت بنا کر کہا!

"درد۔ یعنی کہ پین۔ پتہ نہیں فرانسیسی اور جرمن میں اسے کیا کہتے ہیں"۔

"میں پوچھتا ہوں کہ آپ اس میز پر کیوں آئے ہیں"۔ سرسوکھے میز پر ہاتھ مار کر غرایا!

"انہیں دیکھ کر۔۔!" اجنبی نے جولیا کی طرف اشارہ کیا!

"کیا مطلب۔۔!"

"دیکھنے کا مطلب کیسے سمجھاؤں؟"

"تمہارا دماغ تو نہیں خراب ہوگیا۔۔!"

"اگر کچھ دیر تک آپ اسی قسم کی گفتگو کرتے رہے تو یقیناً خراب ہوجائے گا۔ بھلا کوئی تک ہے۔۔ آخر آپ درد کا مطلب نہیں سمجھتے۔۔دیکھنے کا مطلب نہیں سمجھتے! پھر کیا میں درد کو شکرقند اور دیکھنے کو فلفلانا کہوں۔ واہ بھلا آپ مجھے غصے سے کیوں فلفلا رہے ہیں! میرے پیٹ میں تو شکرقند ہو رہا ہے!"

"تمہاری ایسی کی تیسی"۔ سرسوکھے کرسی کھسکا کر کھڑا ہوگیا اور لگا آستین سمیٹنے!

"ارے۔ تم نے تو میری مٹی پلید کردی جولیا! اجنبی نے جولیا سے کہا۔ " تم نے تو کہا تھا کہ تم کسی سرسوکھے کے ساتھ ملو گی۔ یہ تو سرہاتھی نہیں بلکہ سرپہاڑ ہیں۔ پہلوان بھی معلوم ہوتے ہیں۔ اگر انہوں نے ایک آدھ ہاتھ رکھ ہی دیا ہوتا تو میں کہاں ہوں گا! خدا تمہیں غارت کرے!"

جولیا پیٹ دبائے بےتحاشہ ہنس رہی تھی!

"ارے سرسوکھے! یہ عمران ہے!" بدقت اس نے کہا!

"کیا۔۔! اف فہ۔۔ ہاہا۔۔ ہا ہا۔۔ ہاہا!" سرسوکھے نے بھی منہ پھاڑ دیا۔

لیکن اس کی ہنسی خجالت آمیز تھی۔۔!

پھر وہ بیٹھ گیا! لیکن عمران اب بھی ایسی پوزیشن میں بیٹھا ہوا تھا جیسے اب اٹھ کر بھاگا!

"مائی ڈیئر مسٹر عمران آپ واقعی کمال کے آدمی ہیں!" سرسوکھے نے ہانپتے ہوئے کہا!

وہ اسی طرح ہانپ رہا تھا جیسے دور سے چل کر آیا ہو!

عمران چونکہ میک اپ میں تھا اس لئے حماقت کا اظہار صرف آنکھوں ہی سے ہوسکتا تھا! لیکن اس وقت تو آنکھیں سرسوکھے کا جائزہ لینے میں مصروف تھیں!

"اسمگلنگ کی کہانی میں سن چکا ہوں!" عمران نے کہا۔

"مس جولیا نے آپ کو سب کچھ بتایا ہوگا۔۔!"

"جی ہاں سب کچھ!۔۔ آپ اپنے آدمیوں میں سے کس پر شبہ ہے"۔

"دیکھینے! مجھے تو جس اسٹاف پر شبہ تھا اسے پہلے ہی الگ کردیا تھا! فاورڈنگ اور کلیرنگ کا سیکشن ہی توڑ دیا۔۔ لیکن میں یہ بھی نہیں کہہ سکتا کہ موجودہ اسٹاف بےداغ ہے۔ بھلا کیسے کہہ سکتا ہوں! آپ خود ہی سوچیئے!"

"ٹھیک ہے ایسے حالات میں یقینی طور پر کچھ نہیں کہا جاسکتا"۔ عمران سر ہلا کر بولا!

"پھر آپ میرے لئے کیا کریں گے۔۔؟"

"پکوڑے تلوں گا!" عمران نے سنجیدگی سے کہا اور سرسوکھے بےساختہ ہنس پڑا۔۔

"خیر۔۔ خیر۔۔" اس نے کہا! "میں اب یہ معاملہ آپ پر چھوڑتا ہوں! جس طرح آپ کا دل چاہے اسے ہینڈل کیجیئے!۔۔!"

"آپ کو میرے ساتھ تھوڑی سی دوڑ دھوپ بھی کرنی پڑے گی!"

"اس کی فکر نہ کیجیئے! میں موٹا اور بےہنگم ہی سہی! لیکن چلنے کے معاملے میں کسی سے کم بھی نہیں ہوں! مطلب یہ کہ اگر پیدل بھی چلنا پڑے گا۔ جی ہاں"۔

"سواری کا تو کچومر نکل جائے گا! پیدل ہی ٹھیک ہے"۔ عمران سرہلا کر بولا۔

"میں برا نہیں مانتا!" سرسوکھے نے کھسیانی ہنسی کے ساتھ کہا۔

پتہ نہیں کیوں یک بیک جولیا کو عمران پر تاؤ آنے لگا اور سرسوکھے کے لئے ہمدردی محسوس ہونے لگی!

اس نے کہا۔ "اچھا تو سرسوکھے۔۔ اب ہم اس معاملہ کو دیکھ لیں گے! ہوسکتا ہے کہ آپ بہت مشغول ہوں!"

"اوہ۔ بےحد۔۔ بےحد۔۔ اچھا اب اجازت دیجیئے!" سرسوکھےاٹھتا ہوا بولا۔

عمران اسے جاتے دیکھتا رہا۔۔!

"تم اس کا مضحکہ کیوں اڑا رہے تھے؟" جولیا نے غصیلے لہجے میں پوچھا۔

"پھر کیا کروں؟ اتنے موٹے آدمی کو سر پر بیٹھا لوں!" عمران بھی جھلا کر بولا۔

"مجھے اس سے ہمدردی ہے! اتنے بڑے ڈیل ڈول میں ایک ننھا سا بچہ! بےچارا۔۔!"

"خدا تمہیں بھی بےچاری بننے کی توفیق عطا کرے۔۔ اور آئندہ مجھے کوئی اتنا موٹا بیچارہ نہ دکھائے تو بہتر ہے ورنہ میں تو کہیں کا نہ رہوں گا۔ تم ایسے اوٹ پٹانگ آدمیوں سے ملاتی رہتی ہو۔ اچھا ٹاٹا۔۔!"

پھر جولیا اسے روکتی ہی رہ گئی۔۔ لیکن وہ چھلاوے ہی کی طرح آیا تھا اور اسی طرح یہ جا وہ جا۔۔ نظروں سے غائب۔۔!



\endR

\bigskip\bigskip\bigskip

\centerline{\hdr{(۱۲)}}
\bigskip\bigskip\bigskip
\beginR
\body


دوسری شام جولیا آفس سے گھر آکر لیٹ ہی گئی تھی۔۔! بوریت۔۔! وہ سوچ رہی تھی کہ اس کو ذہنی اضمحلال سے کیسے چھٹکارا ملے گا! آج وہ دن بھر اداس رہتی تھی۔ اس کا کسی کام میں بھی دل نہیں لگا تھا!

عمران۔۔! ان ذہنی الجھنوں کی جڑ عمران ہی تھا! اس کے متعلق کسی ذہنی کشمکش میں پڑ کر وہ اپنی ساری زندہ دلی اور مسرور رہنے کی صلاحیت کھو بیٹھی تھی!

یہ عمران اس کے لئے ایک بہت بڑی مصیبت تھا! اس کی عدم موجودگی میں وہ اس کے لئے بےچین رہتی تھی لیکن جہاں سامنا ہوتا اور وہ اپنے مخصوص لہجے میں گفتگو شروع کرتا تو اس کا یہی جی چاہتا کہ اس وقت جو چیز بھی ہاتھ میں ہو کھینچ مارے! ایسا ہی تاؤ اس کی خاموشی پر بھی آتا تھا! کیونکہ خاموشی حماقت انگیز ہوتی تھی!

جولیا نے کراہ کر کروٹ بدلی۔۔ اور آنکھیں بند کی ہی تھیں کہ فون چیخ پڑا۔۔ وہ اٹھی اور ریسیور اٹھا لیا! دوسری طرف تنویر تھا۔۔!

"اوہو۔۔ تو گھر ہی پر ہو!" اس نے کہا۔ کیا آج سرسوکھے واقعی سوکھتا ہی رہے گا؟"

"کیا مطلب؟" جولیا غرائی!

"سنا ہے آج کل وہ تمہیں بڑی موٹی موٹی رنگینیاں عطا کر رہا ہے۔۔!"

"خاموش رہو بدتمیز۔۔" جولیا بپھر گئی!

"ارے بس۔۔ تھوکو غضہ۔۔ میں نے تو محض عمران کے جملے دہرائے ہیں! ابھی ابھی اس نے فون پر کہا تھا کہ تم تو خیر پہلے ہی ہاتھ دھوچکے تھے اب میں نے بھی دھو لئے ہیں اور اس وقت انہیں تولیے سے خشک کررہا ہوں۔ میں نے پوچھا کیا بکتے ہو کہنے لگا سوکھ رہا ہوں! میں جھنجھلا کر سلسلہ منقطع کرنے ہی والا تھا کہ بولا۔

جولیا آج کل ہمالیاتی عشق کا شکار ہوگئی ہے سرسوکھے اسے عشق کے موٹے موٹے نغمے سناتا ہےاور ایک موٹی سی مسکراہٹ جولیا کے ہونٹوں پر رقص کرنے لگتی ہے اور اسے چاند ستارے، دریا کے کنارے حتیٰ کہ ساون کے نظارے بھی موٹے نظر آنے لگتے ہیں۔۔!"

"شٹ اپ!" جولیا حلق پھاڑ کر چیخی اور سلسلہ منقطع کردیا۔۔

وہ کانپ رہی تھی! اسے ایسا محسوس ہو رہا تھا جیسے رگوں میں خون کی بجائے چنگاریاں دوڑ رہی ہوں!

"سور۔۔ کمینہ۔۔ وحشی۔۔ درندہ!" وہ دانت پیس کر بولی اور منہ کے بل تکیے پر گر گئی۔۔!

تھوڑی دیر تک بےحس وحرکت پڑی رہی! پھر اٹھی اور سرسوکھے کے نمبر ڈائیل کئے! وہ بھی اتفاق سے مل ہی گیا فون پر!

"کون ہے۔۔؟"

"فٹز واٹر۔۔"

"اوہ۔۔ کہیے کہیے۔۔!"

"آپ سے نہیں ملتی تو دل گھبراتا رہتا ہے۔۔!" جولیا ٹھنک کر بولی! اور پھر بڑا برا سا منہ بنایا۔

"اوہو۔۔ تو میں آجاؤں۔۔ یا آپ آرہی ہیں!"

"کسی اچھی جگہ ملیے۔۔!"

"اچھا۔۔ جاگیردار کلب کیسا رہے گا؟"

"اوہو۔۔ بہت شاندار۔۔ پھر آپ کہاں ملیں گے۔۔؟"

"میں آپ کے گھر ہی پر آرہا ہوں!"۔۔ سرسوکھے کا لہجہ بےحد پرمسرت تھا! بالکل ایسا ہی معلوم ہو رہا تھا جیسے کسی بچے سے مٹھائی کا وعدہ کیا گیا ہو!

سلسلہ منقطع کرکے جولیا لباس کا انتخاب کرنے لگی۔۔ یہ عمران آخر خود کو سمجھتا کیا ہے۔ وہ سوچ رہی تھی! بیہودہ کہیں کا۔۔ دوسروں کے جذبات کا احترام کرنا تو آتا ہی نہیں۔۔ جانور۔۔ خیر دیکھوں گی! تم بھی کیا یاد کرو گے۔ اب سرسوکھے ہی سہی۔۔!

سرسوکھے آدھے گھنٹے کےاندر ہی اندر وہاں پہنچ گیا۔ جولیا بےحد دلکش نظر آرہی تھی! اس نے بڑی احتیاط اور توجہ سے میک اپ کیا تھا اور لباس کا تذکرہ ہی فضول ہے کیونکہ گھٹیا سے گھٹیا لباس بھی اس کے جسم پر آنے کے بعد شاندار ہوجاتا تھا۔ وہ ایسی ہی جامہ زیب تھی۔۔!

جاگیردار کلب پہنچنے میں دیر تو نہ لگتی لیکن واقعہ ہی ایسا پیش آیا جو دیر کا سبب تو بن گیا تھا لیکن جولیا کی سمجھ میں نہیں آسکا تھا!

جاگیردار کلب پہنچنے کے لئے ایک ایسی سڑک سے گذرنا پڑتا تھا جو زیادہ کشادہ نہیں تھی اور عموماً سرِشام ہی اپنی رونق کھو بیٹھی تھی! وہ اس سڑک ہی پر تھے کہ جولیا نے محسوس کیا جیسے ان کا تعاقب کیا جارہا ہو! دیر سے ایک کار پیچھے لگی ہوئی تھی!

"شاید یہ آگے جانا چاہتا ہے۔۔ ایک طرف ہوجائیے!" جولیا نے کہا!

سرسوکھے نے بھی پلٹ کر دیکھا۔ پچھلی کار اب زیادہ فاصلے پر نہیں تھی!

اس کے اندر بھی روشنی تھی اور ایک بڑا شاندار آدمی اسٹیرنگ کر رہا تھا!

جولیا کو تو وہ شاندار ہی لگا تھا!

سرسوکھے کے حلق سے عجیب سی آواز نکلی اور پھر جولیا نے محسوس کیا جیسے اس نے اپنے ہونٹ سختی سے بند کرلیئے ہوں! اس نے اپنی گاڑی بائیں کنارے کرلی اور پچھلی کار فراٹے بھرتی ہوئی آگے نکل گئی۔۔!

تھوڑی دیر بعد جولیا نے چونک کر کہا۔ "ارے جاگیردار کلب تو شاید پیچھے ہی رہ گیا۔۔!

"جی ہاں۔۔ بس ابھی واپس ہوتے ہیں! یہ کام اچانک نکل آیا ہے"۔

"میں نہیں سمجھی؟"

"ہوسکتا ہے کہ آپ نے اس آدمی کو بار ہا دیکھا ہو! یہ جو اگلی کار میں ہے!"

"جی نہیں! میں نے تو پہلے کبھی نہیں دیکھا۔۔" جولیا بولی!

"تعجب ہے آپ فارورڈنگ کلیرنگ کا کام کرتی ہیں لیکن اسے نہیں جانتیں میرا خیال تھا کہ یہ بھی آپ کے کاروباری حریفوں میں سے ہوگا!اس کا بھی تو فارورڈنگ کلیرنگ کا بزنس ہے شاید۔۔!"

"پتہ نہیں! میں نہیں جانتی!"

"کسی زمانے میں میرے یہاں اسسٹنٹ منیجر تھا"۔ سرسوکھے نے ٹھنڈا سانس لے کر کہا۔ "لیکن بے ایمان آدمی ہے اس لئے میں نے اسے الگ کردیا تھا!"

"تو کیا آپ اس کا تعاقب کر رہے ہیں!"

"یقیناً کیونکہ میرا خیال ہے کہ وہ میری فرم کے موجودہ جنرل منیجر سے گٹھ جوڑ کئے ہوئے ہے۔ مقصد کیا ہے! میں نہیں جانتا!"

"گٹھ جوڑ کا شبہ کیسے ہوا آپ کو؟"

"جب یہ میرے یہاں تھا تو دونوں ایک دوسرے کے خون کے پیاسے تھے۔۔!"

"تو آپ کس بات کا شبہ کر رہے ہیں۔۔!"

"وہ ایک پرانا اسمگلر ہے۔۔ یہی معلوم ہوجانے پر میں نے اسے اپنی فرم سے الگ کیا تھا۔۔!"

"تب تو پھر اتنے گھماؤ پھراؤ کی بات ہی نہیں تھی! آپ نے پہلے ہی اس کا نام بتایا ہوتا! ہم اسے چیک کرلیتے"۔

"نام تو درجنوں بتائے جاسکتے ہیں! مگر یہ اس وقت میرا تعاقب کیوں کر رہا تھا! مجھے تو یہ دیکھنا ہے۔۔!"

" تو اب آپ اس کا تعاقب کریں گے؟"

قطعی۔۔ قطعی!" وہ بوکھلائے ہوئے لہجے میں بولا! "اس کے علاوہ اور کوئی چارہ نہیں۔ میں نہیں جانتا کہ اب وہ مجھ سے کیا چاہتے ہیں؟ کیا اس لئے میرا تعاقب کیا جا رہا ہے کہ میں نے تم لوگوں سے مدد طلب کی ہے!"

"خیر ایسے لوگوں کے لئے پریشانی کا باعث صرف عمران ہوسکتا ہے!" جولیا نے کہا۔ "کیونکر بعض بڑے جرائم پیشہ اس کی ساکھ سے واقف ہیں!"

"میں یہی کہنا چاہتا تھا مس جولیانا۔۔ آپ کو وہ شام تو یاد ہی ہوگی جب آپ میرے آفس میں میری کہانی سن رہی تھیں۔۔!"

"جی ہاں! میں نے میز پر پائے جانے والے پیر کے نشان کا چربہ عمران کے حوالے کردیا ہے!"

"اوہ۔۔ دیکھیئے وہ کار بائیں جانب مڑرہی ہے۔۔ کیا میں ہیڈلائٹس بجھا دوں"۔

"اگر تعاقب جاری رکھنا ہے تو یہی مناسب ہوگا!" جولیا نے کہا!

سرسوکھے نے اگلی روشنی گل کردی اور پھر وہ بھی بائیں جانب مڑ گیا! تھوڑی دیر بعد وہ پھر شہر کے ایک بھرے پرے حصے میں داخل ہوئے!

"اوہ وہ اپنی گاڑی گرینڈ ہوٹل کی کمپاؤنڈ میں موڑ رہا ہے!" سرسوکھے بڑبڑایا۔۔!

اگلی کار گرینڈ ہوٹل کے پھاٹک میں داخل ہو رہی تھی۔ سرسوکھے نے اپنی گاڑی کی رفتار رینگنے کی حد تک کم کردی۔۔! اگلی کار پارک ہوچکی تھی اس سے وہی آدمی اترا اور بڑے پروقار انداز میں چلتا ہوگا گرینڈ ہوٹل کے صدر دروازے میں داخل ہوگیا۔۔!

ادھر سرسوکھے نے اپنی گاڑی روک دی تھی۔۔!

"اوہ۔۔ میں کیا کروں!" وہ مضطربانہ انداز میں بولا! "آپ ہی بتایئے!"

"کاش میں یہ معلوم کرسکتی کہ آپ کیا چاہتے ہیں"۔

"ہمیشے کے لئے ان بدبختوں کا خاتمہ جن کی وجہ سے نیندیں حرام ہوگئی ہیں مجھ پر۔۔! اس وقت تو میں صرف اپنی جان بچانا چاہتا ہوں! آپ کے لئے کوئی خطرہ نہیں ہے مس جولیا!"

"آپ جو کچھ کہیں۔۔ میں کروں!"

"اوہ دیکھیئے! میں بھی اپنی گاڑی کمپاؤنڈ ہی میں پارک کروں گا اور آپ اسی میں بیٹھ کر میرا انتظار کریں گی!"

"کتنی دیر۔۔!"

"ہوسکتا ہے۔۔ جلد ہی لوٹ آؤں! ہوسکتا ہے دیر ہوجائے"۔

"آپ جائیں گے کہاں۔۔؟"

"اندر۔۔! میں دیکھوں گا کہ وہ کس چکر میں ہے! آپ خود سوچیئے کہ وہ میرا تعاقب کر رہا تھا! پھر آگے نکل آیا۔۔ اب یہاں آرکا ہے۔ کیاوہ میرے گرد کسی قسم کا جال پھیلا رہا ہے!"

جولیا کچھ نہ بولی! سرسوکھے نے گاڑی پھاٹک میں گھمائی اور اسے ایک گوشے میں روکتا ہوا بولا۔

"بس آپ اس کی کار پر نظر رکھیے گا!"

سرسوکھے گاڑی سے اترا اور صدر دروازے کی طرف چل پڑا! اس کی چال میں معمول سے زیادہ تیزی تھی! جولیا کار میں بیٹھی رہی! تقریباً پانچ منٹ گزر گئے! وہ اس آدمی کے متعلق سوچ رہی تھی جسے کار میں دیکھا تھا۔۔ یکایک وہ چونک پڑی ایک نیا سوال اس کے ذہن کے تاریک گوشوں سے ابھرا تھا!۔۔ اگر وہ سرسوکھے کا تعاقب ہی کر رہا تھا تو گاڑی کے اندر روشنی رکھنے کی کیا ضرورت تھی؟

جولیا اس پر غور کرتی رہی! اور اس کا ذہن الجھتا چلا گیا! اب تو ایک نہیں درجنوں سوالات تھے۔۔!

کیا سرسوکھے اسے خطرے میں چھوڑ کر خود کھسک گیا تھا؟ خصوصیت سے اس سوال کا اس کے پاس کوئی جواب نہ تھا! لہذا وہ چپ چاپ سرسوکھے کی گاڑی سے اتر آئی! قریب ہی بڑے بڑے گملوں کی ایک قطار دور تک پھیلی ہوئی تھی! ان میں گنجان اور قدآور پودے تھے جن کی پشت پر تاریکی ہی تھی! جولیا نے سوچا کہ وہ بہ آسانی ان کی آڑ لے سکے گی!

شاید آدھا گھنٹہ گذر چکا تھا لیکن ابھی تک ان دونوں میں سے کسی کی بھی واپسی نہیں ہوئی تھی۔۔!

جولیا سوچنے لگی کہ وہ خواہ مخواہ اپنے پیر تھکا رہی ہے اور اسے ایک بار پھر عمران پر غصہ آگیا۔۔ محض عمران کی وجہ سے وہ اس وقت گھر سے نکل آئی تھی ورنہ دل تو یہی چاہا تھا کہ آفس سے واپسی پر گھنٹوں مسہری پر پڑی رہے گی! تنویر نے فون پر عمران کی گفتگو دہرا کر اسے تاؤ دلا دیا تھا اور وہ سرسوکھے کے ساتھ باہر نکل آئی تھی اور تہیہ کرلیا تھا کہ آئندہ شامیں بھی اسی کے ساتھ گذارے گی!

لیکن اب اسے اپنی جلد بازی کھل رہی تھی! ویسے اس کی ذمہ داری تو عمران ہی پر تھی لہذا وہ سلگتی رہی۔۔!

دفعتاً اسے سرسوکھے نظر آیا جو بڑی تیزی سے اسی کار کی طرف جارہا تھا جس پر تعاقب کرنے والا آیا تھا۔ پھر جولیا نے اسے کار کے انجن میں کچھ کرتے دیکھا اور اس کی آنکھیں حیرت سے پھیل گئیں! آخر وہ کیا کرتا پھر رہا ہے!

اس کے بعد وہ وہیں کھڑے کھڑے اپنی کار کی طرف مڑا اور داہنا ہاتھ اٹھا کر اسے دو تین بار جنبش دی!

غالباً یہ اشارہ جولیا کے لئے تھا کہ وہ ابھی انتظار کرے۔۔ جولیا نے ایک طویل سانس لی۔۔!

سرسوکھے بڑی تیزی سے پھاٹک کی طرف سے چلا جارہا تھا! پھر وہ اس سے گذر کر سڑک پر نکل گیا!

جولیا وہیں کھڑی رہی! پھر اس نے سوچا کہ وہ خواہ مخواہ اپنی ٹانگیں توڑ رہی ہے! جہنم میں گئے سرسوکھے کے معاملات! وہ خود ہی نپٹاتا پھرے گا اسے کیا پڑی ہے کہ خواہ مخواہ اپنا وقت برباد کرے، اپنی انرجی ضائع کرے۔۔ اچانک وہ ایک بار پھر چونک پڑی!
اب وہ آدمی اپنی کار کی طرف جا رہا تھا جو سرسوکھے کی موجودہ بھاگ دور کی وجہ بنا تھا۔۔!

پھر جولیا نے دیکھا کہ وہ کار میں بیٹھ کر اسے اسٹارٹ کرنے کی کوشش کر رہا ہے! تھوڑی ہی دیر بعد وہ انجن کھولے اس پر جھکا ہوا نظر آیا۔۔ اور پھر جب وہ سیدھا کھڑا ہو ا تو اس کے ہاتھوں کی مایوسانہ جنبشیں اس کی بےبسی کا اعلان کر رہی تھیں۔۔!

دفعتاً ایک ٹیکسی ڈرائیور اس کی طرف آیا! دونوں میں گفتگو ہوتی رہی پھر ٹیکسی ڈرائیور نے بھی انجن دیکھا اور کار اسٹارٹ کرنے کی کوشش کی! جولیا محسوس کر رہی تھی کہ وہ آدمی بہت زیادہ پریشان ہے!

پھر ذرا سی دیر بعد اس نے اسے ٹیکسی میں بیٹھتے دیکھا کہ وہ اپنی کار وہیں چھوڑے جا رہا تھا۔۔!

جولیا نے سوچا کہ اب اسے ہر قیمت پر اس کا تعاقب کرنا چاہیئے! ہوسکتا ہے سرسوکھے نے اسے وہاں کچھ دیر روکے رکھنے ہی کے لئے اس کے کار کے انجن میں کوئی خرابی پیدا کی ہو!

اس نے تعاقب کا فیصلہ بہت جلدی میں کیا تھا! کیونکہ ٹیکسی نکلی جارہی تھی ورنہ وہ کوئی قدم اٹھانے سے پہلے مناسب حد تک غور کرنے کی عادی تھی! وہ جھپٹ کر سرسوکھے کی کار میں آ بیٹھی! اور پھر دس منٹ بعد دونوں کاروں کے درمیان صرف سو گز کا فاصلہ رہ گیا! وہ اس فاصلہ کو اس سے بھی زیادہ رکھنا چاہتی تھی لیکن اس بھری پری سڑک پر اس کے امکانات نہیں تھے!

جو ں توں کرکے اس نے تعاقب جاری رکھا! کچھ دیر بعد وہ ٹیکسی شہر کے ایک کم آباد حصے میں داخل ہوئی لیکن یہاں بھی ٹریفک کم نہیں تھا!

دفعتاً وہ ٹیکسی ایک عمارت کی کمپاؤنڈ میں مڑ گئی! پھاٹک کھلا ہی ہوا تھا! جولیا نے اپنی کار کی رفتار کم کرکے اسے سڑک کے نیچے اتار دیا!

دوسری عمارت کی کمپاؤنڈ تاریک پڑی تھی اورچہار دیواری اتنی اونچی تھی کہ اندر کا حال نظر نہیں آسکتا تھا!

پتہ نہیں اس کے سر میں کیا سمائی کہ وہ بھی کار سے اتر کر کمپاؤنڈ میں داخل ہوگئی! چاروں طرف اندھیرا تھا۔ عمارت کی کوئی کھڑکی بھی روشن نہیں تھی!

وہ مہندی کی باڑھ سے لگی ہوئی آگے بڑھ ہی رہی تھی کہ اچانک کوئی سخت سی چیز اس کے بائیں شانے سے کچھ نیچے چھبنے لگی اور ایک تیز قسم کی سرگوشی سنائی دی! "چپ چاپ چلتی رہو۔ یہ پستول بےآواز ہے!"

جولیا کا سرچکرا گیا۔۔ یہ کس مصیبت میں آ پھنسی۔ لیکن وہ چلتی ہی رہی!

اسے ہوش نہیں تھا کہ اندھیرے میں اسے کتنے دروازے طے کرنے پڑے تھے! پھر جب وہ ایک بڑے کمرے میں پہنچی تو اس کی آنکھیں چندھیا کر رہ گئیں۔یہاں متعدد بلب روشن تھے اور ان کی برقی قوت بھی زیادہ تھی!

یہاں اسے وہ آدمی جو ٹیکسی میں بیٹھ کر آیا تھا! تین نقاب پوشو ں میں گھرا ہوا نظر آیا جن کے ہاتھوں میں ریوالور تھے۔۔!

جولیا نے مڑ کر اس کی طرف دیکھا جو اسے یہاں تک لایا تھا۔۔! دوسرے ہی لمحے اس کے حلق سے ایک تحیر زدہ سی چیخ نکلی۔۔! 
یہ سرسوکھے تھا۔۔!

اس کے ہونٹوں پر ایک خونخوار سی مسکراہٹ تھی۔۔! اس نے کہا!

"میں جانتا تھا کہ تم یہی کرو گی۔۔!"

"مم۔۔ مگر۔۔ میں نہیں سمجھی"۔ جولیا ہکلائی!

"ابھی سمجھ جاؤ گئی"۔ سرسوکھے نے خشک لہجے میں کہا! "چپ چاپ یہیں کھڑی رہو! اوہ۔۔ تمہارے ہینڈ بیگ میں ننھا سا پستول ضرور ہوگا! مجھے یقین ہے"۔

اس نے اس کے ہاتھ سے بیگ چھین لیا!

جولیا دم بخود کھڑی رہی! اب وہ پھر اس آدمی کی طرف متوجہ ہوگئی تھی جس کی وجہ سے ان مشکلات میں پڑی تھی۔۔ 
سرسوکھے کا مرکز نگاہ بھی وہیں تھا۔

"کیوں۔۔؟ خفیہ معاہدہ کے کاغذات کہاں ہیں؟" اس نے گرج کر اس آدمی سے پوچھا!

"کیسا خفیہ معاہدہ۔۔ اور کیسے کاغذات؟" وہ آدمی مسکرا کر بولا۔ "میں نہیں جانتا کہ تم کون ہو!"

"اوہ تو کیا تم اسے بھی جھٹلاؤ سکو گے کہ تم رانا تہور علی ہو!"

"اسے جھٹلانے کی ضرورت ہی کیا ہے!"

"کیا لیفٹننٹ واجد والے کاغذات تمہارے پاس نہیں ہیں؟"

" میں جب کسی لیفٹننٹ واجد ہی کو نہیں جانت تو کاغذات کے متعلق کیا بتاؤں۔۔؟"

"تب تو عمران بھی تمہارے لئے اجنبی ہی ہوگا"۔ سرسوکھے کی مسکراہٹ زہریلی تھی!

"یہ کیا چیز ہے۔۔؟"

"خاموش رہو!" سرسوکھے آنکھیں نکال کر چیخا!

"چلو اب خاموش ہی رہوں گا!یقین نہ ہو تو کچھ پوچھ کر آزمالو۔۔!"

"رانا۔۔"

"اب اپنا نام بھی بتا دو۔۔" وہ آدمی مسکرایا! "تاکہ میں بھی تمہیں اتنی ہی بےتکلفی سے مخاطب کرسکوں!"

"رانا تمہارے جسم کا بند بند الگ کردیا جائے گا!"

"ضرور کوشش کرو! میں بھی آدمی کی ٹوٹ پھوٹ کا تجربہ کرنا چاہتا ہوں! میری نظروں سے آج تک کوئی ایسا آدمی نہیں گذرا جس کا بند بند الگ کردیا گیا ہو؟"

"ستون سے باندھ کر کوڑے برساؤ"۔ سرسوکھے نے نقاب پوشوں سے کہا۔

نقاب پوشوں نے اپنے ریوالور جیبوں میں ڈال لئے۔ لیکن اس وقت جولیا کی حیرت کی انتہا نہ رہی جب وہ اس آدمی کی بجائے خود سرسوکھے ہی پر ٹوٹ پڑے۔۔!

"ارے۔۔ ارے! دماغ تو نہیں خراب ہوگیا!" سرسوکھے بوکھلا کر پیچھے ہٹا۔

"ہاں۔۔ دیکھو! دفعتاً وہ آدمی بولا۔ "ہم اسے زندہ چاہتے ہیں! تاکہ اس پر ہودہ کسوا کر سواری کے کام میں لاسکیں۔۔ رانا تہور علی صندوقی کا ہاتھی بھی عام ہاتھیوں سے الگ تھلگ ہونا چاہیئے۔۔!"

جولیا کو تو ابھی بھانت بھانت کی حیرتوں سے دوچار ہونا تھا! سرسوکھے ان تینوں کے لئے لوہے کا چنا ثابت ہوا۔۔!

سارے کمرے میں وہ انہیں نچاتا پھر رہا تھا۔۔ اتنے بھاری جسم والا اتنا پھرتیلا بھی ہوسکتا ہے! حیرت! حیرت!! جولیا کو تو ایسا لگ رہا تھا جیسے کسی بھوت خانے میں آ پھنسی ہو! سرسوکھے آدمی تو نہیں معلوم ہو رہا تھا۔۔!

بالکل ایسا ہی معلوم ہو رہا تھا جیسے کسی ہاتھی نے چیتے کی طرح چھلانگیں لگانی شروع کردی ہو۔۔!

سب سے لمبا نقاب پوش حلق سے طرح طرح کی آوازیں نکالتا ہوا اسے پکڑنے کی کوشش کر رہا تھا۔۔!

رانا تہویر علی ریوالور سنبھالے دروازوں کی روک بنتا پھر رہا تھا کہ کہیں سرسوکھے کسی دروازے سے نکل کر فرار نہ ہوجائے! ویسے اس کی آنکھوں میں کچھ اس قسم کے تاثرات پائے جاتے رہے تھے جیسے اچھی فیلڈنگ کرنے والے کسی چست وچالاک بچے کی آنکھوں میں پائے جاتے ہیں۔ جولیا کبھی اس کی طرف دیکھنے لگی تھی اور کبھی سرسوکھے کی طرف۔۔!

"سرسوکھے تم ابھی تھک جاؤ گے"۔ دفعتاً رانا نے کہا۔

"اسی طرح صبح ہوجائے گی"۔ سرسوکھے نے قہقہہ لگایا۔ تم مجھ پر فائر کیوں نہیں کرتے؟"

"میں ایک بلیک میلر ہوں سرسوکھے!" رانا کہا۔ "کیا تم سودا کرو گے؟"

"میں جانتا تھا!"۔۔ سرسوکھے نے بےتکان قہقہ لگایا۔ وہ اب بھی ان تینوں کو ڈاج دیتا پھر رہا تھا!

جولیا دروازے کی طرف کھسک رہی تھی۔۔ رانا نے اسے للکارا۔

"خبردار اگر تم اپنی جگہ سے ہلیں تو تمہاری لاش یہیں پڑے پڑے سڑ جائے گی!" جولیا ٹھٹک گئی!

"اپنے آدمیوں کو روکو"۔۔ سرسوکھے نے کہا۔

"اوہ۔۔ تم تینوں دفع ہوجاؤ"۔ رانا نے ہاتھ ہلا کر کہا! اور تینوں نقاب پوش اسے چھوڑ کر ایک دروازے سے نکل گئے!

"تم ادھر چلو۔۔!" سرسوکھے نے جولیا سے کہا۔۔ اور رانا نے ریوالور کی نال کو جنبش دے کر سرسوکھے کی تائید کی! جولیا ان کے قریب آگئی!

"تم اسے کہاں لئے پھر رہے ہو سرسوکھے؟ جانتے ہو یہ کون ہے؟" رانا نے پوچھا۔

"میں سب کچھ جانتا ہوں تم معاملے کی بات کرو!"

"ساڑھے تین لاکھ"۔

"بہت ہے۔۔ میں نہیں دے سکتا۔۔!"

"تب پھر میں دوسروں سے بھی بزنس کرسکتا ہوں۔۔! مگر نہیں!

میں تم سے بات ہی کیوں کرو۔۔ معاملہ تو تمہارے چیف ہی سے طے ہو سکے گا۔۔!"

"میرا کوئی چیف نہیں ہے!" سرسوکھے غرایا! "میں مالک ہوں"۔

"تب پھر تم ہی معاملہ طے کرو"۔

"میں ایک لاکھ سے ڈیڑھ لاکھ تک بڑھ سکوں گا! لیکن اس کے بعد گنجائش نہیں ہے!"

"اس سے بہتر تو یہی ہوگا کہ میں عمران سے ہار مان کر اپنا پیچھا چھڑاؤں!"

"تم ایسا نہیں کرسکتے!" سرسوکھے گرجا! "میں کتوں کے راتب میں اضافہ کرنے کی سکت رکھتا ہوں۔۔ ساڑھے تین لاکھ ہی سہی"۔

اچانک رانا نے اچھل کر اس کی توند پر ایک زر دار لات رسید کی۔۔!

اور وہ چیخ کر الٹ گیا! اس کے گرنے سے کسی قسم کی آواز پیدا ہوئی تھی!

جولیا اندازہ نہ کرسکی! عجیب سی آواز تھی۔۔ نہ وہ کسی چٹان کے گرنے کی آواز تھی اور نہ۔۔؟ وہ اندازہ بھی کیسے کرسکتی تھی کیونکہ اس نے آج تک نہ تو گوشت کا پہاڑ دیکھا ہی تھا اور نہ اس کے گرنے کی آواز سنی تھی!

"اب تم اٹھ نہ سکو گے سرسوکھے۔۔! رانا نے قہقہہ لگایا۔ "بس کسی ایسی بطخ کی طرح پڑے رہوجو چت لیٹا کر سینے پر کنکری رکھ دی گئی ہو! مجھے اسی کا انتظار تھا۔ مگر تم تو ایسے بھی ڈفر ہو! تم غالباً یہ سمجھتے تھے کہ رانا اتفاقاً ہاتھ آگیا ہے اسی لئے اس پھر بھی غور نہ کرسکے کہ جو شخص کسی سے چھپتا پھر رہا ہو وہ بھلا کار کے اندر روشنی کیوں رکھنے لگا! کار کے اندر میں نے اس توقع پر روشنی کی تھی کہ شاید تم پھنس ہی جاؤ۔۔ وہی ہوا۔۔ یہاں کچھ دیر پہلے تمہارے آدمی تھے جنہیں میرے آدمیوں نے ٹھکانے لگا کر ان کی جگہ لے لی تھی۔۔ مجھے تمہارے سارے اڈوں کا علم تھا! اس لئے اس وقت ہر اڈے پر میرے ہی آدمی موجود ہوں گے! اتنی دردسری تو محض اس لئے مول لی تھی کہ تمہاری زبان سے اعتراف کراسکوں کہ اس کالی تنظیم کے سربراہ تم ہی ہو۔۔ تم ہی وہ وطن فروش ہو جس نے ملک کو تباہ کردینے کی سازش کی تھی۔۔ ہاہا۔۔ تم اٹھ نہیں سکتے۔۔ بس اسی طرح بےبسی سے ہاتھ پیر مارتے رہو! میں یہ بھی جانتا تھا کہ تم لیٹ جانے پر خود سے نہیں اٹھ سکتے تین چار نوکر تمہیں کینچ کھانچ کر بستر سے اٹھاتے ہیں! اسی کام کے لئے تم نے تین چار پہلوان رکھ چھوڑے ہیں۔۔!"

"مجھے۔۔ اٹھاؤ۔۔ دس لاکھ!" سرسوکھے چیخا!

جولیا پھٹی پھٹی آنکھوں سے اسے دیکھ رہی تھی۔۔!

"تم اس فکر میں تھے کہ مجھے اور عمران دونوں کو ٹھکانے لگا دو۔۔ اس لئے اسمگلنگ کی کہانی لے کر عمران کی بیوی کے پاس پہنچ گئے تھے۔۔!"

"اے۔ تم کیا بکواس کر رہے ہو!" جولیا بگڑ گئی۔۔!

"تم عمران کی بیوی نہیں ہو؟" رانا نے بڑی معصومیت سے پوچھا!

"نہیں۔۔!"

"اوہ۔۔ تو اس نے بکواس کی ہوگی۔۔ بہرحال تو پھر تم اس سے اتنی ہی قریب ہوسکتی ہو کہ سرسوکھے تمہارا سہارا لیتا"۔

"وہ صرف میرا دوست ہے"۔

"شوہر بھی دشمن تو نہیں ہوتا!"

"زبان۔۔ بند کرو۔۔! تم کون ہو؟ اور تمہارا ان معاملات سے کیا تعلق ہے؟"

"زبان بند کرلوں گا تو تم سنو گی! خیر۔۔ تم خود ہی اپنی زبان بند کرو اور مجھے سرسوکھے سے گفتگو کرنے دو! ہاں سوکھے! تم ابھی دس لاکھ کی بات کر رہے تھے! دس کروڑ اور دس ارب کی باتیں شروع کرو پھر شاید مجھے سوچنا پڑے کہ مجھے کیا کرنا چاہیئے!"

"تم کیا چاہتے ہو؟" سرسوکھے نے بےبسی سے پڑے ہوئے بھرائی ہوئی آواز میں پوچھا۔۔!

"تمہارے ہاتھوں کے لئے اسپیشل ہتھکڑیاں بنوائی ہیں! دیکھنا چاہتا ہوں کہ فٹ ہوں گی یا نہیں۔۔؟"

"تم بلیک میلر ہو؟۔۔"

"ہاں میں اپنے ملک و قوم کے لئے سب کچھ کرسکتا ہوں! بلیک میلنگ تو تفریحاً بھی ہوجاتی ہے۔۔!"

"تم کون ہو۔۔؟" سرسوکھے نے خوف زدہ سی آواز میں پوچھا!

"جوزف۔۔!" رانا نے جواب دینے کی بجائے آواز دی!

دوسرے ہی لمحے جوزف کمرے میں تھا اور اس کے ہاتھوں میں بڑی بڑی اور وزنی ہتھکڑیاں تھیں۔۔!

"ہتھکڑی لگا دو!" لیکن خیال رکھنا کہ کہیں وہ تمہارے سہارے اٹھ نہ آئے! ورنہ پھر اس کا پیٹ ہی پھاڑنا پڑے گا! میں اس ہاتھی کو زندہ لے جانا چاہتا ہوں۔۔!"

جوزف اس کا مطلب سمجھ گیا تھا اس لئے وہ کوشش کر رہا تھا کہ قوت صرف کئے بغیر ہی اس کے ہاتھوں میں ہتھکڑیاں ڈال دے۔ لیکن اسے کامیابی نہ ہوئی۔ تب رانا نے صفدر کو آواز دی! اور جولیا چونک کر اسے گھورنے لگی صفدر بھی اندر آیا۔۔!

"چلو بھئی۔۔ تم بھی مدد کرو جوزف کی!" رانا نے کہا اور جولیا کھسک کر اس کے قریب آگئی! وہ آنکھیں پھاڑ پھاڑ کر اسے دیکھ رہی تھی!

"فرمائیے محترمہ۔۔!"

"تم کون ہو؟" جولیا نے آہستہ سے پوچھا!

"ہم۔۔ رانا تہور علی صندوقی ہیں!۔۔ ہمارے حضور ابا۔۔ یعنی کہ آنریبل فادر۔۔"

"تم جھوٹے ہو۔۔!" سرسوکھے حلق پھاڑ کر چیخا!" تم ان لوگوں سے بھی کوئی فراڈ کرو گے۔۔ صفدر تم تو عمران کے ساتھی ہو! جولیا اس کے باتوں پر یقین نہ کرو! یہ تمہیں بھی ڈبوئے گا!"

"مگر کچھ دیر پہلے تو یہ تمہاری فرم کا ایک نالائق ملازم تھا"۔ جولیا نے زہریلے لہجے میں کہا!

"کچھ بھی ہو تم اس سے وفا کی امید نہ رکھنا یہ تمہیں اور صفدر کو یہاں سے زندہ واپس نہ جانے دے گا۔۔!"

"مجھے یقین ہے۔۔ تم بکواس نہ کرو!" صفدر نے اس کے منہ پر گھونسہ مارتے ہوئے کہا! وہ دونوں مل کر اس کے ہاتھوں میں ہتھکڑیاں ڈال چکے تھے!

"پچھتاؤ گے۔۔ تم لوگ پچھتاؤ گے۔۔!" سرسوکھے کراہا!

"تم ڈفر ہو سرسوکھے!"۔۔ دفعتاً رانا نے کہا۔ "عمران اس وقت بہت زیادہ خطرناک ہوجاتا ہے جب اسے خود اپنی ہی تلاش ہو۔۔ کیا سمجھے!"

"میں نہیں سمجھا! تم کیا کہہ رہے ہو؟"

"عمران کو عمران کی تلاش تھی اس لئے تم چکر کھا گئےتھے! سرسوکھے اگر عمران کو عمران کی تلاش نہ ہوتی تو تم کبھی روشنی میں نہ آتے!"

"تم ۔۔ تم۔۔ عمران۔۔!"

"ہاں! میں عمران۔۔!" عمران سینے پر ہاتھ رکھ کر خفیف سا خم ہوا اور پھر سیدھا کھڑا ہوتا ہوا بولا۔ "میں جانتا تھا کہ تم لوگ کیپٹن واجد کی گرفتاری کے بعد سے رانا تہور علی کے پیچھے پڑ جاؤ گے! مجھے سرغنہ پر ہاتھ ڈالنا تھا جو اندھیرے میں تھا! لہذا میں نے کیپٹن واجد کے ان ساتھیوں میں جنہیں میں نے دانستہ نظر انداز کردیا تھا یہ بات پھیلانے کی کوشش کی کیپٹن واجد کے بعض اہم کاغذات رانا تہور علی نے عمران کے ہاتھ لگتے ہی نہیں دیئے۔۔ اور عمران اب رانا تہور علی کی تلاش میں ہے اور رانا تہور علی کوشش کر رہا ہے کہ وہ عمران کو ختم ہی کردے! تم نے سوچا کہ کیوں نہ دونوں ہی کو ختم کردیا جائے! لہذا تم ڈھمپ اینڈ کمپنی جا پہنچے۔ مقصد صرف یہ تھا کہ جولیا کا قرب حاصل کرسکو! ہاں مجھے یہ بھی یاد ہے کہ کسی زمانے میں روشی نے بھی تمہاری فرم کی ملازمت کی تھی! لیکن یہ قطعی غلط ہے کہ تم نے مجھے اسی کے توسط سے پہچانا تھا! سیکرٹ سروس والوں پر تمہاری نظریں پہلے ہی سے تھیں اور تم یہ بھی جانتے تھے کہ میں ان کے لئے کام کرتا ہوں! بہرحال تم اس لئے آئے تھے کہ ہم میں گھل مل کر تم بھی رانا تہور علی کی تلاش کرنے والی مہم میں شریک ہو سکو! اور جب وہ مل جائے تو چپ چاپ اسے اور عمران دونوں کو میٹھی نیند سلا دو۔۔! اس لئے تم نے اپنے آفس کے پراسرار اسمگلروں کی کہانی تراشی تھی۔ تقریب کچھ تو بہرملاقات چاہیئے! تمہیں عمران کی تلاش تھی لیکن وہ ہمیشہ بحیثیت عمران تمہاری نظروں میں رہا ہے تم اسے دیکھتے تھے اور نظرانداز کردیتے تھے۔ اگر ایسا نہ ہوتا تو میں تمہیں دھوکا دینے میں کیسے کامیاب ہوتا! تم یہ کیسے سمجھتے کہ عمران اور تہور علی میں چھڑ گئی ہے! وہ دونوں ایک دوسرے کو رگڑ دینا چاہتے ہیں۔۔!"

سرسوکھے نے آنکھیں بند کرلیں تھیں! ایسا معلوم ہو رہا تھا جیسے وہ شروع سے اب تک کے واقعات کو ذہنی طور پر ترتیب دینے کی کوشش کر رہا ہو۔۔!

عمران نے کچھ دیر خاموش رہ کر قہقہہ لگایا۔ "ہاہا ۔ سوکھے رام! جب میرے کرایہ کے آدمیوں نے ندی کے کنارے مجھ پر اور صفدر پر حملہ کیا تھا تو تم یہی سمجھے تھے کے حملہ رانا تہور علی کی طرف سے ہوا تھا۔۔ وہ ڈرامہ میں نے اسی لئے اسٹیج کیا تھا کہ تم یہی سمجھو! موٹی عقل والے موٹے آدمی تم اتنا نہیں سوچ سکتے تھے کے کھلے میں ہم پر فائرنگ ہوئی تھی۔۔ لیکن اس کے باوجود بھی صفدر بچ نکلا تھا۔۔! میں تو خیر دریا ہی میں کود گیا تھا!"

صفدر نے پلکیں جھپکائیں! اسے وہ واقعہ اب بھی یاد تھا! لیکن اصلیت اسی وقت معلوم ہوئی تھی! اس کے فرشتے بھی اس موقع پر یہ نہ سوچ سکتے کہ جس کا تعاقب کرتے ہوئے وہ ندی تک پہنچے تھے عمران ہی کا آدمی تھا اور وہ فائرنگ بھی مصنوعی ہی تھی! ہوسکتا ہے کہ گولیوں والے کارتوس سرے سے استعمال ہی نہ کئے گئے ہوں! لیکن بچ نکلنے کے بعد تو وہ اسے معجزہ ہی سمجھتا رہا تھا! کیونکہ فائرنگ جھاڑیوں سے ہوئی تھی اور وہ کھلے میدان میں تھے اوٹ کے لئے کوئی جگہ نہیں مل سکتی تھی۔۔! ادھر جولیا کوعمران کی تحریر یاد آگئی جو سرکنڈوں کی جھاڑیوں کے درمیان ملی تھی۔۔!

عمران نے پھر قہقہہ لگایا اور بولا! " میں نے خود ہی تمہیں موقع دیا تھا۔ کہ تم میرے کچھ آدمیوں کو پکڑ لو۔۔تاکہ مجھے تمہارے مختلف اڈوں کا علم ہوسکے اور تم دوسرے چکر میں تھے! تم انہیں پکڑواتے تھے اور پھر ایسے حالات پیدا کرتے تھے کہ وہ نکل جائیں۔۔ اور مجھ تک یہ بات پہنچے کہ وہ لوگ سرسوکھے میں بھی دلچسپی لے رہے ہیں! اور مجھے نہ صرف سرسوکھے کی اسمگلنگ والی کہانی پر یقین آجائے بلکہ میں اس الجھن میں بھی پڑ جاؤں کہ آخر ان اسمگلروں کو رانا تہور علی سے کیا سروکار۔۔! تمہیں یقین تھا کہ اس طرح میں تم پر اعتماد کرکے تمہیں رانا تہور علی والے معاملہ میں بھی شریک کرلوں گا! اس طرح تمہیں رانا تک پہنچنے میں آسانی ہوگی۔۔!"

"باس!" دفعتاً جوزف ہاتھ اٹھا کر بولا۔ "تم نے اس رات اندھیرے میں سبز رنگ کی بوٹ دیکھنے کی ہدایت دی! مجھے بتاؤ کہ میں اندھیرے میں سبز رنگ کیسے دیکھ سکتا تھا؟"

"بکواس بند کرو! یہ میں نے اسی لئے کیا تھا کہ تم یہی پوچھنے کے لئے مجھے تلاش کرتے ہوئے شراب خانے میں آؤ۔۔ اور حلق تک تاڑی ٹھونس لو!"

"میں قسم کہا سکتا ہوں کہ مجھے دس سال بعد تاڑی نصیب ہوئی تھی"۔

جوزف نے غالباً تاڑی کا ذائقہ یاد کرکے اپنے ہونٹ چاٹے تھے!

"بکواس بند کرو!" عمران نے کہا اور پھر سرسوکھے کی طرف دیکھنے لگا جو زمین میں پڑا اس طرح ہانپ رہا تھا جیسے کچھ دیر پہلے کی اچھل کود سے پیدا ہونے کی تھکن اب محسوس ہوئی ہو! دفعتاً اس نے کھنکار کر کہا۔

"میں بہت بڑا آدمی ہوں! تمہیں پچھتانا پڑے گا! اگر تم کسی کو میری کہانی سنانا چاہو گے تو وہ تم پر ہنسے گا۔ تمہیں پاگل سمجھے گا!"

"پاگل تو لوگ ویسے ہی سمجھتے ہیں سوکھے رام۔۔ مجھے بالکل دکھ نہ ہوگا۔ لیکن تم خود ہی عدالت کے لئے اپنے خلاف سارا ثبوت مہیا کرچکے ہو۔ یہاں ایک ٹیپ ریکارڈ بھی موجود ہے جس پر شروع سے اب تک ہماری گفتگو ریکارڈ ہوتی رہی ہے۔۔ اور اب بھی ہو رہی ہے۔۔!"

دفعتاً سرسوکھے پر چنگھاڑنے کا دورہ سا پڑ گی! لیکن ٹیپ ریکارڈر ایک بھی صحیح وسالم گالی ریکارڈ نہ کرسکا ہو! سرسوکھے کی ذہنی حالت اتنی اچھی نہیں معلوم ہوتی تھی کہ وہ مختلف گالیوں کو مربوط کر کے انہیں قابل فہم بنا سکتا۔۔!



\endR

\bigskip\bigskip\bigskip

\centerline{\hdr{(۱۳)}}
\bigskip\bigskip\bigskip
\beginR
\body

دوسرے دن عمران جولیا کے فلیٹ میں نظر آیا! وہ اسے بتا رہا تھا کہ اس نے تنویر کو اسی لئے فون پر بور کیا تھا کہ وہ جولیا کو بور کرے۔ عمران کو یقین تھا کہ جولیا تنویر کی زبانی اس کی بکواس سن کر ضرور تاؤ میں آجائے گی اور نتیجہ یہی ہوگا کہ وہ اسی وقت سرسوکھے کے ساتھ نکل کھڑی ہو۔۔!

"سرسوکھے نے تم سے تعاقب کرنے والے کے متعلق بحث کرکے یہی معلوم کرنا چاہا تھا کہ تم رانا کو پہچانتی ہو یا نہیں۔ تم نہیں پہچانتی تھیں! اس لئے اس نے صحیح اندازہ لگایا اور اپنے کام میں لگ گیا۔۔!"

"ایکس ٹو نے مجھے فون پر ہدایت دی ہے کہ میں رانا کے وجود کو راز ہی رکھوں"۔ جولیا نے کہا۔ "اس کا بیان ہے کہ ہم لوگوں میں سے صرف صفدر اور میں رانا کے وجود سے واقف ہوں! بقیہ لوگ نہیں جانتے! تو کیا تمہارا رانا والا رول ابھی برقرار رہے گا؟"

"فی الحال وہ مستقل ہے!"

"تب پھر یہ سمجھنا چاہیئے کہ اس پارٹی میں سب سے زیادہ اہمیت تمہیں ہی حاصل ہے"۔

"یا پھر میری بیوی کو حاصل ہوسکتی ہے!" عمران بڑی معصومیت سے کہا!

جولیا بڑا برا سا منہ بنا کر دوسری طرف دیکھنے لگی۔ اور عمران اٹھتا ہوا بولا! بہرحال مجھے اس غیر ملکی سازش کی جڑوں کی تلاش تھی۔۔ کتنی موٹی جڑ ہاتھ آئی۔۔ ہاہا۔۔ کاش اسے کسی چڑیاگھر کی زینت بنایاجاسکتا! اس کے پھرتیلے پن نے تو میرے بھی چھکے چھڑا دیئے تھے! لیکن گر جانے کے بعد وہ کس طرح بےبس ہوگیا تھا! دنیا کا آٹھواں عجوبہ۔۔!"

اس کے بعد نہ جولیا نے اسے رسماً ہی روکا اور نہ عمران ہی تفریح کے موڈ میں معلوم ہوتا تھا۔


\endR

\centerline{\hdr{(تمام شد)}}

\centerline{\hdr{۴ نومبر ۱۹۵۸ء}}

\bye


